%\documentclass[11pt]{elsart}
\documentclass[ALICE]{ALICE_analysis_notes}

% Use the option doublespacing or reviewcopy to obtain double line spacing
%\documentclass[doublespacing]{elsart}
\usepackage[american]{babel}
\usepackage[utf8x]{inputenc}
\usepackage[colorlinks=true,bookmarksnumbered=true,bookmarksopen=true]{hyperref}

 \usepackage{graphicx}
 \usepackage{subfigure}
 \usepackage{geometry}
% %\geometry{a4paper,inner=28mm, outer=18mm, top=25mm, bottom=30mm,twoside}
 \geometry{a4paper, inner=28mm, outer=18mm, top=25mm, bottom=20mm, headsep=10mm, footskip=12mm, twoside}
%  \geometry{a4paper, inner=18mm, outer=18mm, top=25mm, bottom=20mm, headsep=10mm, footskip=12mm, twoside}
 \usepackage{amsmath}
 \usepackage{amssymb}
 \usepackage{mathcomp}
 \usepackage{array}
\usepackage{helvet}
 \usepackage{color}
% 	\usepackage{hyperref}
\usepackage[dvips]{epsfig}

\hypersetup{
    pdfauthor={Adrian, Mechler},     % author
    colorlinks=true,       % false: boxed links; true: colored links
    linkcolor=red!60!black,        % color of internal links (change box color with linkbordercolor)
    citecolor=red!60!black,       % color of links to bibliography
    filecolor=red!60!black,     % color of file links
    urlcolor=red!60!black          % color of external links
}
\usepackage[all]{hypcap}
 \usepackage{listings}
 \usepackage{cite}
 \usepackage[printonlyused]{acronym}
 \usepackage{units}
 \usepackage{lineno}
\linenumbers
%  \usepackage[numbers]{natbib}
 \usepackage{ascii}
 \usepackage{feynmf}
\usepackage{dcolumn}
\usepackage[point,rounding]{rccol}
\usepackage{rotating}
\usepackage{subcaption}

\usepackage{todonotes}
\newcommand{\warning}[1]{\todo[inline]{#1}\xspace}
% \newcommand{\warning}[1]{\todo[textsize=tiny]{#1}\xspace}

\renewcommand*{\acsfont}[1]{{\color{black}#1}}
\renewcommand*{\acffont}[1]{{\color{black}#1}}
 \usepackage{booktabs}
\setlength{\abovecaptionskip}{4pt plus 0pt minus 0pt}
\setlength{\parskip}{0pt}

\makeatletter
\AtBeginDocument{%
  \renewcommand*{\AC@hyperlink}[2]{%
    \begingroup
      \hypersetup{hidelinks}%
      \hyperlink{#1}{#2}%
    \endgroup
  }%
}
\makeatother
\PHnumber{ALICE-ANA-XXXX}


%%%%%%%%%%%%%% TODO setup
\usepackage{stackengine}
\setstackgap{L}{.5\baselineskip}
\newcommand\markabove[3][orange]{{\sffamily\color{#1}%\hsmash{}%
  \smash{\toplap{#2}{\bfseries#3}}}}
% \newcommand\markbelow[3][orange]{{\sffamily\color{#1}\hsmash{$\downarrow$}smash{\bottomlap{#2}{\scriptsize\bfseries#3}}}}
\usepackage{ifthen}
\newcounter{TODOpointers}
\newcommand\addTODOpointer[2]{%
  \stepcounter{TODOpointers}%
  \newcounter{todoindex\romannumeral\theTODOpointers}%
  \setcounter{todoindex\romannumeral\theTODOpointers}{0}%
  \expandafter\gdef\csname TODOname\romannumeral\theTODOpointers\endcsname{%
    #1}%
  \expandafter\gdef\csname TODOcolor\romannumeral\theTODOpointers\endcsname{%
    #2}%
  \edef\tmp{\theTODOpointers}%
  \expandafter\edef\csname #1TODO\endcsname{\tmp}%
}
\newcommand\TODO[2]{%
  \edef\thenameindex{\csname #1TODO\endcsname}%
  \stepcounter{todoindex\romannumeral\thenameindex}%
  \def\thetmp{\csname thetodoindex\romannumeral\thenameindex\endcsname}%
  \expandafter\gdef%
    \csname todo.\thenameindex.\romannumeral\thetmp\endcsname{#2}%
  \markabove[\csname TODOcolor\romannumeral\thenameindex\endcsname]%
    {r}{#2}%
}
\newcounter{index}
\newcommand\showTODOs[1]{%
  \edef\thenameindex{\csname #1TODO\endcsname}%
  \def\thetmp{\csname thetodoindex\romannumeral\thenameindex\endcsname}%
  \vspace{5ex}\noindent%
  \rule{20ex}{.1ex}
  \par\noindent
  \hspace{5ex}\textbf{ \Large \textcolor{%
    \csname TODOcolor\romannumeral\thenameindex\endcsname}%
    {\csname TODOname\romannumeral\thenameindex\endcsname%
    }}\par\noindent%
  \setcounter{index}{0}%
  \whiledo{\value{index} < \value{todoindex\romannumeral\thenameindex}}{%
    \addtocounter{index}{1}%
    \markabove[\csname TODOcolor\romannumeral\thenameindex\endcsname]{c}%
       {\arabic{index}} %
       \csname todo.\thenameindex.\romannumeral\theindex\endcsname\\%
  }\par%
}

\addTODOpointer{EMCal}{blue}
\addTODOpointer{PCMEMCal}{cyan!80!black}
% \addTODOpointer{PCMPHOS}{red}
\addTODOpointer{PCM}{violet}
\addTODOpointer{PHOS}{green!60!black}
\addTODOpointer{All}{orange}



\renewcommand{\floatpagefraction}{.82}
\renewcommand{\bottomfraction}{.95}
\renewcommand{\topfraction}{.95}
\renewcommand{\textfraction}{.04}
\newcommand{\pT}{$p_{\mbox{\tiny T}}$\xspace}
\newcommand{\pTHard}{$p_{\mbox{\tiny T,hard}}$\xspace}
\newcommand{\sNN}{$\sqrt{s_{\mbox{\tiny NN}}}$\xspace}
\newcommand{\sNNF}{$\sqrt{s_{\mbox{\tiny NN}}}=2.76\,$TeV\xspace}
\newcommand{\sNNMa}{$\sqrt{s_{\mbox{\tiny NN}}}=5.02\,$TeV\xspace}
\newcommand{\s}{$\sqrt{s}$\xspace}
\newcommand{\sth}{\s~=~13~TeV\xspace}
\newcommand{\sei}{\s~=~8~TeV\xspace}
\newcommand{\stw}{\s~=~2.76~TeV\xspace}
\newcommand{\mT}{$m_{\mbox{\tiny T}}$\xspace}
\newcommand{\qT}{$q_{\mbox{\tiny T}}$\xspace}
\newcommand{\xT}{$x_{\mbox{\tiny T}}$\xspace}
\newcommand{\Pb}{{\mbox{Pb--Pb}}\xspace}
\newcommand{\AACol}{{\mbox{A--A}}\xspace}
\newcommand{\Au}{{\mbox{Au--Au}}\xspace}
\newcommand{\pPb}{{\mbox{p--Pb}}\xspace}
\newcommand{\pA}{{\mbox{p--A}}\xspace}
\newcommand{\pp}{pp\xspace}
\newcommand{\MeVc}{MeV/$c$\xspace}
\newcommand{\GeVc}{GeV/$c$\xspace}
\newcommand{\vtwo}{$\nu_2$\xspace}
\newcommand{\vn}{$\nu_n$\xspace}
% \newcommand{\pi0}{$\pi^0$\xspace}
\newcommand{\mum}{$\mu$m\xspace}
\newcommand{\dEdx}{$\mbox{d}E/\mbox{d}x$\xspace}
\newcommand{\RConv}{R_{\mbox{\tiny conv}}}
\newcommand{\ZConv}{Z_{\mbox{\tiny conv}}}
\newcommand{\XConv}{X_{\mbox{\tiny conv}}}
\newcommand{\YConv}{Y_{\mbox{\tiny conv}}}
\newcommand{\PhiConv}{\phi_{\mbox{\tiny conv}}}
\newcommand{\RSec}{R_{\mbox{\tiny sec}}}
\newcommand{\ZSec}{Z_{\mbox{\tiny sec}}}
\newcommand{\XSec}{X_{\mbox{\tiny sec}}}
\newcommand{\YSec}{Y_{\mbox{\tiny sec}}}
\newcommand{\PhiSec}{\phi_{\mbox{\tiny sec}}}
\newcommand{\dcaZ}{$dca_z$\xspace}
\newcommand{\pTtrack}{p_{\mbox{\tiny T},\mbox{\tiny track}}}
\newcommand{\EtaToPi}{$\eta/\pi^0$\xspace}
\newcommand{\Etat}{\eta_{\mbox{\tiny track, V0}}}
\newcommand{\EtaV}{\eta_{\mbox{\tiny V0}}}


\newcommand{\gsim}        {\stackrel{>}{\sim}}
\newcommand{\lsim}        {\stackrel{<}{\sim}}
\newcommand{\stat}        {({\mathit{ stat.}})}
\newcommand{\syst}        {({\mathit{ sys.}})}
\newcommand{\com}[1]      {}
\newcommand{\PCM}         {\acs{PCM}}
\newcommand{\PCMEMC}      {\acs{PCM}-\acs{EMCal}}
\newcommand{\PCMPHOS}     {\acs{PCM}-\acs{PHOS}}
\newcommand{\PCMDal}      {\acs{PCM}-\acs{Dal}}
\newcommand{\EMC}         {\acs{EMCal}}
\newcommand{\mEMC}        {mEMC}
\newcommand{\PHOS}        {\acs{PHOS}}
\newcommand{\PE}          {\mbox{P-E}}

\setlength{\textfloatsep}{1em}



% \usepackage{floatrow}

% \usepackage{enumitem}\setlist[description]{font=$\bullet$}

% %%%%%%%%%%%%%% Schrift
    \renewcommand*{\familydefault}{\sfdefault}
% \usepackage[labelfont=bf]{caption}  							%change color sheme
\usepackage[labelfont={color=red!60!black,bf}]{caption} 			%change color sheme
\usepackage{sectsty}
% \allsectionsfont{\color{red!60!black}\bfseries}
\sectionfont{\color{red!60!black}\huge\bfseries}
\subsectionfont{\color{red!60!black}\Large}
\subsubsectionfont{\color{red!60!black}\large}

\usepackage{colortbl}
\newcommand{\RowColor}{\rowcolor{red!50}}
\newcommand{\RowColorb}{\rowcolor{red!20}}
\newcommand{\RowColora}{\rowcolor{red!5}}

\usepackage{fancyhdr}
\pagestyle{fancy}

\rfoot{\includegraphics[height=1.cm]{alicelogo.pdf}}

\begin{document}
%\begin{frontmatter}
%\begin{titlepage}
\title{Analysis Note:\\Neutral meson measurements with PCM, EMC, PHOS, PCM-EMC and PCM-PHOS \\
in ALICE in pp collisions at $\sqrt{s}$~=~13~TeV}
\ShortTitle{ALICE-ANA-XXXX}


\begin{Authlist}
H.~Bossi \Iref{a5},
J.~K\"onig\Iref{a2},
J.~L\"uhder\Iref{a1},
A.~Mechler\Iref{a2},
F.~Bock\Iref{a3},
D.~M\"uhlheim\Iref{a1},
N.~Schmidt\Iref{a6},
J.P.~Wessels\Iref{a1},
H.~B\"usching\Iref{a2},\\
C.~Loizides\Iref{a6},
for the ALICE Collaboration \Iref{a4}
\end{Authlist}

\ShortAuthor{ALICE Analysis Note ANA-XXX}
\Instfoot{a5}{Yale University, New Haven, United States}
\Instfoot{a2}{Institut f\"ur Kernphysik, Universit\"at Frankfurt, Frankfurt, Germany}
\Instfoot{a1}{Institut f\"ur Kernphysik, Westf\"alische Wilhelms-Universit\"at M{\"u}nster, M\"unster, Germany}
\Instfoot{a3}{CERN, Geneva, Switzerland}
\Instfoot{a6}{ORNL, United States}
\Instfoot{a4}{\acs{ALICE}}


%% For running titles
% \ShortAuthor{F.~Bock \emph{et al.}}
%


\vspace{4cm}
\begin{abstract}
In this analysis note, the differential invariant cross-sections of the feed-down corrected $\pi^0$ and $\eta$ mesons for \pp collisions at \sth are presented using the \ac{PCM}, \ac{EMCal}, \ac{PHOS} and hybrid methods, with which one referres to the combination of one photon reconstructed with \ac{PCM} and one calorimeter cluster (\PCMEMC, \PCMPHOS).

Please note that this analysis was carried out analogous to the corresponding analysis at \sei \cite{anaNoteEMC8TeV,anaNotePCMEMC8TeV} and \stw \cite{anaNotePCMEMC2760GeV,anaNoteEMC2760GeV}. Therefore, the analysis notes are based on the same template and parts of the text are very close between those notes which will not be cited explicitly in this note. Thus, it has been ensured that it is not necessary to jump between notes in order to fully understand the analysis.
\end{abstract}

% \begin{keyword},
% \PACS
% \end{keyword}
% \end{frontmatter}
%\end{titlepage}
%%%%%%%%%%%%%%%%%%%%%%%%%%%%%%%%%%%%%%%%%%%%%%%%%%%%%%%%%%%%%%%%%%%%%%%%%%%%%

\newpage

\tableofcontents

\include{PhotonDetection_v2}
\include{PhotonDetection_PCM_v1}
\include{PhotonDetection_EDC_v1}
\include{PhotonDetection_PHOS_v1}
\include{PhotonDetection_hybrids_v1}
\include{NeutralMesons_v1}
\section{Combination}
\subsection{Combination of Different Measurements}\label{chapter:combination}

This chapter is dedicated to describe the combination procedure of the different available measurements of neutral pions, eta mesons and \EtaToPi ratios at \s~=~13~TeV.
First of all
% , the results from the respective single measurements are summarized and visualized.
% % Furthermore, invariant mass plots are given for all reconstruction methods and the different triggers used in the respective analyses.
% All example bins are listed, which are candidates for publication.
% Then
, the determination of correlation factors as well as the actual combination procedure including bin width corrections follow in this chapter which concludes with the obtained weights that are needed to combine the different systems.
The results of the combination and comparisons with available references will be shown in the subsequent chapters, beginning from \hyperref[sec:combMeasure]{Section~\ref*{sec:combMeasure}}.

\subsubsection{Correlation Factors}
 As already shown earlier, we have measured neutral meson spectra in seven different systems. These are \acs{PCM} (1), \acs{EMCal} (2), \acs{PCM}-\acs{EMCal} (3), \acs{PHOS} (4), \acs{PCM}-\acs{PHOS} (5), \acs{PCM}-\acs{Dalitz} (6) and merged \acs{EMCal} (7). Each of these systems comes with its own set of systematic and statistical uncertainties, which in principle may be largely correlated between the systems. However, a possible correlation only applies to systematics as the statistical errors of all systems are assumed to be completely uncorrelated. Possible statistical correlations between the measurements, for instance due the conversions at low radii, are negligible due to the small conversion probability and the small likelihood of reconstructing the respective electron in the calorimeters leading to a meson candidate which on top ends up in the respective integration window. On the other hand, the correlations of systematic errors need to be properly taken into account in the combination procedure.

 In order to combine the results, thus, one has to determine the correlation coefficients for the systematic uncertainties of the 'hybrid' \acs{PCM}-\acs{Calo} and the dalitz method with respect to single \acs{PCM} and \acs{EMCal} or \acs{PHOS} measurements. The coefficients of the correlation matrix might be momentum dependent, as seen in \hyperref[eq:FullCorrMat]{Equation~\ref*{eq:FullCorrMat}}. When setting the coefficients of the reconstruction methods that have no correlation in their systematic errors to zero, one ends up with the following correlation matrix:
 			\begin{eqnarray}
 				\tiny
 				\arraycolsep=1.pt\def\arraystretch{2.5}
 				\mathcal{C}(p_{\mbox{\tiny \text{T}}}) =
 				\left( \begin{array}{ccccccc}
 				1 & 0 & cPCM\_PCMEMC(p_{\mbox{\tiny T}}) & 0 & cPCM\_PCMPHOS(p_{\mbox{\tiny T}}) & cPCM\_PCMDalitz(p_{\mbox{\tiny T}}) & 0 \\
 				0 & 1 & cEMC\_PCMEMC(p_{\mbox{\tiny T}}) & 0 & 0 & 0 & cEMC\_mEMC(p_{\mbox{\tiny T}}) \\
 				cPCMEMC\_PCM(p_{\mbox{\tiny T}}) & cPCMEMC\_EMC(p_{\mbox{\tiny T}}) & 1 & 0 & cPCMEMC\_PCMPHOS(p_{\mbox{\tiny T}}) & cPCMEMC\_PCMDalitz(p_{\mbox{\tiny T}}) & 0 \\
 				0 & 0 & 0 & 1 & cPHOS\_PCMPHOS(p_{\mbox{\tiny T}}) & 0 & 0 \\
 				cPCMPHOS\_PCM(p_{\mbox{\tiny T}}) & 0 & cPCMPHOS\_PCMEMC(p_{\mbox{\tiny T}}) & cPCMPHOS\_PHOS(p_{\mbox{\tiny T}}) & 1 & cPCMPHOS\_PCMDalitz(p_{\mbox{\tiny T}}) & 0 \\
 				cPCMDalitz\_PCM(p_{\mbox{\tiny T}}) & 0 & cPCMDalitz\_PCMEMC(p_{\mbox{\tiny T}}) & 0 & cPCMDalitz\_PCMPHOS(p_{\mbox{\tiny T}}) & 1 & 0 \\
 				0 & cmEMC\_EMC(p_{\mbox{\tiny T}}) & 0 & 0 & 0 & 0 & 1 \\
 				\end{array} \right)
 				\label{eq:FullCorrMat}
 			\end{eqnarray}

 The full correlation coefficients $\mathcal{C}_{ij}(p_{\text{\tiny T}})$ can then be calculated according to \hyperref[eq:CijCorr]{Equation~\ref*{eq:CijCorr}}.
 			\begin{eqnarray}
 				\mathcal{C}_{ij}(p_{\text{\tiny T}}) = \frac{\rho_{ij} S_i(p_{\text{\tiny T}}) ~\rho_{ji} S_j(p_{\text{\tiny T}})}{T_i(p_{\text{\tiny T}}) T_j(p_{\text{\tiny T}})}
         \label{eq:CijCorr}
 			\end{eqnarray}
%
 The $T_{x}(p_{\text{\tiny T}})$ represent the total errors, which have been calculated from the quadratic sum of the statistical errors $D_x (p_{\text{\tiny T}})$ and systematic errors $S_{x}(p_{\text{\tiny T}})$, while the two latter errors have been visualized in \hyperref[fig:combStatSysErr]{Figure~\ref*{fig:combStatSysErr}}.\TODO{Combination}{Add eta plots}
 The fraction of the correlated systematic error of a system $x$ with respect to $y$ is given by $\rho_{xy}(p_{\text{\tiny T}})$:
 \begin{equation}
 \rho_{xy}(p_{\text{\tiny T}}) = \frac{\sqrt{S_{x}^{2}(p_{\text{\tiny T}})-U_{xy}^{2}(p_{\text{\tiny T}})}}{S_{x}(p_{\text{\tiny T}})}
 \end{equation}
 where $U_{xy}$ is the uncorrelated systematic error of $x$ with respect to $y$.
 It becomes clear that in general $\rho_{xy}\neq\rho_{yx}$ holds.
 The $\rho_{xy}$ are momentum dependent, hence referred to as $\rho_{xy}(p_{\text{\tiny T}})$.
 Although it has been found that the dependency is rather mild, we are considering the proper momentum dependence.
 In most of the momentum bins, the systematic error will dominate the total errors for the individual and combined measurement, while only for the very first or for the last bins of each trigger this is not the case.
 For each system, the full available transverse momentum range has been considered.
% Furthermore, the uncorrelated systematic errors from \acs{PCM}-\acs{EMCal} with respect to \acs{PCM} are with full size all the calorimeter related error sources as well as trigger and efficiency errors.
% Further details concerning the different systematic error sources are documented in the respective analysis notes quoted in \hyperref[tab:input]{Table~\ref*{tab:input}}.
%
 By careful evaluation of the systematics, the following $\rho_{xy}(p_{\text{\tiny T}})$ have been found as listed in \hyperref[fig:corrFactors]{Figure~\ref*{fig:corrFactors}}\TODO{Combination}{Make this plot!}, where $\rho_{xy}(p_{\text{\tiny T}})$ can be identified with $\rho_{A\_B}$ from the y-axis description. The three figures show the coefficients for $\pi^0$, $\eta$ and \EtaToPi, respectively. Furthermore, they just contain those combinations of detection systems which are at least overlapping in one transverse momentum bin and have at minimum one correlated systematic error source. Therefore, if no markers are shown, the correlation coefficients are found to be zero and hence, there are no markers shown for those cases.
%
 		\begin{figure}[h]
 			\centering
 			\includegraphics[width=0.44\textwidth]{figures/Combination/CorrFactors/pp13TeV_Systems_Pi0_corrFactors.eps}
 			\includegraphics[width=0.44\textwidth]{figures/Combination/CorrFactors/pp13TeV_Systems_Eta_corrFactors.eps}\\
 			\includegraphics[width=0.44\textwidth]{figures/Combination/CorrFactors/pp13TeV_Systems_EtaToPi0_corrFactors.eps}
 			\caption{The momentum dependent factors $\rho_{xy}(p_{\text{\tiny T}})$ (to be identified with: $\rho_{A\_B}$) for the three different combinations of $\pi^0$, $\eta$ and \EtaToPi. The momentum dependence is nicely visualized which is in general rather mild even over broad transverse momentum intervals. The horizontal lines give orientation concerning the momentum dependence and absolute values of the correlation factors. All possible combinations of detection systems which are not shown in the legend show no correlation in between. Therefore, factors have not been visualized since they are anyway '0' in those cases.}
 			\label{fig:corrFactors}
 		\end{figure}

	\begin{figure}[h]
		\centering
		\includegraphics[width=0.44\textwidth]{figures/Combination/Pi0_RelStatErr_Logy.pdf}
		\includegraphics[width=0.44\textwidth]{figures/Combination/Pi0_RelSysErr.pdf}\\

		\caption{The statistical (left) and systematic uncertainty for the different reconstruction methods.}
		\label{fig:combStatSysErr}
	\end{figure}
%
%
\subsubsection{Combination of Measurements}
\label{sec:combMeasure}
 After having derived those correlation coefficients, the combination of the spectra can be calculated based on the standard \ac{BLUE}-algorithm \cite{Lyons1988rp,Valassi2003mu,Lyons1986em,barlow1989statistics,Valassi2013bga}. Here $Q_{a}(p_{\text{\tiny T}})$ represents the $a^{th}$ individual measurements and $0 < a < n$. Each of them comes with an error $T_a (p_{\text{\tiny T}})$, as defined before. The $Q_a$ can be summarized in the vector $\mathbf{Q}$, while the respective weight can be written as vector $\pmb{\omega}$ and its components are written as $\omega_{a}$.
 		\begin{eqnarray}
 			\pmb{\omega}(p_{\text{\tiny T}})		&=& \mathcal{C}^{-1}\textbf{U}/(\textbf{U}^{T}\mathcal{C}^{-1}\textbf{U}) \hspace{0.5cm}\text{, with } \textbf{U}= \text{ unity vector }\\
 			\omega_{a}(p_{\text{\tiny T}})					&\equiv& \frac{\sum_{b=1}^{n} \mathcal{H}_{ab}}{\sum_{a,b=1}^{n} \mathcal{H}_{ab}} \hspace{1.32cm} \text{, with } \mathcal{H} = \mathcal{C}^{-1} \text{ and } \mathcal{H}_{ab} \text{ its elements} \\
 			\langle Q (p_{\text{\tiny T}}) \rangle  &=& \pmb{\omega}^{T}(p_{\text{\tiny T}}) \mathbf{Q}(p_{\text{\tiny T}}) \\
 																							&=&	\sum_{a=1}^{n} \omega_{a}(p_{\text{\tiny T}}) Q_{a} (p_{\text{\tiny T}})
 		\end{eqnarray}

 The obtained weights ($\omega_{a}$) for the different measurements are shown in \hyperref[fig:Weights]{Figure~\ref*{fig:Weights}}. The plots from that figure also show the respective ranges in which the spectra of each individual trigger are considered for the final result in the end. The total combined errors as well as combined statistical and systematic errors can be found in \hyperref[fig:combTotErr]{Figure~\ref*{fig:combTotErr}}.
 \TODO{Combination}{Add Eta To Pi0 plot!!!}
 		\begin{figure}[h]
 		\centering
 			\includegraphics[width=0.49\textwidth]{figures/Combination/Pi0_WeightsMethods.pdf}
 			\includegraphics[width=0.49\textwidth]{figures/Combination/Eta_WeightsMethods.pdf}\\
 			\includegraphics[width=0.49\textwidth]{figures/Combination/EtaToPi0_WeightsMethods.pdf}
 			\caption{Obtained weights using the BLUE-algorithm for the neutral pion (top left), eta meson (top right) as well as the \EtaToPi-ratio (bottom plot) for \pp collisions at \s~=~13~TeV. }
 			\label{fig:Weights}
 		\end{figure}



%%%%%%%%%  chapter originatly from meson part
\subsubsection{Correction for Finite Bin Width}
\label{sec:BinShift}
The neutral meson spectra are binned in transverse momentum while the bin width of those slices varies for different \pT. The underlying spectrum, however, is steeply falling with increasing transverse momenta.
Thus, the yield measured in a \pT interval does not correspond to the yield which is measured at the center of a \pT bin.
Therefore, the true \pT value for the yield, which has been measured in the given slice, needs to be determined \cite{Lafferty1994cj}.
This can be done with two different approaches. Either shifting the data points horizontally in \pT, such that they correspond to the true \pT for a given bin or by shifting the data points vertically to the true yield for the \pT at the bin center.
Both methods depend on the same underlying model assumptions.
For the spectra, the points are shifted in the \pT direction (also referred to as 'x'-direction), to not change the measured yield, while for the calculation of \EtaToPi-ratio, they will be shifted in the 'y'-direction.
The shifting in the vertical direction for the \EtaToPi-ratio has been chosen due to the fact that the spectra for the neutral pion and eta meson may have different shapes and therefore a shift in x-direction would result in a mismatch of the bin centers of the $\eta$ and the $\pi0$ spectrum.
Thus, for the same binning the points would shift by a different amount in the \pT direction, making it harder to correctly calculate the \EtaToPi-ratio.


 The spectra have been shifted assuming the two-component model fit (\acs{TCM}), developed by A. Bylinkin and A. Rostovtsev \cite{Bylinkin2012bz,Bylinkin2011yi,Bylinkin2015xya} as an approximation of the underlying spectrum.
 		\begin{eqnarray}
 		\label{eq:TCM_Binshift}
  			E \frac{\mbox{d}^3 \sigma}{\mbox{d}p^3} = A_e \exp{\frac{-(\sqrt{p_{\text{\tiny T}}^2 + M^2} - M)}{T_{2}}} + \frac{A}{\left( 1+ \frac{p_{\text{\tiny T}}^2}{T^2 n}\right)^{-n}}
  			\label{eq:TwoComponentModel}
  		\end{eqnarray}
% % 		\begin{equation}
% % 		 	\frac{\mbox{d}^2N}{\mbox{d}y \mbox{d}p_{\mbox{\tiny T}}} = \frac{(n_{\mbox{\tiny Tsallis}}-1)(n_{\mbox{\tiny Tsallis}}-2)}{n_{\mbox{\tiny Tsallis}}T [n_{\mbox{\tiny Tsallis}}T + m(n_{\mbox{\tiny Tsallis}}-2)]} \times A \times p_{\mbox{\tiny T}} \times \left(1 + \frac{m_{T}-m}{n_{\mbox{\tiny Tsallis}}T}\right)^{-n_{\mbox{\tiny Tsallis}}}.
% % 			\label{eq:Tsallis}
% % 		\end{equation}
% % The parameters $m$ and $m_{T} = \sqrt{m^2 + p_{T}^2}$ correspond to the particle mass and the transverse mass, while $A$, $T$ and $n_{\mbox{\tiny Tsallis}}$ are the free parameters of the fit function.
% % In order to be able to combine different measurements in \acs{ALICE}, this bin width correction is calculated based on the final combined spectra and then applied to the individual spectra based on the combined fit.
% % Moreover, this allows to constrain the fit even more as each individual measurement covers a slightly different transverse momentum region, resulting in smaller correction factors.
% % Hence, the described bin width corrections for the two neutral mesons are only applied for the plots shown in the combination note for \s~=~13~TeV which may be found under \cite{anaNoteComb8TeV}.
% % In this note, only the \EtaToPi-ratio has been shifted according to this chapter and the related final plots from \hyperref[fig:EtaToPi0Final]{Figure~\ref*{fig:EtaToPi0Final}} already contain 'y'-shift spectra for the \EtaToPi-ratio.
% %
% \hyperref[fig:binShiftEtaPi0]{Figure~\ref*{fig:binShiftEtaPi0}} visualizes the bin width corrections for the final combined \EtaToPi-ratios as well as separately for the respective triggers. The bin shift has been performed in 'y' as already explained. Thus, bin shift correction factors are above one. \hyperref[fig:binShiftEtaPi0]{Figure~\ref*{fig:binShiftEtaPi0}} shows that corrections are of the order of 1\% almost throughout the whole transverse momentum region covered and go up to about 3\% for low momenta.
% %
% % 		\begin{figure}[t]
% % 			\centering
% % 			\includegraphics[width=0.49\textwidth]{figures/Meson/MesonTrigger/EtaToPi0_data_BinShiftCorrection.eps}
% % 			\includegraphics[width=0.49\textwidth]{figures/Meson/MesonTrigger/EtaToPi0_data_CombinedBinShiftCorrection.eps}
% % 			\caption{By diving the 'y'-bin values before and after the bin shift procedure, the size of bin width corrections for the combined \EtaToPi-ratios for this analysis can be read off from these histograms for the different \EtaToPi-ratios obtained for the different triggers in this analysis (left) as well as for the combination (right).}
% % 			\label{fig:binShiftEtaPi0}
% % 		\end{figure}
% %
% %
%

%%%%% original chapter from combination note
% \subsubsection{Correction for Finite Bin Width}
% \label{sec:BinShift}
% The neutral meson spectra are binned in transverse momentum while the bin width of those slices increases with rising \pT. The underlying spectrum, however, is steeply falling with increasing transverse momenta.
% Thus, the yield measured in a \pT interval does not correspond to the yield which is measured at the center of a \pT bin.
% Therefore, the true \pT value for the yield, which has been measured in the given slice, needs to be determined \cite{Lafferty1994cj}.
% This can be done with two different approaches. Either shifting the data points horizontally in \pT, such that they correspond to the true \pT for a given bin or by shifting the data points vertically to the true yield for the \pT at the bin center.
% Both methods depend on the same underlying model assumptions. For the spectra, the points are shifted in the \pT direction (also referred to as 'x'-direction) while for the calculation of \EtaToPi-ratio, they will be shifted in the 'y'-direction.
% The shifting in the vertical direction for the \EtaToPi-ratio has been chosen due to the fact that the spectra for the neutral pion and eta meson may have different shapes. Thus, for the same binning the points would shift by a different amount in the \pT direction, making it harder to correctly calculate the \EtaToPi-ratio.
%
% The spectra have been shifted assuming the Tsallis function \cite{Tsallis1987eu} as an approximation of the underlying spectrum.
% 		\begin{equation}
% 		 	\frac{\mbox{d}^2N}{\mbox{d}y \mbox{d}p_{\mbox{\tiny T}}} = \frac{(n_{\mbox{\tiny Tsallis}}-1)(n_{\mbox{\tiny Tsallis}}-2)}{n_{\mbox{\tiny Tsallis}}T [n_{\mbox{\tiny Tsallis}}T + m(n_{\mbox{\tiny Tsallis}}-2)]} \times A \times p_{\mbox{\tiny T}} \times \left(1 + \frac{m_{T}-m}{n_{\mbox{\tiny Tsallis}}T}\right)^{-n_{\mbox{\tiny Tsallis}}}.
% 			\label{eq:Tsallis}
% 		\end{equation}
% The parameters $m$ and $m_{T} = \sqrt{m^2 + p_{T}^2}$ correspond to the particle mass and the transverse mass, while $A$, $T$ and $n_{\mbox{\tiny Tsallis}}$ are the free parameters of the fit function.
% In order to be able to combine different measurements in \acs{ALICE}, this bin width correction is calculated based on the combined spectra and then applied to the individual spectra based on the combined fit.
% Moreover, this allows to constrain the fit even more as each individual measurement covers a slightly different transverse momentum region, resulting in smaller correction factors.
% Hence, the described bin width corrections are already applied for all quantities/spectra shown in the following \hyperref[sec:combResults]{Section~\ref*{sec:combResults}}.
%
 \hyperref[fig:binShift]{Figure~\ref*{fig:binShift}} shows the bin width corrections for the final combined neutral pion and eta meson spectra. Corrections are of the order of 1\% throughout the whole transverse momentum region covered.

 		\begin{figure}[h]
 			\centering
 			\includegraphics[width=0.49\textwidth]{figures/Combination/BinShiftCorrection_Pi0.pdf}
 			\includegraphics[width=0.49\textwidth]{figures/Combination/BinShiftCorrection_Eta.pdf}
 			\caption{By diving the 'x'-bin values before and after the bin shift procedure, the size of bin width corrections for the combined final neutral pion and eta meson spectra can be read off from these histograms.}
 			\label{fig:binShift}
 		\end{figure}

% The next \hyperref[fig:binShiftEtaPi0]{Figure~\ref*{fig:binShiftEtaPi0}} visualizes the bin width corrections for the final combined \EtaToPi-ratios for \acs{EMCal} as well as \acs{PCM}-\acs{EMCal}. Note that corrections are now positive. Instead of shifting in 'x' as in the cases before, the bin shift for the \EtaToPi-ratio has been performed in 'y'. Thus, bin shift correction factors are now above one. \hyperref[fig:binShiftEtaPi0]{Figure~\ref*{fig:binShiftEtaPi0}} shows that corrections are of the order of 1\% almost throughout the whole transverse momentum region covered and go up to $\approx$2\% for low momenta in the latter case.
%
% 		\begin{figure}[h]
% 			\centering
% 			\includegraphics[width=0.49\textwidth]{figures/Combination/EtaToPi0_data_CombinedBinShiftCorrection.eps}
% 			\includegraphics[width=0.49\textwidth]{figures/Combination/EtaToPi0_data_CombinedBinShiftCorrection_EMCAL.eps}
% 			\caption{By diving the 'y'-bin values before and after the bin shift procedure, the size of bin width corrections for the combined final \EtaToPi-ratios can be read off from these histograms. Please note the different applied shift direction (in 'y' instead of in 'x') compared to \hyperref[fig:binShift]{Figure~\ref*{fig:binShift}}.}
% 			\label{fig:binShiftEtaPi0}
% 		\end{figure}
%
%
\subsection{Final Combined Results}
\label{sec:combResults}

		\begin{figure}[h]
	 		\centering
	 			\includegraphics[width=0.49\textwidth]{figures/Combination/Pi0_InvXSectionCompAllSystems_Comb.pdf}
	 			\includegraphics[width=0.49\textwidth]{figures/Combination/Eta_InvXSectionCompAllSystems_Comb.pdf}
	 			\caption{The invariant cross section for $\pi^0$ (left) and $\eta$ (right) measured with the different reconstruction techniques as well as the combined result. A TCM parametrisation to the combined spectrum is also shown.}
	 			\label{fig:CombSpectrumAndIndividual}
	 		\end{figure}

\hyperref[fig:CombSpectrumAndIndividual]{Figure~\ref*{fig:CombSpectrumAndIndividual}} shows the bin shift corrected invariant cross section for $\pi^0$ (left) and $\eta$ (right) measured in pp collisions at $\sqrt{s} = 13$ TeV for the different reconstruction techniques as well as the combined Spectrum. As this is the main result of the Analysis these spectra will be described in greaer detail in this section. Starting with the uncertainties of the measurement and comparing it to different parametrizations we will later also make comparisons to theory calculations.

 \hyperref[fig:combTotErr]{Figure~\ref*{fig:combTotErr}} shows the relative combined statistical and systematic errors for $\pi^0$, $\eta$ and \EtaToPi after the combination has been performed according to the procedure described in \hyperref[sec:combMeasure]{Section~\ref*{sec:combMeasure}} using the respective weights shown in \hyperref[fig:Weights]{Figure~\ref*{fig:Weights}}.

 		\begin{figure}[h]
 			\centering
 			\includegraphics[width=0.49\textwidth]{figures/Combination/Pi0_Reldecomp.pdf}
 			\includegraphics[width=0.49\textwidth]{figures/Combination/Eta_Reldecomp.pdf}\\
 			\includegraphics[width=0.49\textwidth]{figures/Combination/EtaToPi0_Reldecomp.pdf}
 			\caption{Relative total, statistical and systematic errors for $\pi^0$, $\eta$ and \EtaToPi, after combinations of available measurements in the given system have been performed.}
 			\label{fig:combTotErr}
 		\end{figure}
%
% %The following \hyperref[fig:Pi0_Eta_diffComb]{Figure~\ref*{fig:Pi0_Eta_diffComb}} is actually analogue to \hyperref[fig:Pi0_Eta_diff]{Figure~\ref*{fig:Pi0_Eta_diff}}, but shows in addition the combined spectra of all systems which has already been bin-shifted in 'x'-direction according to \hyperref[sec:BinShift]{Section~\ref*{sec:BinShift}}.
%
 The combined spectra have been fitted with the two-component model fit (\acs{TCM}), developed by A. Bylinkin and A. Rostovtsev \cite{Bylinkin2012bz,Bylinkin2011yi,Bylinkin2015xya},
 		\begin{eqnarray}
 			E \frac{\mbox{d}^3 \sigma}{\mbox{d}p^3} = A_e \exp{\frac{-(\sqrt{p_{\text{\tiny T}}^2 + M^2} - M)}{T_{2}}} + \frac{A}{\left( 1+ \frac{p_{\text{\tiny T}}^2}{T^2 n}\right)^{-n}}
 			\label{eq:TwoComponentModel}
 		\end{eqnarray}
 which can be looked at in \hyperref[fig:Fits]{Figure~\ref*{fig:Fits}}. The resulting parameters for the $\pi^0$ and $\eta$ are summarized in the following \hyperref[tab:fitParam]{Table~\ref*{tab:FitParam}}.
 %Moreover, Tsallis functions \cite{Tsallis1987eu} have been fitted to the spectra using \hyperref[eq:Tsallis]{Equation~\ref*{eq:Tsallis}} which may also be found in that table. \TODO{Combination}{Update Table!!}

 \renewcommand{\arraystretch}{1.3}
 \begin{table}[t]%[h!]
   \begin{center}
   \scalebox{0.92}{
     \begin{tabular}{c||cccccc}
      \hline
      TCM & $A_{e}$ (pb~GeV$^{-2}c^{3}$) & $T_{e}$ (GeV) & $A$ (pb~GeV$^{-2}c^{3}$) & $T$ (GeV) & $n$ & $\chi^2_{\rm red}$ \\
      \hline\hline
      $\pi^{0}$ & (5.23$\pm$3.05)$\times10^{11}$ & 0.154$\pm$0.028 & (2.68$\pm$0.68)$\times10^{10}$ & 0.667$\pm$0.033 & 2.995$\pm$0.013 & 0.07 \\
      \hline
      $\eta$ & (6.93$\pm$4.59)$\times10^{9}$ & 0.168$\pm$0.049 & (2.69$\pm$0.59)$\times10^{9}$ & 0.848$\pm$0.041 & 2.944$\pm$0.019 & 0.33\\
      \hline
%      \multicolumn{6}{c}{}\\
%      %
%      %
%      \hline
%      Tsallis & \multicolumn{2}{c}{$C$ (pb)} & $T$ (GeV) & \multicolumn{2}{c}{$n$} & $\chi^2_{\rm red}$ \\
%      \hline\hline
%      $\pi^{0}$ & \multicolumn{2}{c}{(2.46$\pm$0.18)$\times10^{11}$} & 0.121$\pm$0.004 & \multicolumn{2}{c}{6.465$\pm$0.042} & 0.57 \\
%      \hline
%      $\eta$ & \multicolumn{2}{c}{(1.56$\pm$0.19)$\times10^{10}$} & 0.221$\pm$0.012 & \multicolumn{2}{c}{6.560$\pm$0.113} & 0.59\\
%      \hline
     \end{tabular}}
     \caption{Parameters of the fits to the $\pi^{0}$ and $\eta$ invariant differential cross sections using the TCM fit \cite{Bylinkin2012bz,Bylinkin2011yi,Bylinkin2015xya} from Eq. \ref{eq:TwoComponentModel} as well as using a Tsallis fit \cite{Tsallis1987eu} from Eq.~\ref{eq:Tsallis}.}
     \label{tab:FitParam}
   \end{center}
 \end{table}
 \renewcommand{\arraystretch}{1.0}

 		\begin{figure}[h]
 		\centering
 			\includegraphics[width=0.49\textwidth]{figures/Combination/Pi0_RatioOfCombToCombFit.pdf}
 			\includegraphics[width=0.49\textwidth]{figures/Combination/Eta_RatioOfCombToCombFit_PP13TeV.pdf}
 			\caption{Two-component model (\acs{TCM}) fits for the pion (left) and eta (right) to the fully combined ALICE results at \s~=~13~TeV. }
 				%The fit parameters are given in the plots as well as in \hyperref[tab:FitParam]{Table~\ref*{tab:FitParam}}. As there are just a limited number of points with relatively bigger errors (in comparison to pion) available for the eta meson, the \acs{TCM} fit ends up with bigger errors in that case, but still converges with physical values. In blue and green, the two different components of \acs{TCM} have been decomposed. The red line displays the full \acs{TCM} fit, which is the sum of blue and green.}
 			\label{fig:Fits}
 		\end{figure}

% 		\begin{figure}[h]
% 		\centering
% 			\includegraphics[width=0.49\textwidth]{figures/Combination/ComparisonWithFit_Tsallis_Pi0_8TeV.eps}
% 			\includegraphics[width=0.49\textwidth]{figures/Combination/ComparisonWithFit_Tsallis_Eta_8TeV.eps}
% 			\caption{Tsallis \cite{Tsallis1987eu} fits for the pion (left) and eta (right) to the fully combined ALICE results at \s~=~8~TeV. The fit parameters are given in the plots as well as in \hyperref[tab:FitParam]{Table~\ref*{tab:FitParam}}.}
% 			\label{fig:FitsTsallis}
% 		\end{figure}
%
 The ratios of the single measurements to the combined TCM fits are given in \hyperref[fig:SingleRatiosToCombine]{Figure~\ref*{fig:SingleRatiosToCombine}} for the $\pi^0$ and $\eta$ mesons. A good agreement within the respective uncertainties can be stated.
%
 		\begin{figure}[h]
 		\centering
 			\includegraphics[width=0.49\textwidth]{figures/Combination/Pi0_RatioOfIndividualMeasToCombFit_PP13TeV.pdf}
 			\includegraphics[width=0.49\textwidth]{figures/Combination/Eta_RatioOfIndividualMeasToCombFit_PP13TeV.pdf}
 			\caption{Ratios of the single measurements to the \acs{TCM} fits of the fully combined results for the neutral pion (left) and the eta meson (right).}
 			\label{fig:SingleRatiosToCombine}
 		\end{figure}

 The ratios of the combined measurements to the combined fits are given in \hyperref[fig:RatioCombine]{Figure~\ref*{fig:RatioCombine}} for the $\pi^0$ and $\eta$, respectively.
%
% 	\begin{figure}[h]
% 	\centering
% 		\includegraphics[width=0.49\textwidth]{figures/Combination/Pi0_RatioOfCombToCombFit_PP8TeV.eps}
% 		\includegraphics[width=0.49\textwidth]{figures/Combination/Eta_RatioOfCombToCombFit_PP8TeV.eps}
% 		\caption{Ratios of the combined results to the \acs{TCM} fits of the fully combined results for the neutral pion (left) and the eta meson (right).}
% 		\label{fig:RatioCombine}
% 	\end{figure}
%
%
% The ratios of the single measurements to the combined Tsallis fits are given in \hyperref[fig:SingleRatiosToCombineTsallis]{Figure~\ref*{fig:SingleRatiosToCombineTsallis}} for the $\pi^0$ and $\eta$ mesons.
%
% 	\begin{figure}[h]
% 	\centering
% 		\includegraphics[width=0.45\textwidth]{figures/Combination/Pi0_RatioOfIndividualMeasToCombTsallisFit_PP.eps}
% 		\includegraphics[width=0.45\textwidth]{figures/Combination/Eta_RatioOfIndividualMeasToCombTsallisFit_PP.eps}
% 		\caption{Ratios of the single measurements to the Tsallis fits of the fully combined results for the neutral pion (left) and the eta meson (right).}
% 		\label{fig:SingleRatiosToCombineTsallis}
% 	\end{figure}
%
% The ratios of the combined measurements to the combined Tsallis fits are given in \hyperref[fig:RatioCombineTsallis]{Figure~\ref*{fig:RatioCombineTsallis}} for the $\pi^0$ and $\eta$, respectively.
%
% 	\begin{figure}[h]
% 	\centering
% 		\includegraphics[width=0.45\textwidth]{figures/Combination/Pi0_RatioOfCombToCombTsallisFit_PP8TeV.eps}
% 		\includegraphics[width=0.45\textwidth]{figures/Combination/Eta_RatioOfCombToCombTsallisFit_PP8TeV.eps}
%      	\caption{Ratios of the combined results to the Tsallis fits of the fully combined results for the neutral pion (left) and the eta meson (right).}
% 		\label{fig:RatioCombineTsallis}
% 	\end{figure}
%
% An overview over the ratios of the combined measurements to the combined TCM, Tsallis, mod. Hagedorn and power law fits are given in \hyperref[fig:RatioCombineOverview]{Figure~\ref*{fig:RatioCombineOverview}} for the $\pi^0$ and $\eta$, respectively.
%
% 		\begin{figure}[h]
% 		\centering
% 			\includegraphics[width=0.49\textwidth]{figures/Combination/Pi0_RatioOfCombToDifferentFits_PP8TeV.eps}
% 			\includegraphics[width=0.49\textwidth]{figures/Combination/Eta_RatioOfCombToDifferentFits_PP8TeV.eps}
% 			\caption{Ratios of the combined results to TCM, Tsallis, mod. Hagedorn and power law fits of the fully combined results for the neutral pion (left) and the eta meson (right).}
% 			\label{fig:RatioCombineOverview}
% 		\end{figure}
%
% The following \hyperref[fig:Pi0Final]{Figure~\ref*{fig:Pi0Final}} and \hyperref[fig:EtaFinal]{Figure~\ref*{fig:EtaFinal}} show the combined results for $\pi^0$ and $\eta$, together with comparisons to theory calculations and different Monte Carlo generators. The \acs{TCM} fits to the combined results are also plotted.
%
% 		\begin{figure}
% 		    \vspace{-0.8cm}
% 			\centering
% 			\includegraphics[width=0.9\textwidth]{figures/Combination/Pi0_InvXSectionWithRatios_Paper.eps}
% 			\vspace{-0.2cm}
% 			\caption{The plot shows the fully combined invariant neutral pion cross section together with the corresponding \acs{TCM} fit as well as ratios to it. Furthermore, the given NLO calculation using \acs{PDF}:MSTW08 and \acs{FF}:DSS14 is compared to the fit as well as predictions from PYTHIA8.210 'Tune 4C' and 'Monash 2013'.}
% 			\label{fig:Pi0Final}
% 		\end{figure}
%
% 		\begin{figure}
% 		    \vspace{-0.8cm}
% 			\centering
% 			\includegraphics[width=0.9\textwidth]{figures/Combination/Eta_InvXSectionWithRatios_Paper.eps}
% 			\vspace{-0.2cm}
% 			\caption{The plot shows the fully combined invariant eta meson cross section together with the corresponding \acs{TCM} fit as well as ratios to it. Furthermore, the given NLO calculation using \acs{PDF}:CTEQ6M5 and \acs{FF}:AESSS is compared to the fit as well as predictions from PYTHIA8.210 'Tune 4C' and 'Monash 2013'.}
% 			\label{fig:EtaFinal}
% 		\end{figure}
%
%
% \hyperref[fig:EtaToPi0Final]{Figure~\ref*{fig:EtaToPi0Final}} visualizes the fully combined \EtaToPi ratio at \s~=~8~TeV and compares with \acs{ALICE} results at two different center of mass energies of \s~=~2.76~TeV (see \cite{anaNoteComb} as well as \cite{anaNoteEMC2760GeV}, \cite{anaNotePCMEMC2760GeV} and \cite{anaNote1052} for references) and \s~=~7~TeV (see \cite{anaNote7TeV2012} for reference) which are found to be consistent within errors.
%
% 		\begin{figure}[h]
% 			\centering
% 			\includegraphics[width=0.9\textwidth]{figures/Combination/EtaToPi0_Paper.eps}
% 			\caption{Combined \EtaToPi ratio at \s~=~8~TeV compared to the ratios obtained at \s~=~2.76~TeV and \s~=~7~TeV.}
% 			\label{fig:EtaToPi0Final}
% 		\end{figure}
%
%
\subsection{Comparison to Theory Calculations}
 \hyperref[fig:Pi0_Theory]{Figure~\ref*{fig:Pi0_Theory}} shows the ratio of the combined neutral pion cross section, the PYTHIA8.210 \TODO{Combination}{which PYTHIA?} calculation as well as a NLO calculation using \acs{FF}:DSS14 to the combined fit.

 		\begin{figure}[h]
 			\centering
 			\includegraphics[width=0.49\textwidth]{figures/Combination/Pi0_RatioTheoryToData_PP.pdf}
 			\caption{\newline This plot shows ratios of the given NLO calculation using \acs{PDF}:MSTW08 and \acs{FF}:DSS14 as well as s prediction from PYTHIA8.210 'Monash 2013' to the combined \acs{TCM} fit.}
 			\label{fig:Pi0_Theory}
 		\end{figure}

\hyperref[fig:Eta_Theory]{Figure~\ref*{fig:Eta_Theory} (left)} shows the ratio of the combined eta meson cross section, the PYTHIA8.210 calculation as well as a NLO calculation using \acs{PDF}:CTEQ6M5 and \acs{FF}:AESSS to the combined fit.
%
% The combined \EtaToPi-ratio is shown in \hyperref[fig:EtaToPi0_Theory]{Figure~\ref*{fig:EtaToPi0_Theory} (right)} with ALICE measurements at different center of mass energies and other related experiments, together a $m_{T}$ scaled \EtaToPi-ratio computed from combined neutral pion spectrum at \s~=~8~TeV by applying the empirical scaling law.
%
% The combined \EtaToPi-ratio is shown in \hyperref[fig:EtaToPi0_Paper]{Figure~\ref*{fig:EtaToPi0_Paper}} with PYTHIA~8.210 predictions as well as NLO calculations using \acs{PDF}:CTEQ6M5, together with \acs{FF}:DSS07 for the neutral pion and \acs{FF}:AESSS for the eta meson.
%
% The combined \EtaToPi-ratio alone is shown in \hyperref[fig:EtaToPi0_Combined]{Figure~\ref*{fig:EtaToPi0_Combined}}.
%
 		\begin{figure}[h]
 			\centering
 			\includegraphics[width=0.49\textwidth]{figures/Combination/Eta_RatioTheoryToData_PP.pdf}
 			\caption{\newline (left) This plot shows ratios of the given NLO calculation using \acs{PDF}:CTEQ6M5 and \acs{FF}:AESSS as well as a prediction from PYTHIA8.210 'Monash 2013' to the combined \acs{TCM} fit.}
 			\label{fig:Eta_Theory}
 			\label{fig:EtaToPi0_Theory}
 		\end{figure}
%
% 		\begin{figure}[h]
% 			\centering
% 			\includegraphics[width=0.49\textwidth]{figures/Combination/EtaToPi0_Theory_Paper.eps}
% 			\includegraphics[width=0.49\textwidth]{figures/Combination/EtaToPi0_Theory_Paper_LIN.eps}
% 			\caption{These plots (left with x-axis in log-scale, right in linear) display the combined \EtaToPi-ratio as well as a NLO prediction using \acs{PDF}:CTEQ6M5, \acs{FF}:DSS07 for the neutral pion and \acs{FF}:AESSS for the eta meson. The predictions from the two different PYTHIA tunes are also shown.}
% 			\label{fig:EtaToPi0_Paper}
% 		\end{figure}
%
% 		\begin{figure}[h]
% 			\centering
% 			\includegraphics[width=0.49\textwidth]{figures/Combination/EtaToPi0_Combined.eps}
% 			\includegraphics[width=0.49\textwidth]{figures/Combination/EtaToPi0_Combined_LIN.eps}
% 			\caption{The combined \EtaToPi-ratio is shown in log- and linear-scale without any further additions.}
% 			\label{fig:EtaToPi0_Combined}
% 		\end{figure}
%
%
  The combined $\pi^0$ and $\eta$ spectra are shown in \hyperref[fig:Pi0Eta_Theory]{Figure~\ref*{fig:Pi0Eta_Theory}} together with PYTHIA as well as NLO pQCD predictions. The TCM and Tsallis fits are also plotted. \hyperref[fig:Pi0Eta]{Figure~\ref*{fig:Pi0Eta}} shows comparable figures and contains only selected parts of \hyperref[fig:Pi0Eta_Theory]{Figure~\ref*{fig:Pi0Eta_Theory}}.

 		\begin{figure}[h]
 			\centering
 			\includegraphics[width=0.49\textwidth]{figures/Combination/InvXSection_Pi0_Eta_Theory.pdf}
 			\caption{Combined $\pi^0$ and $\eta$ spectra with fits and theory predictions.}
 			\label{fig:Pi0Eta_Theory}
 		\end{figure}

%  	    \begin{figure}[h]
% 			\centering
% 			\includegraphics[width=0.49\textwidth]{figures/Combination/InvXSection_Pi0_Eta_Fits.eps}
% 			\includegraphics[width=0.49\textwidth]{figures/Combination/InvXSection_Pi0_Eta.eps}
% 			\caption{Combined $\pi^0$ and $\eta$ spectra with fits on left side and without on right side.}
% 			\label{fig:Pi0Eta}
% 		\end{figure}
%
%
\subsubsection{\texorpdfstring{Comparison with $m_{\rm T}$ scaling prediction}{Comparison with mT scaling prediction}}
 The validity of $m_{\rm T}$ scaling was tested by comparing the measured $\eta/\pi^{0}$ ratio with the ratio of the $\eta$ spectrum, derived from the $m_{\rm T}$-scaled TCM parametrization of the $\pi^0$ spectrum, to the fit of the measured $\pi^0$ spectrum, which is shown in \hyperref[fig:mTscaling]{Figure~\ref*{fig:mTscaling}}.
 The quoted $\pi^0$ fit parameters from \hyperref[tab:FitParam]{Table~\ref*{tab:FitParam}} were used as input for $m_{\rm T}$ scaling, further replacing the $\pi^0$ mass by the $\eta$ mass and using the normalization ratio $C^{\eta}/C^{\pi^0}=0.47$.

  	    \begin{figure}[h]
 			\centering
% 			\includegraphics[width=0.75\textwidth]{figures/Combination/EtaToPi0_mT_ratio.eps}
			\includegraphics[width=0.75\textwidth]{figures/Combination/EtaToPi0_Theory.pdf}
 			\caption{.}
 			\label{fig:mTscaling}
 		\end{figure}

As expected from other measurements of the $\eta/\pi^{0}$-ratio, the ratio is constant above $p_{T} \approx$ 4 GeV/c. A constant function fitted in a range from $4 \text{GeV}/c \leq p_{T} < 30 \text{GeV}/c$ gives a value of $C^{\eta}/C^{\pi^0}=0.47 \pm 0.00 (stat) \pm 0.01 (sys)$.\\
Additionally the prediction from Pythia is shown in  \hyperref[fig:mTscaling]{Figure~\ref*{fig:mTscaling}}. Pythia underestimates the $\eta/\pi^{0}$-ratio at high $p_{\text{T}}$ and overestimated it at low $p_{\text{T}}$.
% For each single point below 3.5~\GeVc in \hyperref[fig:mTscaling]{Figure~\ref*{fig:mTscaling}}, the following deviations in n$\sigma$ have been found:
%
% \begin{itemize}
% \item \pT: 0.65; ytrue: 1; ymeas: 0.448; y\_sig: 0.214; y\_sig\_rel: 0.479; \textbf{n$\sigma$: 2.56}
% \item \pT: 0.95; ytrue: 1; ymeas: 0.434; y\_sig: 0.139; y\_sig\_rel: 0.321; \textbf{n$\sigma$: 4.04}
% \item \pT: 1.25; ytrue: 1; ymeas: 0.702; y\_sig: 0.112; y\_sig\_rel: 0.160; \textbf{n$\sigma$: 2.63}
% \item \pT: 1.55; ytrue: 1; ymeas: 0.678; y\_sig: 0.089; y\_sig\_rel: 0.132; \textbf{n$\sigma$: 3.57}
% \item \pT: 1.85; ytrue: 1; ymeas: 0.906; y\_sig: 0.098; y\_sig\_rel: 0.108; \textbf{n$\sigma$: 0.94}
% \item \pT: 2.20; ytrue: 1; ymeas: 0.755; y\_sig: 0.067; y\_sig\_rel: 0.089; \textbf{n$\sigma$: 3.61}
% \item \pT: 2.60; ytrue: 1; ymeas: 0.848; y\_sig: 0.068; y\_sig\_rel: 0.080; \textbf{n$\sigma$: 2.21}
% \item \pT: 3.00; ytrue: 1; ymeas: 0.881; y\_sig: 0.077; y\_sig\_rel: 0.087; \textbf{n$\sigma$: 1.52}
% \item \pT: 3.40; ytrue: 1; ymeas: 0.933; y\_sig: 0.088; y\_sig\_rel: 0.094; \textbf{n$\sigma$: 0.75}
% \end{itemize}
%
% In total, a disagreement of measured $\eta/\pi^{0}$ ratio with respect to $m_{\rm T}$ scaling hypothesis of $6.2\sigma$ for \pT~$<$~3.5\GeVc has been found, clearly breaking this empirical scaling law for the quoted momentum region.
%
%




\begin{appendix}

\section{Software Version \& Input Data}
\label{sec:software}
In the following, the software versions and input files are listed which have been used to generate the results shown in this note:

\begin{enumerate}
\item[-] AliPhysics \cite{AliPhysicsPCM2015} version:
\begin{center}
Software tags:\\
data: ''''\\
MB MC:  ''''\\
JetJet MC: ''''
\end{center}

\item[-] Input data / MC:
\begin{center}
data trains: GA\_pp\_AOD  \\
MB MC trains: GA\_pp\_MC\_AOD \\
JetJet MC trains: GA\_pp\_MC\_AOD
\end{center}

\item[-] Offline software (PCM AnalysisSoftware) \cite{ConvSoftware} used:
\begin{center}
Tag ''XXXX' in master branch of quoted git repository
\end{center}

\item[-] Cocktail output (Grid running) \cite{ConvSoftware} used:
\begin{center}
Software tag ''''\\
train: MBGen\_pp
\end{center}

\end{enumerate}

\section{List of Good Runs}
\label{sec:goodRuns}


 \subsection{LHC16d anchored to: LHC17f6, LHC17f6\_extra, LHC17f6\_extra2}

 \textbf{PCM:}\\
252235, 252238, 252248, 252271, 252310, 252313, 252315, 252317, 252319, 252322, 252325, 252326, 252330, 252332, 252336, 252368, 252370, 252371, 252374, 252375\\

 \subsection{LHC16g anchored to: LHC17d17, LHC17d17\_extra, LHC17d17\_extra2}

 \textbf{EDC:}\\
254128, 254147, 254148, 254149, 254174, 254199, 254204, 254205, 254293, 254302, 254303, 254304, 254330, 254331, 254332\\

 \textbf{PCMEDC:}\\
254128, 254147, 254148, 254149, 254174, 254199, 254204, 254205, 254293, 254302, 254303, 254304, 254330, 254331, 254332\\

 \textbf{PCM:}\\
254128, 254147, 254148, 254149, 254174, 254175, 254178, 254193, 254199, 254204, 254205, 254293, 254302, 254303, 254304, 254330, 254331, 254332\\

 \textbf{PHOS:}\\
254128, 254147, 254148, 254149, 254174, 254175, 254178, 254193, 254199, 254204, 254205, 254304, 254330, 254331\\

 \subsection{LHC16h anchored to: LHC17f5, LHC17f5\_extra, LHC17f5\_extra2}

 \textbf{EDC:}\\
254604, 254606, 254608, 254621, 254629, 254630, 254632, 254640, 254644, 254646, 254648, 254649, 254651, 254652, 254653, 254654, 255248, 255249, 255251, 255252, 255253, 255255, 255256, 255275, 255276, 255350, 255351, 255352, 255418, 255419, 255420, 255421, 255440, 255463, 255465, 255466, 255467\\

 \textbf{PCMEDC:}\\
254604, 254606, 254621, 254629, 254630, 254632, 254640, 254644, 254646, 254648, 254649, 254651, 254652, 254653, 254654, 255249, 255251, 255252, 255253, 255255, 255256, 255275, 255276, 255350, 255351, 255352, 255418, 255419, 255420, 255421, 255463, 255465, 255466, 255467\\

 \textbf{PCM:}\\
254418, 254419, 254422, 254604, 254606, 254608, 254621, 254629, 254630, 254632, 254640, 254644, 254646, 254648, 254649, 254651, 254652, 254653, 254654, 255079, 255082, 255085, 255086, 255091, 255111, 255154, 255159, 255162, 255167, 255171, 255173, 255174, 255176, 255177, 255240, 255242, 255247, 255248, 255249, 255251, 255252, 255253, 255255, 255256, 255275, 255276, 255280, 255283, 255350, 255351, 255352, 255398, 255402, 255407, 255415, 255418, 255419, 255420, 255421, 255440, 255442, 255447, 255463, 255465, 255466, 255467, 255469\\

 \textbf{PHOS:}\\
255275, 254983, 254984, 255242, 254604, 254606, 254608, 255249, 255251, 255252, 255253, 255255, 255256, 255177, 254621, 255466, 255008, 255009, 255010, 254629, 254630, 254378, 255111, 255276, 254381, 255407, 254640, 255240, 255154, 255283, 254644, 255467, 254646, 255159, 254648, 254649, 254394, 254395, 254396, 254653, 254654, 255167, 255171, 255173, 255176, 254476, 255180, 255182, 255280, 255440, 255442, 254419, 254422, 255447, 254479, 255162, 255248, 255398, 255079, 255465, 255082, 255463, 255085, 255086, 255247, 255091, 255350, 255351, 255402, 255011, 255415, 255418, 255419, 255420\\

 \subsection{LHC16i anchored to: LHC17d3, LHC17d3\_extra, LHC17d3\_extra2}

 \textbf{EDCtrigger:}\\
255539, 255540, 255541, 255542, 255543, 255577, 255582, 255583, 255591, 255592, 255614, 255615, 255616, 255617, 255618\\

 \textbf{EDC:}\\
255539, 255540, 255541, 255542, 255543, 255577, 255582, 255583, 255591, 255592, 255614, 255615, 255616, 255617, 255618\\

 \textbf{PCMEDC:}\\
255539, 255540, 255541, 255542, 255543, 255577, 255582, 255583, 255591, 255592, 255614, 255615, 255616, 255617, 255618\\

 \textbf{PCM:}\\
255539, 255540, 255541, 255542, 255543, 255577, 255582, 255583, 255591, 255592, 255614, 255615, 255616, 255617, 255618\\

 \textbf{PHOS:}\\
255615, 255614, 255592, 255591, 255582, 255577, 255543, 255542, 255540, 255538, 255537, 255535, 255534, 255533, 255583, 255539\\

 \subsection{LHC16j anchored to: LHC17e5, LHC17e5\_extra, LHC17e5\_extra2}

 \textbf{EDCtrigger:}\\
256207, 256231, 256281, 256282, 256283, 256284, 256289, 256290, 256292, 256295, 256297, 256298, 256299, 256302, 256307, 256309, 256311, 256356, 256357, 256361, 256362, 256364, 256365, 256366, 256371, 256372, 256373, 256415, 256417, 256418, 256420\\

 \textbf{EDC:}\\
256207, 256223, 256225, 256231, 256281, 256282, 256283, 256284, 256289, 256290, 256292, 256295, 256297, 256298, 256299, 256302, 256307, 256309, 256311, 256356, 256357, 256361, 256362, 256363, 256364, 256365, 256366, 256371, 256372, 256373, 256415, 256417, 256418, 256420 \\

 \textbf{PCMEDC:}\\
256207, 256223, 256225, 256231, 256281, 256282, 256283, 256284, 256289, 256290, 256292, 256295, 256297, 256298, 256299, 256302, 256307, 256309, 256311, 256356, 256357, 256361, 256362, 256364, 256365, 256366, 256371, 256372, 256373, 256415, 256417, 256418, 256420\\

 \textbf{PCM:}\\
256204, 256207, 256210, 256212, 256213, 256215, 256219, 256222, 256223, 256225, 256227, 256228, 256231, 256281, 256282, 256283, 256284, 256287, 256289, 256290, 256292, 256295, 256297, 256298, 256299, 256302, 256307, 256309, 256311, 256356, 256357, 256361, 256362, 256363, 256364, 256365, 256366, 256368, 256371, 256372, 256373, 256415, 256417, 256418, 256420\\

 \textbf{PHOS:}\\
256219, 256223, 256225, 256227, 256228, 256231, 256283, 256284, 256289, 256290, 256292, 256297, 256299, 256302, 256307, 256309, 256311, 256357, 256363, 256364, 256365, 256366, 256368, 256371, 256415, 256210, 256212, 256213, 256215, 256282, 256287, 256361, 256362, 256417, 256418\\

 \subsection{LHC16k anchored to: LHC18f1, LHC18f1\_extra}

 \textbf{EDCtrigger:}\\
258537, 258499, 258498, 258477, 258456, 258454, 258426, 258393, 258388, 258387, 258359, 258336, 258299, 258280, 258278, 258274, 258273, 258271, 258270, 258258, 258257, 258256, 258204, 258203, 258202, 258198, 258197, 258178, 258117, 258114, 258113, 258109, 258108, 258107, 258063, 258062, 258059, 258049, 258048, 258045, 258042, 258019, 258017, 258014, 258012, 257963, 257960, 257958, 257957, 257939, 257937, 257936, 257893, 257855, 257850, 257803, 257800, 257799, 257798, 257797, 257773, 257765, 257754, 257737, 257735, 257734, 257733, 257724, 257697, 257694, 257692, 257691, 257689, 257687, 257682, 257642, 257606, 257605, 257594, 257590, 257587, 257566, 257562, 257561, 257560, 257541, 257540, 257539, 257537, 257531, 257530, 257492, 257491, 257490, 257487, 257474, 257457, 257260, 257224, 257209, 257206, 257204, 257145, 257144, 257142, 257141, 257140, 257139, 257138, 257137, 257136, 257100, 257092, 257084, 257083, 257082, 257080, 257077, 257026, 257021, 257012, 257011, 256944, 256942, 256941, 256697, 256695, 256694, 256691, 256684, 256681, 256677, 256676, 256658, 256620, 256619, 256591, 256589, 256567, 256565, 256564, 256562, 256561, 256560, 256556, 256554, 256552, 256514, 256512, 256510, 256506, 256504\\

 \textbf{EDC:}\\
256504, 256506, 256510, 256512, 256552, 256554, 256556, 256560, 256561, 256562, 256564, 256565, 256567, 256589, 256619, 256620, 256658, 256676, 256677, 256681, 256684, 256691, 256694, 256695, 256697, 256941, 256942, 256944, 257012, 257021, 257026, 257071, 257077, 257080, 257082, 257083, 257084, 257092, 257100, 257136, 257137, 257138, 257139, 257140, 257141, 257142, 257144, 257145, 257204, 257206, 257209, 257224, 257260, 257318, 257320, 257322, 257364, 257433, 257457, 257474, 257487, 257490, 257491, 257492, 257530, 257531, 257537, 257539, 257540, 257541, 257560, 257561, 257562, 257566, 257587, 257590, 257594, 257605, 257606, 257630, 257642, 257682, 257684, 257687, 257688, 257689, 257691, 257692, 257694, 257697, 257724, 257725, 257733, 257734, 257735, 257737, 257754, 257765, 257773, 257797, 257798, 257799, 257800, 257803, 257804, 257850, 257853, 257855, 257893, 257936, 257937, 257939, 257957, 257958, 257960, 257963, 257979, 258012, 258014, 258017, 258019, 258039, 258041, 258042, 258045, 258048, 258049, 258059, 258060, 258062, 258063, 258107, 258108, 258109, 258113, 258114, 258117, 258178, 258197, 258198, 258202, 258203, 258204, 258256, 258257, 258258, 258270, 258271, 258273, 258274, 258278, 258299, 258336, 258359, 258393, 258426, 258454, 258456, 258477, 258499, 258537\\

 \textbf{PCMEDC:}\\
258537, 258499, 258477, 258456, 258454, 258426, 258393, 258359, 258336, 258299, 258278, 258274, 258273, 258271, 258270, 258258, 258257, 258256, 258204, 258203, 258202, 258198, 258197, 258178, 258117, 258114, 258113, 258109, 258108, 258107, 258063, 258062, 258059, 258049, 258048, 258045, 258042, 258019, 258017, 258014, 258012, 257963, 257960, 257958, 257957, 257939, 257937, 257936, 257893, 257855, 257850, 257803, 257800, 257799, 257798, 257797, 257773, 257765, 257754, 257737, 257735, 257734, 257733, 257724, 257697, 257694, 257692, 257691, 257689, 257687, 257682, 257642, 257606, 257605, 257594, 257590, 257587, 257566, 257562, 257561, 257560, 257541, 257540, 257539, 257537, 257531, 257530, 257492, 257491, 257490, 257487, 257474, 257457, 257433, 257364, 257322, 257320, 257318, 257260, 257224, 257209, 257206, 257204, 257145, 257144, 257142, 257141, 257140, 257139, 257138, 257137, 257136, 257100, 257092, 257084, 257083, 257082, 257080, 257077, 257026, 257021, 257012, 256944, 256942, 256941, 256697, 256695, 256694, 256691, 256684, 256681, 256677, 256676, 256658, 256620, 256619, 256589, 256567, 256565, 256564, 256562, 256561, 256560, 256556, 256554, 256552, 256512, 256510, 256506, 256504\\

 \textbf{PCM:}\\
256504, 256506, 256510, 256512, 256552, 256554, 256556, 256557, 256560, 256561, 256562, 256564, 256565, 256567, 256589, 256619, 256620, 256658, 256676, 256677, 256681, 256684, 256691, 256692, 256694, 256695, 256697, 256941, 256942, 256944, 257012, 257021, 257026, 257071, 257077, 257080, 257082, 257083, 257084, 257086, 257092, 257095, 257100, 257136, 257137, 257138, 257139, 257140, 257141, 257142, 257144, 257145, 257204, 257206, 257209, 257224, 257260, 257318, 257320, 257322, 257330, 257358, 257364, 257433, 257457, 257468, 257474, 257487, 257488, 257490, 257491, 257492, 257530, 257531, 257537, 257539, 257540, 257541, 257560, 257561, 257562, 257566, 257587, 257588, 257590, 257592, 257594, 257595, 257604, 257605, 257606, 257630, 257632, 257635, 257636, 257642, 257644, 257682, 257684, 257685, 257687, 257688, 257689, 257691, 257692, 257694, 257697, 257724, 257725, 257727, 257733, 257734, 257735, 257737, 257754, 257757, 257765, 257773, 257797, 257798, 257799, 257800, 257803, 257804, 257850, 257851, 257853, 257855, 257893, 257901, 257912, 257932, 257936, 257937, 257939, 257957, 257958, 257960, 257963, 257979, 257986, 257989, 257992, 258003, 258008, 258012, 258014, 258017, 258019, 258039, 258041, 258042, 258045, 258048, 258049, 258053, 258059, 258060, 258062, 258063, 258107, 258108, 258109, 258113, 258114, 258117, 258178, 258197, 258198, 258202, 258203, 258204, 258256, 258257, 258258, 258270, 258271, 258273, 258274, 258278, 258299, 258301, 258302, 258303, 258306, 258307, 258332, 258336, 258359, 258391, 258393, 258426, 258452, 258454, 258456, 258477, 258499, 258537\\

 \textbf{PHOS:}\\
257682, 257684, 257685, 257687, 257688, 257689, 257692, 257697, 257724, 257725, 257727, 257733, 257734, 257735, 257737, 257754, 257757, 257765, 257773, 257797, 257798, 257799, 257800, 257803, 257804, 257850, 257851, 257853, 257893, 257901, 257912, 257932, 257936, 257937, 257939, 257957, 257958, 257960, 257963, 257986, 257989, 257992, 258003, 258008, 258019, 258039, 258041, 258045, 258048, 258049, 258053, 258059, 258060, 258062, 258063, 258107, 258108, 258109, 258113, 258114, 258117, 258178, 258197, 258198, 258202, 258203, 258204, 258256, 258257, 258258, 258270, 258271, 258274, 258278, 258280, 258301, 258302, 258303, 258307, 258332, 258336, 258359, 258388, 258391, 258393, 258452, 258454, 258456, 258477, 258498, 258499\\

 \subsection{LHC16l anchored to: LHC18d8, LHC18d8\_extra}

 \textbf{EDC:}\\
258962, 258964, 259088, 259090, 259091, 259096, 259099, 259117, 259118, 259164, 259204, 259257, 259263, 259269, 259270, 259271, 259272, 259273, 259274, 259302, 259305, 259307, 259334, 259336, 259339, 259340, 259341, 259342, 259378, 259382, 259388, 259389, 259394, 259395, 259396, 259473, 259477, 259649, 259650, 259668, 259697, 259703, 259704, 259711, 259747, 259748, 259750, 259751, 259752, 259756, 259781, 259788, 259822, 259841, 259842, 259860, 259866, 259867, 259868, 259888\\

 \textbf{PCMEDC:}\\
258962, 258964, 259088, 259090, 259091, 259096, 259099, 259117, 259118, 259164, 259204, 259257, 259263, 259269, 259270, 259271, 259272, 259273, 259274, 259302, 259305, 259307, 259334, 259336, 259339, 259340, 259341, 259342, 259378, 259382, 259388, 259389, 259394, 259395, 259396, 259473, 259477, 259649, 259650, 259668, 259697, 259703, 259704, 259711, 259747, 259748, 259750, 259751, 259752, 259756, 259781, 259788, 259822, 259841, 259842, 259860, 259866, 259867, 259868, 259888\\

 \textbf{PCM:}\\
258962, 258964, 259086, 259088, 259090, 259091, 259096, 259099, 259117, 259118, 259162, 259164, 259204, 259257, 259261, 259263, 259264, 259269, 259270, 259271, 259272, 259273, 259274, 259302, 259303, 259305, 259307, 259334, 259336, 259339, 259340, 259341, 259342, 259378, 259381, 259382, 259388, 259389, 259394, 259395, 259396, 259473, 259477, 259649, 259650, 259668, 259697, 259700, 259703, 259704, 259705, 259711, 259713, 259747, 259748, 259750, 259751, 259752, 259756, 259781, 259788, 259789, 259822, 259841, 259842, 259860, 259866, 259867, 259868, 259888\\

 \textbf{PHOS:}\\
258883, 258884, 258885, 258886, 258889, 258890, 258919, 258920, 258921, 258923, 258926, 258931, 258962, 258964, 259086, 259088, 259090, 259096, 259099, 259117, 259162, 259164, 259204, 259261, 259263, 259264, 259269, 259270, 259271, 259272, 259273, 259274, 259302, 259303, 259305, 259307, 259334, 259336, 259339, 259340, 259341, 259342, 259378, 259381, 259389, 259394, 259395, 259473, 259477, 259649, 259650, 259668, 259697, 259703, 259704, 259705, 259711, 259713, 259747, 259750, 259751, 259752, 259756, 259788, 259789, 259822, 259841, 259842, 259860, 259866, 259867, 259868, 259888, 259961, 259979, 260010, 260011\\

 \subsection{LHC16o anchored to: LHC17d16, LHC17d16\_extra, LHC17d16\_extra2}

 \textbf{EDCtrigger:}\\
264035, 264033, 263985, 263984, 263981, 263979, 263978, 263977, 263923, 263917, 263916, 263905, 263866, 263863, 263861, 263803, 263793, 263792, 263790, 263786, 263785, 263784, 263744, 263743, 263741, 263739, 263738, 263737, 263691, 263690, 263689, 263682, 263663, 263662, 263657, 263654, 263653, 263652, 262778, 262777, 262776, 262768, 262760, 262727, 262719, 262717, 262713, 262708, 262706, 262705, 262635, 262632, 262628, 262624, 262594, 262593, 262583, 262574, 262572, 262571, 262570, 262569, 262567, 262563, 262533, 262532, 262492, 262451, 262450\\

 \textbf{EDC:}\\
262450, 262451, 262492, 262528, 262532, 262533, 262563, 262567, 262569, 262570, 262571, 262572, 262574, 262583, 262593, 262594, 262624, 262628, 262632, 262635, 262705, 262706, 262708, 262713, 262717, 262719, 262723, 262727, 262760, 262768, 262776, 262777, 262778, 262841, 262844, 262849, 262853, 262858, 263332, 263487, 263496, 263529, 263647, 263652, 263653, 263654, 263657, 263662, 263663, 263682, 263689, 263690, 263691, 263737, 263738, 263739, 263741, 263743, 263744, 263784, 263785, 263786, 263787, 263790, 263792, 263793, 263803, 263861, 263863, 263866, 263905, 263916, 263917, 263977, 263978, 263979, 263981, 263984, 263985, 264033, 264035\\

 \textbf{PCMEDC:}\\
264035, 264033, 263985, 263984, 263981, 263979, 263978, 263977, 263917, 263916, 263905, 263866, 263863, 263861, 263803, 263793, 263792, 263790, 263787, 263786, 263785, 263784, 263744, 263743, 263741, 263739, 263738, 263737, 263691, 263690, 263689, 263682, 263663, 263662, 263657, 263654, 263653, 263652, 263496, 263487, 263332, 262858, 262853, 262849, 262844, 262841, 262778, 262777, 262776, 262768, 262760, 262727, 262719, 262717, 262713, 262708, 262706, 262705, 262635, 262632, 262628, 262624, 262594, 262593, 262583, 262574, 262572, 262571, 262570, 262569, 262567, 262563, 262533, 262532, 262492, 262451, 262450\\

 \textbf{PCM:}\\
262450, 262451, 262487, 262489, 262490, 262492, 262528, 262532, 262533, 262537, 262563, 262567, 262568, 262569, 262570, 262571, 262572, 262574, 262578, 262583, 262593, 262594, 262624, 262628, 262632, 262635, 262705, 262706, 262708, 262713, 262717, 262719, 262723, 262725, 262727, 262760, 262768, 262776, 262777, 262778, 262841, 262842, 262844, 262847, 262849, 262853, 262855, 262858, 263332, 263487, 263490, 263496, 263497, 263529, 263647, 263652, 263653, 263654, 263657, 263662, 263663, 263682, 263689, 263690, 263691, 263737, 263738, 263739, 263741, 263743, 263744, 263784, 263785, 263786, 263787, 263790, 263792, 263793, 263803, 263810, 263861, 263863, 263866, 263905, 263916, 263917, 263920, 263977, 263978, 263979, 263981, 263984, 263985, 264033, 264035\\

 \textbf{PHOS:}\\
262395, 262396, 262397, 262398, 262399, 262418, 262419, 262422, 262423, 262424, 262425, 262426, 262428, 262430, 262450, 262451, 262487, 262489, 262490, 262492, 262528, 262532, 262533, 262537, 262563, 262567, 262568, 262569, 262570, 262571, 262572, 262574, 262578, 262583, 262593, 262624, 262628, 262632, 262635, 262705, 262706, 262708, 262713, 262717, 262723, 262725, 262727, 262768, 262776, 262777, 262778, 262842, 262844, 262847, 262849, 262853, 262855, 263487, 263490, 263497, 263529, 263647, 263652, 263654, 263657, 263662, 263663, 263682, 263689, 263690, 263691, 263737, 263739, 263741, 263743, 263784, 263785, 263786, 263787, 263790, 263792, 263803, 263810, 263813, 263823, 263824, 263829, 263830, 263861, 263863, 263866, 263905, 263916, 263917, 263920, 263977, 263978, 263979, 263981, 263984, 263985\\

 \subsection{LHC16p anchored to: LHC17d18, LHC17d18\_extra, LHC17d18\_extra2}

 \textbf{EDCtrigger:}\\
264347, 264346, 264345, 264336, 264312, 264306, 264305, 264281, 264279, 264277, 264267, 264266, 264265, 264264, 264262, 264261, 264260, 264259, 264238, 264235, 264233, 264232, 264198, 264197, 264190, 264188, 264168, 264164, 264139, 264138, 264137, 264129, 264110, 264109, 264086, 264085, 264082, 264078, 264076\\

 \textbf{EDC:}\\
264347, 264346, 264345, 264336, 264312, 264306, 264305, 264281, 264279, 264277, 264267, 264266, 264265, 264264, 264262, 264261, 264260, 264259, 264238, 264235, 264233, 264232, 264198, 264197, 264190, 264188, 264168, 264164, 264139, 264138, 264137, 264129, 264110, 264109, 264086, 264085, 264082, 264078, 264076\\

 \textbf{PCMEDC:}\\
264347, 264346, 264345, 264336, 264312, 264306, 264305, 264281, 264279, 264277, 264267, 264266, 264265, 264264, 264262, 264261, 264260, 264259, 264238, 264235, 264233, 264232, 264198, 264197, 264190, 264188, 264168, 264164, 264139, 264138, 264137, 264129, 264110, 264109, 264086, 264085, 264082, 264078, 264076\\

 \textbf{PCM:}\\
264076, 264078, 264082, 264085, 264086, 264109, 264110, 264129, 264137, 264138, 264139, 264164, 264188, 264190, 264194, 264197, 264198, 264232, 264233, 264235, 264238, 264259, 264260, 264261, 264262, 264264, 264265, 264266, 264267, 264273, 264277, 264279, 264281, 264305, 264306, 264312, 264336, 264341, 264345, 264346, 264347\\

 \textbf{PHOS:}\\
264076, 264078, 264082, 264085, 264086, 264109, 264110, 264129, 264137, 264138, 264139, 264164, 264168, 264188, 264190, 264197, 264198, 264232, 264233, 264235, 264238, 264259, 264260, 264261, 264262, 264264, 264265, 264266, 264267, 264277, 264279, 264281, 264306, 264312, 264336, 264341, 264345, 264346\\

 \subsection{LHC17c anchored to: LHC18d3, LHC18d3\_extra}

 \textbf{EDC:}\\
270661, 270663, 270665\\

 \textbf{PCMEDC:}\\
270661, 270663, 270665\\

 \textbf{PCM:}\\
270581, 270661, 270663, 270665, 270667\\

 \textbf{PHOS:}\\
270667, 270665, 270663, 270661, 270581\\

 \subsection{LHC17e anchored to:  }

 \textbf{PCM:}\\
270822, 270824, 270827, 270828, 270830\\

 \textbf{PHOS:}\\
270830, 270828, 270827, 270822\\

 \subsection{LHC17f anchored to:  }

 \textbf{PCM:}\\
270854, 270855, 270856, 270861, 270865\\

 \textbf{PHOS:}\\
270854, 270855, 270856, 270861, 270865\\

 \subsection{LHC17h anchored to: LHC18c12, LHC18c12\_extra}

 \textbf{EDCtrigger:}\\
272076, 272100, 272101, 272123, 272151, 272152, 272153, 272154, 272155, 272156, 272194, 272335, 272340, 272359, 272360, 272388, 272389, 272394, 272395, 272399, 272400, 272411, 272413, 272414, 272417, 272461, 272462, 272463, 272466, 272468, 272469, 272521, 272574, 272575, 272585, 272607, 272608, 272610, 272620, 272690, 272691, 272692, 272712, 272746, 272747, 272749, 272760, 272762, 272763, 272764, 272782, 272783, 272828, 272903, 272905, 272932, 272933, 272934, 272947, 272949, 273009, 273077, 273099, 273100\\

 \textbf{EDC:}\\
271870, 271871, 271873, 271874, 271880, 271886, 272018, 272020, 272076, 272100, 272101, 272123, 272151, 272152, 272153, 272154, 272155, 272156, 272194, 272335, 272340, 272359, 272360, 272388, 272389, 272394, 272395, 272399, 272400, 272411, 272413, 272461, 272462, 272463, 272466, 272468, 272521, 272574, 272575, 272585, 272607, 272608, 272610, 272620, 272690, 272691, 272712, 272747, 272749, 272760, 272763, 272764, 272782, 272783, 272784, 272828, 272903, 272905, 272932, 272933, 272934, 272935, 272947, 272949, 273009, 273077, 273099, 273100, 272949, 273009, 273077, 273099, 273100\\

 \textbf{PCMEDC:}\\
271868, 271870, 271871, 271873, 271874, 271880, 271881, 271886, 272018, 272020, 272076, 272100, 272101, 272123, 272151, 272152, 272153, 272154, 272155, 272156, 272194, 272335, 272340, 272359, 272360, 272388, 272389, 272394, 272395, 272399, 272400, 272411, 272413, 272414, 272417, 272461, 272462, 272463, 272466, 272468, 272469, 272521, 272574, 272575, 272585, 272607, 272608, 272610, 272620, 272690, 272691, 272692, 272712, 272746, 272747, 272749, 272760, 272762, 272763, 272764, 272782, 272783, 272784, 272828, 272903, 272905, 272932, 272933, 272934, 272947, 272949, 273009, 273077, 273099, 273100\\

 \textbf{PCM:}\\
273103, 273100, 273099, 273077, 273010, 273009, 272985, 272983, 272976, 272949, 272947, 272939, 272935, 272934, 272933, 272932, 272905, 272903, 272880, 272873, 272871, 272870, 272836, 272834, 272833, 272829, 272828, 272784, 272783, 272782, 272764, 272763, 272760, 272749, 272747, 272712, 272691, 272690, 272620, 272610, 272608, 272607, 272585, 272577, 272575, 272574, 272521, 272468, 272466, 272463, 272462, 272461, 272413, 272411, 272400, 272399, 272395, 272394, 272389, 272388, 272360, 272359, 272340, 272335, 272194, 272156, 272155, 272154, 272153, 272152, 272151, 272123, 272101, 272100, 272076, 272042, 272040, 272039, 272038, 272036, 272020, 272018, 271886, 271880, 271874, 271873, 271871, 271870\\

 \textbf{PHOS:}\\
272151, 272152, 272153, 272154, 272155, 272156, 272194, 272335, 272340, 272359, 272360, 272388, 272389, 272394, 272395, 272399, 272400, 272411, 272413, 272414, 272417, 272461, 272462, 272463, 272466, 272468, 272469, 272521, 272574, 272575, 272577, 272585, 272607, 272608, 272610, 272620, 272690, 272691, 272692, 272712, 272746, 272747, 272749, 272760, 272762, 272763, 272764, 272782, 272783, 272784, 272828, 272829, 272833, 272834, 272836, 272870, 272871, 272880, 272903, 272905, 272932, 272933, 272934, 272935, 272947, 272949, 272976, 272983, 272985, 273009, 273010, 273077, 273099, 273100, 273101, 273103\\

 \subsection{LHC17i anchored to: LHC17k4, LHC17k4\_extra}

 \textbf{EDCtrigger:}\\
273824, 273825, 273885, 273886, 273887, 273889, 273918, 273942, 273943, 273946, 273985, 273986, 274063, 274064, 274092, 274125, 274147, 274148, 274174, 274212, 274232, 274258, 274259, 274263, 274264, 274269, 274270, 274271, 274278, 274280, 274281, 274283, 274329, 274351, 274352, 274355, 274363, 274364, 274442\\

 \textbf{EDC:}\\
273591, 273592, 273593, 273653, 273654, 273687, 273689, 273690, 273695, 273709, 273711, 273719, 273824, 273825, 273885, 273886, 273887, 273889, 273918, 273942, 273943, 273946, 273985, 273986, 274063, 274064, 274092, 274125, 274147, 274148, 274174, 274212, 274232, 274258, 274259, 274263, 274264, 274269, 274270, 274271, 274278, 274280, 274281, 274283, 274329, 274351, 274352, 274355, 274363, 274364, 274385, 274386, 274387, 274388, 274389, 274390, 274442\\

 \textbf{PCMEDC:}\\
273591, 273592, 273593, 273653, 273654, 273687, 273689, 273690, 273695, 273709, 273711, 273719, 273824, 273825, 273885, 273886, 273887, 273889, 273918, 273942, 273943, 273946, 273985, 273986, 274064, 274092, 274125, 274147, 274148, 274174, 274212, 274232, 274258, 274259, 274263, 274264, 274269, 274270, 274271, 274278, 274280, 274281, 274283, 274329, 274351, 274352, 274363, 274364, 274385, 274386, 274387, 274388, 274389, 274390, 274442\\

 \textbf{PCM:}\\
273591, 273592, 273593, 273653, 273654, 273687, 273689, 273690, 273695, 273709, 273711, 273719, 273824, 273825, 273885, 273886, 273887, 273889, 273918, 273942, 273943, 273946, 273985, 273986, 274058, 274064, 274092, 274125, 274147, 274148, 274174, 274212, 274232, 274258, 274259, 274263, 274264, 274266, 274268, 274269, 274270, 274271, 274276, 274278, 274280, 274281, 274283, 274329, 274351, 274352, 274360, 274363, 274364, 274385, 274386, 274387, 274388, 274389, 274390, 274442\\

 \textbf{PHOS:}\\
273824, 273825, 273885, 273886, 273887, 273889, 273918, 273942, 273943, 273946, 273985, 273986, 274058, 274063, 274064, 274092, 274094, 274125, 274147, 274148, 274174, 274212, 274232, 274258, 274259, 274263, 274264, 274266, 274268, 274269, 274270, 274271, 274276, 274278, 274279, 274280, 274281, 274283, 274329, 274351, 274352, 274355, 274357, 274360, 274363, 274364, 274442\\

 \subsection{LHC17j anchored to: LHC17h11, LHC17h11\_extra}

 \textbf{EDCtrigger:}\\
274593, 274595, 274596, 274601, 274653, 274657, 274667, 274669, 274671\\

 \textbf{EDC:}\\
274593, 274595, 274596, 274601, 274653, 274657, 274667, 274669, 274671\\

 \textbf{PCMEDC:}\\
274593, 274595, 274596, 274601, 274653, 274657, 274667, 274669, 274671\\

 \textbf{PCM:}\\
274593, 274594, 274595, 274596, 274601, 274653, 274657, 274667, 274669, 274671\\

 \textbf{PHOS:}\\
274593, 274594, 274595, 274596, 274601, 274653, 274657, 274667, 274669, 274671\\

 \subsection{LHC17k anchored to: LHC18c13, LHC18c13\_extra}

 \textbf{EDCtrigger:}\\
274708, 274801, 274802, 274806, 274815, 274821, 274822, 274877, 274878, 274882, 275067, 275068, 275073, 275075, 275076, 275149, 275150, 275151, 275173, 275174, 275177, 275239, 275245, 275246, 275247, 275283, 275314, 275324, 275326, 275328, 275332, 275360, 275361, 275369, 275372, 275401, 275406, 275448, 275452, 275453, 275456, 275457, 275471, 275472, 275515, 275558, 275559, 275612, 275621, 275622, 275623, 275624, 275648, 275661, 275664, 276097, 276098, 276102, 276104, 276135, 276140, 276145, 276166, 276169, 276170, 276178, 276205, 276230, 276257, 276259, 276290, 276292, 276294, 276302, 276348, 276351, 276435, 276437, 276438, 276439, 276462, 276506, 276507, 276508\\

 \textbf{EDC:}\\
274690, 274708, 274801, 274802, 274806, 274815, 274821, 274822, 274877, 274878, 274882, 275067, 275068, 275073, 275075, 275076, 275149, 275150, 275151, 275173, 275174, 275177, 275239, 275245, 275246, 275247, 275283, 275314, 275324, 275326, 275328, 275332, 275360, 275361, 275369, 275372, 275401, 275406, 275443, 275448, 275452, 275453, 275456, 275457, 275471, 275472, 275515, 275558, 275559, 275612, 275621, 275622, 275623, 275624, 275648, 275661, 275664, 275847, 276097, 276098, 276102, 276104, 276135, 276140, 276145, 276166, 276169, 276170, 276178, 276205, 276230, 276257, 276259, 276290, 276292, 276294, 276302, 276348, 276351, 276435, 276437, 276438, 276439, 276462, 276506, 276507, 276508\\

 \textbf{PCMEDC:}\\
274690, 274708, 274801, 274802, 274806, 274815, 274821, 274822, 274877, 274878, 274882, 275067, 275068, 275073, 275075, 275076, 275149, 275150, 275151, 275173, 275174, 275177, 275239, 275245, 275246, 275247, 275283, 275314, 275324, 275326, 275328, 275332, 275360, 275361, 275369, 275372, 275401, 275406, 275443, 275448, 275452, 275453, 275456, 275457, 275471, 275472, 275515, 275558, 275559, 275612, 275621, 275622, 275623, 275624, 275648, 275661, 275664, 275847, 276097, 276098, 276102, 276104, 276135, 276140, 276145, 276166, 276169, 276170, 276178, 276205, 276230, 276257, 276259, 276290, 276292, 276294, 276302, 276348, 276351, 276435, 276437, 276438, 276439, 276462, 276506, 276507, 276508\\

 \textbf{PCM:}\\
274690, 274708, 274801, 274802, 274803, 274806, 274807, 274811, 274815, 274817, 274821, 274822, 274877, 274878, 274882, 274886, 274978, 274979, 275067, 275068, 275073, 275075, 275076, 275149, 275150, 275151, 275173, 275174, 275177, 275180, 275184, 275188, 275239, 275245, 275246, 275247, 275283, 275314, 275322, 275324, 275326, 275328, 275332, 275333, 275360, 275361, 275369, 275372, 275394, 275395, 275401, 275404, 275406, 275443, 275448, 275452, 275453, 275456, 275457, 275459, 275467, 275471, 275472, 275515, 275558, 275559, 275612, 275617, 275621, 275622, 275623, 275624, 275647, 275648, 275650, 275657, 275661, 275664, 275847, 275924, 275925, 276012, 276013, 276017, 276019, 276020, 276040, 276041, 276045, 276097, 276098, 276099, 276102, 276104, 276135, 276140, 276145, 276166, 276169, 276170, 276177, 276178, 276205, 276230, 276257, 276259, 276290, 276291, 276292, 276294, 276297, 276302, 276307, 276312, 276348, 276351, 276435, 276437, 276438, 276439, 276462, 276506, 276507, 276508\\

 \textbf{PHOS:}\\
274736, 274801, 274802, 274803, 274806, 274807, 274811, 274815, 274817, 274821, 274822, 274877, 274878, 274882, 274883, 274884, 274886, 274889, 274978, 274979, 275067, 275068, 275073, 275075, 275076, 275149, 275150, 275151, 275173, 275174, 275177, 275180, 275184, 275188, 275245, 275246, 275247, 275283, 275314, 275322, 275324, 275326, 275328, 275332, 275333, 275360, 275361, 275369, 275372, 275394, 275395, 275401, 275404, 275406, 275443, 275448, 275452, 275453, 275456, 275457, 275459, 275467, 275471, 275472, 275515, 275558, 275559, 275612, 275617, 275621, 275622, 275623, 275624, 275647, 275648, 275650, 275657, 275661, 275664, 275924, 275925, 276040, 276041, 276045, 276097, 276098, 276099, 276102, 276104, 276105, 276108, 276135, 276140, 276141, 276145, 276166, 276169, 276170, 276177, 276178, 276205, 276230, 276257, 276259, 276290, 276291, 276292, 276294, 276297, 276302, 276307, 276348, 276351, 276435, 276438, 276439, 276462, 276506, 276507, 276508\\

 \subsection{LHC17l anchored to: LHC18a8, LHC18a8\_extra}

 \textbf{EDCtrigger:}\\
276552, 276553, 276556, 276557, 276608, 276669, 276670, 276671, 276672, 276674, 276675, 276762, 276916, 276917, 276920, 276970, 276971, 276972, 277015, 277016, 277017, 277037, 277073, 277076, 277079, 277087, 277091, 277121, 277155, 277180, 277181, 277182, 277183, 277188, 277189, 277193, 277194, 277196, 277197, 277256, 277257, 277262, 277293, 277310, 277312, 277314, 277360, 277383, 277384, 277385, 277389, 277472, 277473, 277476, 277477, 277478, 277479, 277530, 277536, 277537, 277574, 277575, 277576, 277577, 277721, 277722, 277723, 277725, 277745, 277746, 277747, 277749, 277794, 277799, 277800, 277801, 277802, 277834, 277842, 277845, 277847, 277848, 277876, 277897, 277898, 277899, 277900, 277901, 277903, 277907, 277930, 277952, 277987, 277988, 277989, 277996, 278121, 278122, 278126, 278127, 278158, 278163, 278164, 278165, 278166, 278167, 278189, 278191, 278215\\

 \textbf{EDC:}\\
276551, 276552, 276553, 276556, 276557, 276608, 276644, 276670, 276671, 276672, 276674, 276675, 276762, 276916, 276917, 276920, 276967, 276969, 276970, 276971, 276972, 277015, 277016, 277017, 277037, 277073, 277076, 277079, 277087, 277091, 277121, 277155, 277180, 277181, 277182, 277183, 277188, 277189, 277193, 277194, 277196, 277197, 277256, 277257, 277262, 277293, 277310, 277312, 277314, 277360, 277383, 277384, 277385, 277389, 277416, 277417, 277418, 277472, 277473, 277476, 277477, 277478, 277479, 277530, 277531, 277534, 277536, 277537, 277574, 277575, 277576, 277577, 277721, 277722, 277723, 277725, 277745, 277746, 277747, 277749, 277794, 277799, 277800, 277801, 277802, 277834, 277841, 277842, 277845, 277847, 277848, 277876, 277897, 277898, 277899, 277900, 277901, 277903, 277907, 277930, 277952, 277987, 277989, 277996, 278121, 278122, 278126, 278127, 278158, 278163, 278164, 278165, 278166, 278167, 278189, 278191, 278215, 278216\\

 \textbf{PCMEDC:}\\
276552, 276553, 276556, 276557, 276608, 276644, 276669, 276670, 276671, 276672, 276674, 276675, 276762, 276916, 276917, 276920, 276967, 276969, 276970, 276971, 276972, 277015, 277016, 277017, 277037, 277073, 277076, 277079, 277087, 277091, 277121, 277155, 277180, 277181, 277182, 277183, 277188, 277189, 277193, 277194, 277196, 277197, 277256, 277257, 277262, 277293, 277310, 277312, 277314, 277360, 277383, 277384, 277385, 277389, 277416, 277417, 277418, 277470, 277472, 277473, 277476, 277477, 277478, 277479, 277530, 277531, 277534, 277536, 277537, 277574, 277575, 277576, 277577, 277721, 277722, 277723, 277725, 277745, 277746, 277747, 277749, 277794, 277799, 277800, 277801, 277802, 277834, 277841, 277842, 277845, 277847, 277848, 277876, 277897, 277898, 277899, 277900, 277901, 277903, 277907, 277930, 277952, 277987, 277988, 277989, 277996, 278121, 278122, 278126, 278127, 278158, 278163, 278164, 278165, 278166, 278167, 278189, 278191, 278215, 278216\\

 \textbf{PCM:}\\
276551, 276552, 276553, 276556, 276557, 276608, 276644, 276670, 276671, 276672, 276674, 276675, 276762, 276916, 276917, 276920, 276967, 276969, 276970, 276971, 276972, 277015, 277016, 277017, 277037, 277073, 277076, 277079, 277082, 277087, 277091, 277117, 277121, 277155, 277180, 277181, 277182, 277183, 277184, 277188, 277189, 277193, 277194, 277196, 277197, 277256, 277257, 277262, 277293, 277310, 277312, 277314, 277360, 277383, 277384, 277385, 277386, 277389, 277416, 277417, 277418, 277472, 277473, 277476, 277477, 277478, 277479, 277530, 277531, 277534, 277536, 277537, 277574, 277575, 277576, 277577, 277721, 277722, 277723, 277725, 277745, 277746, 277747, 277749, 277794, 277795, 277799, 277800, 277801, 277802, 277805, 277834, 277836, 277841, 277842, 277845, 277847, 277848, 277870, 277876, 277897, 277898, 277899, 277900, 277901, 277903, 277904, 277907, 277930, 277952, 277987, 277989, 277991, 277996, 278121, 278122, 278123, 278126, 278127, 278158, 278163, 278164, 278165, 278166, 278167, 278189, 278191, 278215, 278216\\

 \textbf{PHOS:}\\
276551, 276552, 276553, 276556, 276557, 276608, 276644, 276669, 276670, 276671, 276672, 276674, 276675, 276762, 276916, 276917, 276920, 276967, 276969, 276970, 276971, 276972, 277015, 277016, 277017, 277037, 277073, 277076, 277079, 277082, 277087, 277091, 277117, 277121, 277155, 277180, 277181, 277182, 277183, 277184, 277188, 277189, 277193, 277194, 277196, 277197, 277256, 277257, 277262, 277293, 277310, 277312, 277314, 277360, 277383, 277384, 277385, 277386, 277389, 277416, 277417, 277418, 277470, 277472, 277473, 277476, 277477, 277478, 277479, 277530, 277531, 277536, 277537, 277574, 277575, 277576, 277577, 277721, 277722, 277723, 277725, 277745, 277746, 277747, 277749, 277794, 277795, 277799, 277800, 277801, 277802, 277805, 277834, 277836, 277841, 277842, 277845, 277847, 277848, 277870, 277876, 277897, 277898, 277899, 277900, 277901, 277903, 277904, 277907, 277930, 277952, 277987, 277988, 277991, 278121, 278122, 278123, 278126, 278127, 278130, 278158, 278163, 278164, 278165, 278166, 278167, 278189, 278191, 278215, 278216\\

 \subsection{LHC17m anchored to: LHC17l5, LHC17l5\_extra}

 \textbf{EDCtrigger:}\\
 279005, 279008, 279035, 279036, 279041, 279043, 279044, 279068, 279069, 279073, 279074, 279075, 279106, 279107, 279117, 279118, 279122, 279123, 279130, 279155, 279157, 279199, 279201, 279207, 279208, 279232, 279234, 279235, 279242, 279264, 279265, 279268, 279270, 279273, 279274, 279309, 279310, 279312, 279342, 279344, 279348, 279355, 279391, 279410, 279439, 279441, 279483, 279487, 279491, 279550, 279559, 279632, 279641, 279642, 279676, 279677, 279682, 279687, 279688, 279689, 279718, 279719, 279747, 279749, 279773, 279826, 279853, 279855, 280051, 280052, 280066, 280107, 280108, 280111, 280114, 280126, 280131, 280134, 280140\\

 \textbf{EDC:}\\
278915, 278936, 278939, 278941, 278959, 278960, 278963, 278964, 278999, 279000, 279005, 279007, 279008, 279035, 279036, 279041, 279043, 279044, 279068, 279069, 279073, 279074, 279075, 279106, 279107, 279117, 279118, 279122, 279123, 279130, 279155, 279157, 279199, 279201, 279207, 279208, 279232, 279234, 279235, 279242, 279264, 279265, 279268, 279270, 279273, 279274, 279309, 279310, 279312, 279342, 279344, 279348, 279354, 279355, 279391, 279410, 279439, 279441, 279483, 279487, 279488, 279491, 279550, 279559, 279630, 279632, 279641, 279642, 279676, 279677, 279682, 279683, 279687, 279688, 279689, 279718, 279719, 279747, 279749, 279773, 279826, 279830, 279853, 279855, 280051, 280052, 280066, 280107, 280108, 280111, 280114, 280126, 280131, 280134, 280140\\

 \textbf{PCMEDC:}\\
278915, 278936, 278939, 278941, 278959, 278960, 278963, 278964, 278999, 279000, 279005, 279007, 279008, 279035, 279036, 279041, 279043, 279044, 279068, 279069, 279073, 279074, 279075, 279106, 279107, 279117, 279118, 279122, 279123, 279130, 279155, 279157, 279199, 279201, 279207, 279208, 279232, 279234, 279235, 279242, 279264, 279265, 279268, 279270, 279273, 279274, 279309, 279310, 279312, 279342, 279344, 279348, 279354, 279355, 279391, 279410, 279439, 279441, 279483, 279487, 279488, 279491, 279550, 279559, 279630, 279632, 279641, 279642, 279676, 279677, 279682, 279683, 279687, 279688, 279689, 279718, 279719, 279747, 279749, 279773, 279826, 279830, 279853, 279855, 280051, 280052, 280066, 280107, 280108, 280111, 280114, 280126, 280131, 280134, 280140\\

 \textbf{PCM:}\\
280140, 280135, 280134, 280131, 280126, 280118, 280114, 280111, 280108, 280107, 280066, 280052, 280051, 279879, 279855, 279854, 279853, 279830, 279827, 279826, 279773, 279749, 279747, 279719, 279718, 279715, 279689, 279688, 279687, 279684, 279683, 279682, 279679, 279677, 279676, 279642, 279641, 279632, 279630, 279559, 279550, 279491, 279488, 279487, 279483, 279441, 279439, 279435, 279410, 279391, 279355, 279354, 279349, 279348, 279344, 279342, 279312, 279310, 279309, 279274, 279273, 279270, 279268, 279267, 279265, 279264, 279242, 279238, 279235, 279234, 279232, 279208, 279207, 279201, 279199, 279157, 279155, 279130, 279123, 279122, 279118, 279117, 279107, 279106, 279075, 279074, 279073, 279069, 279068, 279044, 279043, 279041, 279036, 279035, 279008, 279007, 279005, 279000, 278999, 278964, 278963, 278960, 278959, 278941, 278939, 278936, 278915, 278914\\

 \textbf{PHOS:}\\
279234, 279235, 279238, 279242, 279264, 279265, 279267, 279268, 279270, 279273, 279274, 279309, 279310, 279312, 279342, 279348, 279349, 279354, 279355, 279391, 279435, 279439, 279441, 279483, 279487, 279488, 279491, 279550, 279559, 279630, 279632, 279641, 279642, 279676, 279679, 279682, 279683, 279684, 279687, 279688, 279689, 279715, 279718, 279719, 279747, 279749, 279773, 279826, 279827, 279830, 279853, 279855, 279879, 279880, 280051, 280052, 280066, 280107, 280108, 280111, 280114, 280118, 280126, 280131, 280134, 280135, 280140\\

 \subsection{LHC17o anchored to: LHC18a9, LHC18a9\_extra}

 \textbf{EDCtrigger:}\\
280282, 280283, 280284, 280285, 280286, 280348, 280349, 280350, 280351, 280352, 280374, 280375, 280403, 280412, 280413, 280415, 280419, 280518, 280519, 280546, 280547, 280550, 280551, 280574, 280576, 280581, 280583, 280613, 280634, 280636, 280637, 280647, 280648, 280650, 280671, 280673, 280676, 280679, 280705, 280706, 280729, 280753, 280754, 280755, 280756, 280757, 280761, 280762, 280763, 280764, 280765, 280766, 280767, 280768, 280786, 280787, 280793, 280842, 280844, 280845, 280847, 280848, 280849, 280854, 280856, 280880, 280881, 280897, 280936, 280943, 280947, 280994, 280996, 280997, 280998, 280999, 281032, 281033, 281035, 281036, 281060, 281061, 281062, 281080, 281081, 281179, 281180, 281181, 281189, 281190, 281191, 281212, 281213, 281240, 281241, 281242, 281243, 281244, 281271, 281273, 281275, 281277, 281301, 281321, 281415, 281443, 281444, 281446, 281449, 281450, 281563, 281568, 281569, 281574, 281580, 281705, 281706, 281707, 281713, 281750, 281751, 281753, 281892, 281893, 281894, 281895, 281915, 281916, 281918, 281920, 281928, 281931, 281932, 281939, 281961\\

 \textbf{EDC:}\\
280282, 280283, 280284, 280285, 280286, 280348, 280349, 280350, 280351, 280352, 280374, 280375, 280403, 280412, 280413, 280415, 280419, 280499, 280518, 280519, 280546, 280547, 280550, 280551, 280574, 280575, 280576, 280581, 280583, 280613, 280634, 280636, 280637, 280647, 280648, 280650, 280671, 280673, 280676, 280679, 280705, 280706, 280729, 280753, 280754, 280755, 280756, 280757, 280761, 280762, 280763, 280764, 280765, 280766, 280767, 280768, 280786, 280787, 280792, 280793, 280842, 280844, 280845, 280847, 280848, 280849, 280854, 280856, 280880, 280881, 280897, 280936, 280943, 280947, 280994, 280996, 280997, 280998, 280999, 281032, 281033, 281035, 281036, 281060, 281061, 281062, 281080, 281081, 281179, 281180, 281181, 281189, 281190, 281191, 281212, 281213, 281240, 281241, 281242, 281243, 281244, 281271, 281273, 281275, 281277, 281301, 281321, 281415, 281441, 281443, 281444, 281446, 281449, 281450, 281475, 281557, 281563, 281568, 281569, 281574, 281580, 281581, 281583, 281592, 281633, 281705, 281706, 281707, 281713, 281750, 281751, 281753, 281754, 281892, 281893, 281894, 281895, 281915, 281916, 281918, 281920, 281928, 281931, 281932, 281939, 281940, 281961\\

 \textbf{PCMEDC:}\\
280282, 280283, 280284, 280285, 280286, 280348, 280349, 280350, 280351, 280374, 280375, 280403, 280412, 280413, 280415, 280419, 280499, 280518, 280519, 280546, 280547, 280550, 280551, 280574, 280575, 280576, 280581, 280583, 280613, 280634, 280636, 280637, 280647, 280671, 280679, 280705, 280706, 280729, 280753, 280754, 280755, 280756, 280757, 280761, 280762, 280763, 280764, 280765, 280766, 280767, 280768, 280786, 280787, 280792, 280793, 280842, 280844, 280845, 280847, 280848, 280849, 280854, 280856, 280880, 280881, 280897, 280936, 280943, 280947, 280994, 280996, 280997, 280998, 280999, 281032, 281033, 281035, 281036, 281060, 281061, 281062, 281080, 281081, 281179, 281180, 281181, 281189, 281190, 281191, 281212, 281213, 281240, 281241, 281242, 281243, 281244, 281271, 281273, 281275, 281277, 281301, 281321, 281415, 281441, 281443, 281444, 281446, 281449, 281450, 281475, 281557, 281563, 281568, 281569, 281574, 281583, 281592, 281633, 281892, 281893, 281894, 281895, 281915, 281916, 281918, 281920, 281928, 281931, 281932, 281939, 281940, 281961\\

 \textbf{PCM:}\\
280282, 280283, 280284, 280285, 280286, 280290, 280310, 280312, 280348, 280349, 280350, 280351, 280374, 280375, 280403, 280405, 280406, 280412, 280413, 280415, 280419, 280443, 280445, 280446, 280447, 280448, 280490, 280499, 280518, 280519, 280546, 280547, 280550, 280551, 280574, 280575, 280576, 280581, 280583, 280613, 280634, 280636, 280637, 280639, 280645, 280647, 280671, 280679, 280681, 280705, 280706, 280729, 280753, 280754, 280755, 280756, 280757, 280761, 280762, 280763, 280764, 280765, 280766, 280767, 280768, 280786, 280787, 280792, 280793, 280842, 280844, 280845, 280847, 280848, 280849, 280854, 280856, 280880, 280881, 280890, 280897, 280936, 280940, 280943, 280947, 280990, 280994, 280996, 280997, 280998, 280999, 281032, 281033, 281035, 281036, 281060, 281061, 281062, 281080, 281081, 281179, 281180, 281181, 281189, 281190, 281191, 281212, 281213, 281240, 281241, 281242, 281243, 281244, 281271, 281273, 281275, 281277, 281301, 281321, 281415, 281441, 281443, 281444, 281446, 281449, 281450, 281475, 281477, 281509, 281511, 281557, 281562, 281563, 281568, 281569, 281574, 281583, 281592, 281633, 281892, 281893, 281894, 281895, 281915, 281916, 281918, 281920, 281928, 281931, 281932, 281939, 281940, 281953, 281956, 281961\\

 \textbf{PHOS:}\\
280282, 280283, 280284, 280285, 280286, 280290, 280310, 280312, 280348, 280349, 280350, 280351, 280352, 280374, 280375, 280403, 280405, 280406, 280412, 280413, 280415, 280419, 280443, 280445, 280446, 280447, 280448, 280490, 280499, 280518, 280519, 280546, 280547, 280551, 280574, 280575, 280576, 280581, 280583, 280613, 280634, 280636, 280637, 280639, 280645, 280647, 280648, 280650, 280671, 280673, 280676, 280679, 280681, 280705, 280729, 280753, 280755, 280756, 280757, 280761, 280762, 280763, 280764, 280765, 280766, 280767, 280768, 280786, 280787, 280792, 280793, 280842, 280844, 280845, 280847, 280848, 280849, 280854, 280856, 280880, 280881, 280890, 280897, 280936, 280940, 280943, 280947, 280990, 280994, 280996, 280997, 280998, 280999, 281032, 281033, 281035, 281036, 281060, 281061, 281062, 281080, 281081, 281179, 281180, 281181, 281189, 281190, 281191, 281212, 281213, 281240, 281241, 281242, 281243, 281244, 281271, 281273, 281275, 281277, 281301, 281321, 281415, 281441, 281443, 281444, 281446, 281449, 281450, 281475, 281477, 281509, 281557, 281562, 281563, 281592, 281633, 281705, 281706, 281707, 281709, 281713, 281741, 281750, 281751, 281753, 281754, 281755, 281756, 281892, 281893, 281894, 281895, 281915, 281916, 281918, 281931, 281932, 281939, 281940, 281953, 281956, 281961\\

 \subsection{LHC17r anchored to: LHC18a1, LHC18a1\_extra}

 \textbf{EDCtrigger:}\\
282544, 282545, 282546, 282579, 282580, 282606, 282607, 282608, 282609, 282620, 282622, 282651, 282666, 282667, 282670, 282671, 282673, 282676, 282677, 282700, 282702, 282703, 282704\\

 \textbf{EDC:}\\
282544, 282545, 282546, 282573, 282579, 282580, 282606, 282607, 282608, 282609, 282618, 282620, 282622, 282629, 282651, 282666, 282667, 282670, 282671, 282673, 282676, 282677, 282700, 282702, 282703, 282704\\

 \textbf{PCMEDC:}\\
 282544, 282545, 282546, 282573, 282579, 282580, 282606, 282607, 282608, 282609, 282618, 282620, 282622, 282629, 282651, 282666, 282667, 282670, 282671, 282673, 282676, 282677, 282700, 282702, 282703, 282704\\

 \textbf{PCM:}\\
282528, 282544, 282545, 282546, 282573, 282575, 282579, 282580, 282606, 282607, 282608, 282609, 282618, 282620, 282622, 282629, 282651, 282653, 282666, 282667, 282668, 282670, 282671, 282673, 282676, 282677, 282700, 282702, 282703, 282704\\

 \textbf{PHOS:}\\
282704, 282703, 282702, 282700, 282677, 282676, 282671, 282670, 282668, 282667, 282666, 282653, 282651, 282629, 282620, 282609, 282608, 282607, 282606, 282580, 282579, 282575, 282573, 282546, 282545, 282544, 282528\\


 \subsection{LHC18b anchored to:  }

  \textbf{EDCtrigger:}\\
 285009, 285011, 285012, 285013, 285014, 285015, 285106, 285108, 285125, 285127, 285165, 285200, 285202, 285203, 285222, 285224, 285327, 285328, 285347, 285364, 285365, 285396 \\

 \textbf{EDC:}\\
 285009, 285011, 285012, 285013, 285014, 285015, 285106, 285108, 285125, 285127, 285165, 285200, 285202, 285203, 285222, 285224, 285327, 285328, 285347, 285364, 285365, 285396 \\

 \textbf{PCMEDC:}\\
 285009, 285011, 285012, 285013, 285014, 285015, 285106, 285108, 285125, 285127, 285165, 285200, 285202, 285203, 285222, 285224, 285327, 285328, 285347, 285364, 285365, 285396  \\

 \textbf{PHOS:}\\
285008, 285009, 285010, 285011, 285012, 285013, 285014, 285015, 285064, 285065, 285066, 285106, 285108, 285125, 285127, 285165, 285200, 285202, 285203, 285222, 285224, 285327, 285328, 285347, 285364, 285365, 285396, 285447\\

 \subsection{LHC18d anchored to:  }

 \textbf{EDCtrigger:}\\
 286014, 286025, 286064, 286124, 286127, 286129, 286130, 286159, 286198, 286201, 286202, 286203, 286229, 286230, 286231, 286254, 286255, 286257, 286258, 286261, 286263, 286282, 286284, 286287, 286288, 286289, 286308, 286309, 286310, 286311, 286313, 286314, 286336, 286337, 286340, 286341, 286345, 286348, 286349, 286350\\

 \textbf{EDC:}\\
 286014, 286025, 286064, 286124, 286127, 286129, 286130, 286159, 286198, 286201, 286202, 286203, 286229, 286230, 286231, 286254, 286255, 286257, 286258, 286261, 286263, 286282, 286284, 286287, 286288, 286289, 286308, 286309, 286310, 286311, 286313, 286314, 286336, 286337, 286340, 286341, 286345, 286348, 286349, 286350, 286348, 286349, 286350 \\


 \textbf{PCMEDC:}\\
 286014, 286025, 286064, 286124, 286127, 286129, 286130, 286159, 286198, 286201, 286202, 286203, 286229, 286230, 286231, 286254, 286255, 286257, 286258, 286261, 286263, 286282, 286284, 286287, 286288, 286289, 286308, 286309, 286310, 286311, 286313, 286314, 286336, 286337, 286340, 286341, 286345, 286348, 286349, 286350 \\

 \textbf{PHOS:}\\
285978, 285979, 285980, 286014, 286025, 286027, 286028, 286030, 286064, 286124, 286127, 286129, 286130, 286159, 286198, 286199, 286201, 286202, 286203, 286229, 286230, 286231, 286254, 286255, 286257, 286258, 286261, 286263, 286282, 286284, 286287, 286288, 286289, 286308, 286309, 286310, 286311, 286312, 286313, 286314, 286336, 286337, 286340, 286341, 286345, 286348, 286349, 286350\\

 \subsection{LHC18e anchored to:  }

  \textbf{EDCtrigger:}\\
 286380, 286426, 286428, 286454, 286455, 286482, 286501, 286502, 286508, 286509, 286511, 286566, 286567, 286568, 286569, 286591, 286592, 286653, 286661, 286695, 286731, 286799, 286801, 286846, 286848, 286850, 286852, 286874, 286876, 286907, 286908, 286910, 286911, 286930, 286931, 286932, 286936, 286937 \\

 \textbf{EDC:}\\
 286380, 286426, 286427, 286428, 286454, 286455, 286482, 286501, 286502, 286508, 286509, 286511, 286566, 286567, 286568, 286569, 286591, 286592, 286594, 286653, 286661, 286695, 286731, 286799, 286801, 286846, 286848, 286850, 286852, 286874, 286876, 286877, 286907, 286908, 286910, 286911, 286930, 286931, 286932, 286936, 286937 \\


 \textbf{PCMEDC:}\\
 286380, 286426, 286428, 286454, 286455, 286482, 286501, 286502, 286508, 286509, 286511, 286566, 286567, 286568, 286569, 286591, 286592, 286653, 286661, 286695, 286731, 286799, 286801, 286846, 286848, 286850, 286852, 286874, 286876, 286907, 286908, 286910, 286911, 286930, 286931, 286932, 286936, 286937 \\

 \textbf{PHOS:}\\
286380, 286426, 286427, 286428, 286454, 286455, 286482, 286501, 286502, 286508, 286509, 286511, 286566, 286567, 286568, 286569, 286591, 286592, 286594, 286633, 286653, 286661, 286695, 286731, 286799, 286801, 286805, 286809, 286810, 286846, 286848, 286850, 286852, 286874, 286876, 286877, 286907, 286908, 286910, 286911, 286930, 286931, 286932, 286933, 286936, 286937\\

 \subsection{LHC18f anchored to:  }

	\textbf{EDCtrigger:}\\
	287000, 287021, 287063, 287064, 287071, 287077, 287137, 287155, 287185, 287201, 287209, 287248, 287249, 287250, 287254, 287283, 287323, 287324, 287325, 287343, 287346, 287347, 287349, 287353, 287355, 287356, 287360, 287380, 287381, 287385, 287387, 287388, 287389, 287413, 287480, 287481, 287484, 287486, 287513, 287516, 287517, 287520, 287521, 287573, 287575, 287576, 287578, 287657, 287783, 287876, 287877, 287883, 287884, 287885, 287911, 287912, 287913, 287915, 287923, 287941, 287975, 287977 \\

	\textbf{EDC:}\\
	287000, 287021, 287063, 287064, 287071, 287072, 287077, 287137, 287155, 287185, 287201, 287208, 287209, 287248, 287249, 287250, 287254, 287283, 287323, 287324, 287325, 287343, 287346, 287347, 287349, 287353, 287355, 287356, 287360, 287380, 287381, 287385, 287387, 287388, 287389, 287413, 287480, 287481, 287484, 287486, 287513, 287516, 287517, 287518, 287520, 287521, 287524, 287573, 287575, 287576, 287578, 287657, 287658, 287783, 287784, 287876, 287877, 287883, 287884, 287885, 287911, 287912, 287913, 287915, 287923, 287941, 287975, 287977, 287923, 287941, 287975, 287977 \\

	\textbf{PCMEDC:}\\
	287000, 287021, 287063, 287064, 287071, 287077, 287137, 287155, 287185, 287201, 287209, 287248, 287249, 287250, 287254, 287283, 287323, 287324, 287325, 287343, 287346, 287347, 287349, 287353, 287355, 287356, 287360, 287380, 287381, 287385, 287387, 287388, 287389, 287413, 287480, 287481, 287484, 287486, 287513, 287516, 287517, 287520, 287521, 287573, 287575, 287576, 287578, 287657, 287783, 287876, 287877, 287883, 287884, 287885, 287911, 287912, 287913, 287915, 287923, 287941, 287975, 287977  \\

 \textbf{PHOS:}\\
287000, 287021, 287063, 287064, 287066, 287071, 287072, 287077, 287137, 287155, 287185, 287201, 287202, 287203, 287204, 287208, 287209, 287248, 287249, 287250, 287251, 287254, 287283, 287323, 287324, 287325, 287343, 287344, 287346, 287347, 287349, 287353, 287355, 287356, 287360, 287380, 287381, 287385, 287387, 287388, 287389, 287413, 287451, 287480, 287481, 287484, 287486, 287513, 287516, 287517, 287518, 287520, 287521, 287524, 287573, 287575, 287576, 287578, 287616, 287654, 287656, 287657, 287658, 287783, 287784, 287876, 287877, 287883, 287884, 287885, 287911, 287912, 287913, 287915, 287923, 287941, 287975, 287977\\

 \subsection{LHC18g anchored to:  }

 \textbf{EDCtrigger:}\\
 288619, 288650 \\

 \textbf{EDC:}\\
 288619, 288650 \\

 \textbf{PCMEDC:}\\
 288619, 288650 \\

 \textbf{PHOS:}\\
288619, 288640, 288642, 288644, 288650, 288687, 288689, 288690, 288743, 288748, 288750\\

 \subsection{LHC18h anchored to:  }

 \textbf{EDCtrigger:}\\
 288804, 288806\\

 \textbf{EDC:}\\
 288804, 288806\\

 \textbf{PCMEDC:}\\
288804, 288806\\

 \textbf{PHOS:}\\
288804, 288806\\

 \subsection{LHC18i anchored to:  }

 \textbf{EDCtrigger:}\\
  288861, 288862, 288864, 288868, 288897, 288903, 288908, 288909\\

 \textbf{EDC:}\\
 288861, 288862, 288864, 288868, 288897, 288903, 288908, 288909\\

 \textbf{PCMEDC:}\\
288861, 288862, 288864, 288868, 288897, 288903, 288908, 288909 \\

 \textbf{PHOS:}\\
288861, 288862, 288863, 288864, 288868, 288897, 288902, 288903, 288908, 288909\\

 \subsection{LHC18j anchored to:  }

 \textbf{EDCtrigger:}\\
 288943\\

 \textbf{EDC:}\\
 288943\\

 \textbf{PCMEDC:}\\
288943\\

 \textbf{PHOS:}\\
288943\\

 \subsection{LHC18k anchored to:  }

 \textbf{EDCtrigger:}\\
 289165, 289166, 289169, 289172, 289176, 289177, 289198, 289199, 289200, 289201\\

 \textbf{EDC:}\\
 289165, 289166, 289169, 289172, 289176, 289177, 289198, 289199, 289200, 289201\\

 \textbf{PCMEDC:}\\
289165, 289166, 289169, 289172, 289176, 289177, 289198, 289199, 289200, 289201  \\

 \textbf{PHOS:}\\
289165, 289166, 289167, 289169, 289172, 289175, 289176, 289177, 289198, 289199, 289200, 289201\\

 \subsection{LHC18l anchored to:  }

 \textbf{EDCtrigger:}\\
 289240, 289242, 289243, 289253, 289275, 289276, 289277, 289278, 289280, 289281, 289300, 289303, 289306, 289308, 289309, 289354, 289355, 289356, 289365, 289366, 289367, 289368, 289369, 289370, 289373, 289374, 289426, 289444, 289462, 289463, 289465, 289466, 289468, 289493, 289494, 289521, 289547, 289721, 289729, 289731, 289732, 289757, 289775, 289808, 289811, 289815, 289816, 289817, 289818, 289830, 289849, 289854, 289855, 289856, 289857, 289879, 289880, 289884, 289928, 289935, 289940, 289943, 289965, 289966, 289971 \\

 \textbf{EDC:}\\
 289240, 289242, 289243, 289253, 289275, 289276, 289277, 289278, 289280, 289281, 289300, 289303, 289306, 289308, 289309, 289354, 289355, 289356, 289365, 289366, 289367, 289368, 289369, 289370, 289373, 289374, 289426, 289444, 289462, 289463, 289465, 289466, 289468, 289493, 289494, 289521, 289547, 289721, 289729, 289731, 289732, 289757, 289775, 289808, 289811, 289815, 289816, 289817, 289818, 289830, 289849, 289854, 289855, 289856, 289857, 289879, 289880, 289884, 289928, 289935, 289940, 289943, 289965, 289966, 289971, 289241, 289254, 289353, 289724, 289814, 289966\\

 \textbf{PCMEDC:}\\
 289240, 289242, 289243, 289253, 289275, 289276, 289277, 289278, 289280, 289281, 289300, 289303, 289306, 289308, 289309, 289354, 289355, 289356, 289365, 289366, 289367, 289368, 289369, 289370, 289373, 289374, 289426, 289444, 289462, 289463, 289465, 289466, 289468, 289493, 289494, 289521, 289547, 289721, 289729, 289731, 289732, 289757, 289775, 289808, 289811, 289815, 289816, 289817, 289818, 289830, 289849, 289854, 289855, 289856, 289857, 289879, 289880, 289884, 289928, 289935, 289940, 289943, 289965, 289966, 289971  \\

 \textbf{PHOS:}\\
289240, 289241, 289242, 289243, 289247, 289249, 289253, 289254, 289275, 289276, 289277, 289278, 289280, 289281, 289300, 289303, 289306, 289308, 289309, 289330, 289331, 289353, 289354, 289355, 289356, 289363, 289365, 289366, 289367, 289368, 289369, 289370, 289373, 289374, 289426, 289444, 289462, 289463, 289465, 289466, 289468, 289493, 289494, 289521, 289547, 289574, 289576, 289577, 289582, 289625, 289632, 289634, 289657, 289658, 289659, 289660, 289664, 289666, 289721, 289723, 289724, 289729, 289731, 289732, 289757, 289775, 289808, 289811, 289814, 289815, 289816, 289817, 289818, 289830, 289849, 289852, 289854, 289855, 289856, 289857, 289879, 289880, 289884, 289928, 289931, 289935, 289940, 289941, 289943, 289965, 289966, 289971\\

 \subsection{LHC18m anchored to:  }
\textbf{EDCtrigger:}\\
290323, 290324, 290327, 290350, 290374, 290375, 290399, 290401, 290411, 290412, 290425, 290426, 290427, 290428, 290456, 290458, 290459, 290469, 290499, 290500, 290538, 290539, 290540, 290544, 290549, 290550, 290553, 290588, 290590, 290612, 290613, 290614, 290615, 290627, 290632, 290645, 290658, 290660, 290665, 290941, 290943, 290944, 290948, 290974, 290975, 290976, 290979, 290980, 291002, 291003, 291004, 291005, 291035, 291037, 291041, 291065, 291066, 291069, 291093, 291100, 291110, 291111, 291116, 291143, 291188, 291209, 291240, 291257, 291262, 291263, 291265, 291266, 291282, 291283, 291284, 291285, 291286, 291360, 291361, 291362, 291373, 291375, 291377, 291397, 291399, 291402, 291417, 291419, 291420, 291424, 291446, 291447, 291453, 291456, 291457, 291481, 291482, 291484, 291485, 292242, 292397, 292398, 292405, 292406, 292428, 292429, 292430, 292432, 292434, 292456, 292457, 292460, 292461, 292495, 292496, 292497, 292500, 292521, 292523, 292524, 292526, 292553, 292554, 292557, 292559, 292560, 292563, 292584, 292586, 292693, 292695, 292696, 292698, 292701, 292704, 292737, 292739, 292747, 292748, 292750, 292809, 292810, 292811, 292831, 292832, 292836\\

\textbf{EDC:}\\
290323, 290324, 290327, 290350, 290374, 290375, 290376, 290399, 290401, 290411, 290412, 290425, 290426, 290427, 290428, 290456, 290458, 290459, 290469, 290499, 290500, 290538, 290539, 290540, 290544, 290549, 290550, 290553, 290588, 290590, 290612, 290613, 290614, 290615, 290627, 290632, 290645, 290658, 290660, 290665, 290687, 290689, 290692, 290696, 290699, 290721, 290742, 290764, 290766, 290769, 290774, 290787, 290790, 290841, 290843, 290846, 290848, 290860, 290862, 290886, 290887, 290892, 290894, 290895, 290932, 290935, 290941, 290943, 290944, 290948, 290974, 290975, 290976, 290979, 290980, 291002, 291003, 291004, 291005, 291035, 291037, 291041, 291065, 291066, 291069, 291093, 291100, 291101, 291110, 291111, 291116, 291143, 291188, 291209, 291240, 291257, 291262, 291263, 291265, 291266, 291282, 291283, 291284, 291285, 291286, 291360, 291361, 291362, 291373, 291375, 291377, 291397, 291399, 291402, 291416, 291417, 291419, 291420, 291424, 291446, 291447, 291453, 291456, 291457, 291481, 291482, 291484, 291485, 291590, 291614, 291615, 291618, 291622, 291624, 291626, 291657, 291661, 291665, 291690, 291692, 291694, 291697, 291698, 291702, 291706, 291729, 291755, 291756, 291760, 291768, 291769, 291795, 291796, 291803, 291942, 291943, 291944, 291945, 291946, 291948, 291953, 291976, 291977, 291982, 292012, 292040, 292060, 292061, 292062, 292067, 292075, 292080, 292081, 292106, 292107, 292108, 292109, 292114, 292115, 292140, 292160, 292161, 292162, 292163, 292164, 292166, 292167, 292168, 292192, 292218, 292240, 292242, 292265, 292270, 292273, 292274, 292298, 292397, 292398, 292405, 292406, 292428, 292429, 292430, 292432, 292434, 292456, 292457, 292460, 292461, 292495, 292496, 292497, 292500, 292521, 292523, 292524, 292526, 292553, 292554, 292557, 292559, 292560, 292563, 292584, 292586, 292693, 292695, 292696, 292698, 292701, 292704, 292737, 292739, 292747, 292748, 292750, 292803, 292804, 292809, 292810, 292811, 292831, 292832, 292836, 292839, 292832, 292836, 292839 \\

\textbf{PCMEDC:}\\
290323, 290324, 290327, 290350, 290374, 290375, 290376, 290399, 290401, 290411, 290412, 290425, 290426, 290427, 290428, 290456, 290458, 290459, 290469, 290499, 290500, 290538, 290539, 290540, 290544, 290549, 290550, 290553, 290588, 290590, 290612, 290613, 290614, 290615, 290627, 290632, 290645, 290658, 290660, 290665, 290687, 290689, 290692, 290696, 290699, 290721, 290764, 290766, 290769, 290774, 290787, 290790, 290841, 290843, 290846, 290848, 290860, 290862, 290886, 290887, 290892, 290894, 290895, 290932, 290935, 290941, 290943, 290944, 290948, 290974, 290975, 290976, 290979, 290980, 291002, 291003, 291004, 291005, 291035, 291037, 291041, 291065, 291066, 291069, 291093, 291100, 291110, 291111, 291116, 291143, 291188, 291209, 291240, 291257, 291262, 291263, 291265, 291266, 291282, 291283, 291284, 291285, 291286, 291360, 291361, 291362, 291373, 291375, 291377, 291397, 291399, 291402, 291417, 291419, 291420, 291424, 291446, 291447, 291453, 291456, 291457, 291481, 291482, 291484, 291485, 291614, 291615, 291618, 291622, 291624, 291626, 291657, 291661, 291665, 291690, 291692, 291694, 291698, 291706, 291729, 291755, 291756, 291760, 291768, 291795, 291796, 291803, 291942, 291943, 291944, 291945, 291946, 291948, 291953, 291976, 291977, 291982, 292012, 292040, 292060, 292061, 292062, 292067, 292075, 292080, 292081, 292106, 292107, 292108, 292109, 292114, 292115, 292140, 292160, 292161, 292162, 292163, 292164, 292166, 292167, 292168, 292192, 292218, 292240, 292242, 292265, 292270, 292273, 292274, 292298, 292397, 292398, 292405, 292406, 292428, 292429, 292430, 292432, 292434, 292456, 292457, 292460, 292461, 292495, 292496, 292497, 292500, 292521, 292523, 292524, 292526, 292553, 292554, 292557, 292559, 292560, 292563, 292584, 292586, 292693, 292695, 292696, 292698, 292701, 292704, 292737, 292739, 292747, 292748, 292750, 292809, 292810, 292811, 292831, 292832, 292836 \\

 \subsection{LHC18n anchored to:  }

 \textbf{EDCtrigger:}\\
 293357, 293359\\

 \textbf{EDC:}\\
 293357, 293359\\

 \textbf{PCMEDC:}\\
293357, 293359\\

 \textbf{PHOS:}\\
293357, 293359\\

 \subsection{LHC18o anchored to:  }

 \textbf{EDCtrigger:}\\
 293368, 293386, 293392, 293413, 293424, 293474, 293475, 293494, 293496, 293497, 293570, 293571, 293582, 293587, 293588, 293686, 293690, 293691, 293695, 293696, 293740, 293741, 293770, 293773, 293774, 293776, 293799, 293805, 293806, 293809, 293829, 293830, 293831, 293856, 293886, 293891, 293893, 293896, 293898\\

 \textbf{EDC:}\\
 293368, 293386, 293392, 293413, 293424, 293474, 293475, 293494, 293496, 293497, 293570, 293571, 293578, 293582, 293587, 293588, 293686, 293689, 293690, 293691, 293695, 293696, 293740, 293741, 293770, 293773, 293774, 293776, 293799, 293805, 293806, 293809, 293829, 293830, 293831, 293856, 293886, 293891, 293893, 293896, 293898\\

 \textbf{PCMEDC:}\\
 293368, 293386, 293392, 293413, 293424, 293474, 293475, 293494, 293496, 293497, 293570, 293571, 293582, 293587, 293588, 293686, 293690, 293691, 293695, 293696, 293740, 293741, 293770, 293773, 293774, 293776, 293799, 293805, 293806, 293809, 293829, 293830, 293831, 293856, 293886, 293891, 293893, 293896, 293898 \\

 \textbf{PHOS:}\\
293368, 293386, 293392, 293413, 293424, 293474, 293475, 293494, 293496, 293497, 293570, 293571, 293573, 293578, 293579, 293582, 293583, 293587, 293588, 293686, 293689, 293690, 293691, 293692, 293695, 293696, 293698, 293740, 293741, 293770, 293773, 293774, 293776, 293799, 293802, 293805, 293806, 293807, 293809, 293829, 293830, 293831, 293856, 293886, 293891, 293893, 293896, 293898\\

 \subsection{LHC18p anchored to:  }

 \textbf{EDCtrigger:}\\
 294009, 294010, 294011, 294012, 294013, 294131, 294152, 294154, 294155, 294156, 294199, 294200, 294208, 294210, 294212, 294241, 294242, 294305, 294307, 294310, 294524, 294525, 294526, 294529, 294530, 294531, 294553, 294556, 294558, 294562, 294563, 294586, 294587, 294588, 294590, 294591, 294593, 294632, 294633, 294634, 294636, 294653, 294715, 294716, 294721, 294722, 294741, 294742, 294743, 294744, 294745, 294746, 294747, 294769, 294774, 294775, 294805, 294815, 294816, 294817, 294818, 294852, 294875, 294877, 294883, 294884, 294925 \\

 \textbf{EDC:}\\
 294009, 294010, 294011, 294012, 294013, 294131, 294152, 294154, 294155, 294156, 294199, 294200, 294208, 294210, 294212, 294241, 294242, 294305, 294307, 294310, 294524, 294525, 294526, 294529, 294530, 294531, 294553, 294556, 294558, 294562, 294563, 294586, 294587, 294588, 294590, 294591, 294593, 294632, 294633, 294634, 294636, 294653, 294715, 294716, 294721, 294722, 294741, 294742, 294743, 294744, 294745, 294746, 294747, 294769, 294774, 294775, 294805, 294815, 294816, 294817, 294818, 294852, 294875, 294877, 294883, 294884, 294925\\

 \textbf{PCMEDC:}\\
294009, 294010, 294011, 294012, 294013, 294131, 294152, 294154, 294155, 294156, 294199, 294200, 294208, 294210, 294212, 294241, 294242, 294305, 294307, 294310, 294524, 294525, 294526, 294529, 294530, 294531, 294553, 294556, 294558, 294562, 294563, 294586, 294587, 294588, 294590, 294591, 294593, 294632, 294633, 294634, 294636, 294653, 294715, 294716, 294721, 294722, 294741, 294742, 294743, 294744, 294745, 294746, 294747, 294769, 294774, 294775, 294805, 294815, 294816, 294817, 294818, 294852, 294875, 294877, 294883, 294884, 294925\\

 \textbf{PHOS:}\\
294009, 294010, 294011, 294012, 294013, 294128, 294131, 294152, 294154, 294155, 294156, 294199, 294200, 294201, 294205, 294208, 294210, 294212, 294241, 294242, 294305, 294307, 294308, 294310, 294502, 294503, 294524, 294525, 294526, 294527, 294529, 294530, 294531, 294553, 294556, 294558, 294562, 294563, 294586, 294587, 294588, 294590, 294591, 294593, 294620, 294632, 294633, 294634, 294636, 294653, 294703, 294710, 294715, 294716, 294718, 294721, 294722, 294741, 294742, 294743, 294744, 294745, 294746, 294747, 294749, 294769, 294772, 294774, 294775, 294805, 294809, 294813, 294815, 294816, 294817, 294818, 294852, 294875, 294877, 294880, 294883, 294884, 294916, 294925\\

%
%

\subsection{LHC16x, MB INT7}
\label{subsec:goodRuns16}

Using that selection of runs, the quality assurance has been performed as it is described in the respective chapters that can be found within this note. Although all different observables (see \hyperref[tab:trackSelection]{Table~\ref*{tab:trackSelection}},  \hyperref[tab:PID]{Table~\ref*{tab:PID}}, \hyperref[tab:PhotonCuts]{Table~\ref*{tab:PhotonCuts}} and \hyperref[tab:clusterSelection]{Table~\ref*{tab:clusterSelection}}) have been checked in great detail for their periodwise and runwise agreement, the most important observables to distinguish ``good`` and ``bad`` runs are: \vspace{-0.2cm}
\begin{itemize}
\item[-] Number of $\pi^{0}$ candidates per event, $\pi^{0}$ mass and width \vspace{-0.3cm}
\item[-] Conversion Photon Candidates (short: conversions) per event, \acs{EMCal} Photon Candidates (short: clusters) per event \vspace{-0.3cm}
\item[-] Energy distributions of conversions, clusters \vspace{-0.3cm}
\item[-] $\eta$/$\phi$ distributions of selected conversions, clusters \vspace{-0.3cm}
\item[-] ``Event properties `` like number of charged tracks, z-vertex distributions as well as fraction of events discarded due to pileup or missing primary vertices \vspace{-0.2cm}
\end{itemize}
If a run deviates from the global mean value of the respective observable for more than $2\sigma$, it is considered as ``bad``. This is being done unless there is no clear explanation for a different behavior like some modules of \acs{EMCal} being switched off, for example, which leads to a decrease in the number of good \acs{EMCal} photon candidates. Therefore, we always compared data and Monte Carlo as well to ensure that Monte Carlo describes the recorded data reasonably well. Also, if trending points in data and Monte Carlo do not show the same behavior, which means if values of observables change and if they are not properly followed by the other respective set, the run is also considered as ``bad`` unless the reason for the deviation cannot be clearly identified and fixed.\TODO{All}{Add Runlists}
%
% \begin{table}[h]
% \hspace*{-0.2cm}
% \small
% \centering
% \begin{tabular}{ll}
% \toprule
% \textbf{LHC16d - PCM Good Runs PCM INT7} \\ \midrule
%     252235, 252238, 252248, 252271, 252310, 252313, 252315, 252317, 252319, 252322, \\ \midrule
%     252325, 252326, 252330, 252332, 252336, 252368, 252370, 252371, 252374, 252375 \\
%     \bottomrule
%     \end{tabular}
%     \caption{Runs used in this analysis for LHC16x, minimum bias trigger INT7 in PCM. The set of runs is identical for data and corresponding Monte Carlo productions that are used.}
%     \label{tab:runs16d}
% \end{table}
%
% \begin{table}[h]
% \hspace*{-0.2cm}
% \small
% \centering
% \begin{tabular}{ll}
% \toprule
% \textbf{LHC16e - PCM Good Runs PCM INT7} \\ \midrule
%     252858, 252867, 253437, 253478, 253481, 253482, 253488, 253517, 253529, 253530, \\ %\midrule
%     253563, 253589, 253591 \\
%     \bottomrule
%     \end{tabular}
%     \caption{Runs used in this analysis for LHC16x, minimum bias trigger INT7 in PCM. The set of runs is identical for data and corresponding Monte Carlo productions that are used.}
%     \label{tab:runs16e}
% \end{table}

% 	\begin{table}[h]
% 	\hspace*{-0.2cm}
% 	\small
% 	\centering
% 	\begin{tabular}{ll}
% 	\toprule
% 	\textbf{LHC16g - PCM Good Runs PCM INT7} \\ \midrule
% 		254128, 254147, 254148, 254149, 254174, 254175, 254178, 254193, 254199, 254204, \\ \midrule
% 		254205, 254293, 254302, 254303, 254304, 254330, 254331, 254332 \\
% 		\bottomrule
% 		\end{tabular}
% 		\caption{Runs used in this analysis for LHC16x, minimum bias trigger INT7 in PCM. The set of runs is identical for data and corresponding Monte Carlo productions that are used.}
% 		\label{tab:runs16g}
% 	\end{table}
%
% 	\begin{table}[h]
% 	\hspace*{-0.2cm}
% 	\small
% 	\centering
% 	\begin{tabular}{ll}
% 	\toprule
% 	\textbf{LHC16h - PCM Good Runs PCM INT7} \\ \midrule
% 		254418, 254419, 254422, 254604, 254606, 254608, 254621, 254629, 254630, 254632, \\ \midrule
% 		254640, 254644, 254646, 254648, 254649, 254651, 254652, 254653, 254654, 255079, \\ \midrule
% 		255082, 255085, 255086, 255091, 255111, 255154, 255159, 255162, 255167, 255171, \\ \midrule
% 		255173, 255174, 255176, 255177, 255240, 255242, 255247, 255248, 255249, 255251, \\ \midrule
% 		255252, 255253, 255255, 255256, 255275, 255276, 255280, 255283, 255350, 255351, \\ \midrule
% 		255352, 255398, 255402, 255407, 255415, 255418, 255419, 255420, 255421, 255440, \\ \midrule
% 		255442, 255447, 255463, 255465, 255466, 255467 \\
% 		\bottomrule
% 		\end{tabular}
% 		\caption{Runs used in this analysis for LHC16x, minimum bias trigger INT7 in PCM. The set of runs is identical for data and corresponding Monte Carlo productions that are used.}
% 		\label{tab:runs16h}
% 	\end{table}
%
% 	\begin{table}[h]
% 	\hspace*{-0.2cm}
% 	\small
% 	\centering
% 	\begin{tabular}{ll}
% 	\toprule
% 	\textbf{LHC16i - PCM Good Runs PCM INT7} \\ \midrule
% 		255539, 255540, 255541, 255542, 255543, 255577, 255582, 255583, 255591, 255614, \\ \midrule
% 		255615, 255616, 255617, 255618 \\
% 		\bottomrule
% 		\end{tabular}
% 		\caption{Runs used in this analysis for LHC16x, minimum bias trigger INT7 in PCM. The set of runs is identical for data and corresponding Monte Carlo productions that are used.}
% 		\label{tab:runs16i}
% 	\end{table}
%
% 	\begin{table}[h]
% 	\hspace*{-0.2cm}
% 	\small
% 	\centering
% 	\begin{tabular}{ll}
% 	\toprule
% 	\textbf{LHC16j - PCM Good Runs PCM INT7} \\ \midrule
% 		256204, 256207, 256210, 256212, 256213, 256215, 256219, 256222, 256223, 256227, \\ \midrule
% 		256228, 256231, 256281, 256282, 256283, 256284, 256287, 256289, 256290, 256292, \\ \midrule
% 		256295, 256297, 256298, 256299, 256302, 256307, 256309, 256311, 256356, 256357, \\ \midrule
% 		256361, 256362, 256363, 256364, 256365, 256366, 256368, 256371, 256372, 256373, \\ \midrule
% 		256415, 256417, 256418, 256420 \\
% 		\bottomrule
% 		\end{tabular}
% 		\caption{Runs used in this analysis for LHC16x, minimum bias trigger INT7 in PCM. The set of runs is identical for data and corresponding Monte Carlo productions that are used.}
% 		\label{tab:runs16j}
% 	\end{table}
%
% 	\begin{table}[h]
% 	\hspace*{-0.2cm}
% 	\small
% 	\centering
% 	\begin{tabular}{ll}
% 	\toprule
% 	\textbf{LHC16k - PCM Good Runs PCM INT7} \\ \midrule
% 		256504, 256506, 256510, 256512, 256552, 256554, 256556, 256557, 256560, 256561, \\ \midrule
% 		256562, 256564, 256565, 256567, 256589, 256619, 256620, 256658, 256676, 256677, \\ \midrule
% 		256681, 256684, 256691, 256692, 256694, 256695, 256697, 256941, 256942, 256944, \\ \midrule
% 		257012, 257021, 257026, 257077, 257080, 257082, 257083, 257084, 257086, 257092, \\ \midrule
% 		257095, 257100, 257136, 257137, 257138, 257139, 257140, 257141, 257142, 257144, \\ \midrule
% 		257145, 257204, 257206, 257209, 257224, 257260, 257318, 257320, 257322, 257330, \\ \midrule
% 		257358, 257364, 257433, 257457, 257468, 257474, 257487, 257488, 257490, 257491, \\ \midrule
% 		257492, 257530, 257531, 257537, 257539, 257540, 257541, 257560, 257561, 257562, \\ \midrule
% 		257566, 257587, 257588, 257590, 257592, 257594, 257595, 257604, 257605, 257606, \\ \midrule
% 		257630, 257632, 257635, 257636, 257642, 257644, 257682, 257684, 257685, 257687, \\ \midrule
% 		257688, 257689, 257691, 257692, 257694, 257697, 257724, 257725, 257727, 257733, \\ \midrule
% 		257734, 257735, 257737, 257754, 257757, 257765, 257773, 257797, 257798, 257799, \\ \midrule
% 		257800, 257803, 257804, 257850, 257851, 257853, 257855, 257893, 257901, 257912, \\ \midrule
% 		257932, 257936, 257937, 257939, 257957, 257958, 257960, 257963, 257979, 257986, \\ \midrule
% 		257989, 257992, 258003, 258008, 258012, 258014, 258017, 258019, 258039, 258041, \\ \midrule
% 		258042, 258045, 258048, 258049, 258053, 258059, 258060, 258062, 258063, 258107, \\ \midrule
% 		258108, 258109, 258113, 258114, 258117, 258178, 258197, 258198, 258202, 258203, \\ \midrule
% 		258204, 258256, 258257, 258258, 258270, 258271, 258273, 258274, 258278, 258299, \\ \midrule
% 		258301, 258302, 258303, 258306, 258307, 258332, 258336, 258359, 258391, 258393, \\ \midrule
% 		258426, 258452, 258454, 258456, 258477, 258499, 258537 \\
% 		\bottomrule
% 		\end{tabular}
% 		\caption{Runs used in this analysis for LHC16x, minimum bias trigger INT7 in PCM. The set of runs is identical for data and corresponding Monte Carlo productions that are used.}
% 		\label{tab:runs16k}
% 	\end{table}
%
% 	\begin{table}[h]
% 	\hspace*{-0.2cm}
% 	\small
% 	\centering
% 	\begin{tabular}{ll}
% 	\toprule
% 	\textbf{LHC16l - PCM Good Runs PCM INT7} \\ \midrule
% 		258962, 258964, 259088, 259090, 259091, 259096, 259099, 259117, 259118, 259162, \\ \midrule
% 		259164, 259204, 259257, 259261, 259263, 259264, 259269, 259270, 259271, 259272, \\ \midrule
% 		259273, 259274, 259302, 259303, 259305, 259307, 259334, 259336, 259339, 259340, \\ \midrule
% 		259341, 259342, 259378, 259382, 259388, 259389, 259394, 259395, 259396, 259473, \\ \midrule
% 		259477, 259649, 259650, 259668, 259697, 259700, 259703, 259704, 259705, 259711, \\ \midrule
% 		259713, 259747, 259748, 259750, 259751, 259752, 259756, 259781, 259788, 259789, \\ \midrule
% 		259822, 259841, 259842, 259860, 259866, 259867, 259868, 259888 \\
% 		\bottomrule
% 		\end{tabular}
% 		\caption{Runs used in this analysis for LHC16x, minimum bias trigger INT7 in PCM. The set of runs is identical for data and corresponding Monte Carlo productions that are used.}
% 		\label{tab:runs16l}
% 	\end{table}
%
% 	\begin{table}[h]
% 	\hspace*{-0.2cm}
% 	\small
% 	\centering
% 	\begin{tabular}{ll}
% 	\toprule
% 	\textbf{LHC16o - PCM Good Runs PCM INT7} \\ \midrule
% 		262395, 262396, 262397, 262398, 262399, 262418, 262419, 262422, 262423, 262424, \\ \midrule
% 		262425, 262426, 262428, 262430, 262487, 262489, 262490, 262492, 262563, 263331, \\ \midrule
% 		263813, 263814, 263823, 263824, 263829, 263830, 263923, 263979 \\
% 		\bottomrule
% 		\end{tabular}
% 		\caption{Runs used in this analysis for LHC16x, minimum bias trigger INT7 in PCM. The set of runs is identical for data and corresponding Monte Carlo productions that are used.}
% 		\label{tab:runs16o}
% 	\end{table}
%
% 	\begin{table}[h]
% 	\hspace*{-0.2cm}
% 	\small
% 	\centering
% 	\begin{tabular}{ll}
% 	\toprule
% 	\textbf{LHC16p - PCM Good Runs PCM INT7} \\ \midrule
% 		264076, 264078, 264082, 264085, 264086, 264109, 264110, 264129, 264137, 264138, \\ \midrule
% 		264139, 264164, 264188, 264190, 264194, 264197, 264198, 264232, 264233, 264235, \\ \midrule
% 		264238, 264259, 264260, 264261, 264262, 264264, 264265, 264266, 264267, 264273, \\ \midrule
% 		264277, 264279, 264281, 264305, 264306, 264312, 264336, 264341, 264345, 264346, \\ \midrule
% 		264347 \\
% 		\bottomrule
% 		\end{tabular}
% 		\caption{Runs used in this analysis for LHC16x, minimum bias trigger INT7 in PCM. The set of runs is identical for data and corresponding Monte Carlo productions that are used.}
% 		\label{tab:runs16p}
% 	\end{table}
%

\clearpage
\section{Quality Assurance on Cell Level}
\label{chap:extCellQA}
\subsection{Periodwise Cell QA}
The periodwise quality assurance on cell level is being performed using a multiple set of cuts on the following observables:\vspace{-0.2cm}
\begin{itemize}
\item Integrated occurrences of cells in events, see \hyperref[fig:cellHotCells]{Figure~\ref*{fig:cellHotCells}} \vspace{-0.3cm}
\item Cell energy distribution - mean / sigma, see \hyperref[fig:cellEnergy]{Figure~\ref*{fig:cellEnergy}} \vspace{-0.3cm}
\item Cell time distribution - mean / sigma, see \hyperref[fig:cellTime]{Figure~\ref*{fig:cellTime}} \vspace{-0.3cm}
\item Cell time distribution - ratio of fired cells within $20~ns$ to $500~ns$, see \hyperref[fig:cellHotCellsTime1D]{Figure~\ref*{fig:cellHotCellsTime1D}} \vspace{-0.3cm}
\item Cell energy distribution - ratio of fired cells above threshold to below threshold, see \hyperref[fig:cellHotCells2D]{Figure~\ref*{fig:cellHotCells2D}} \vspace{-0.2cm}
\end{itemize}
Only cells which fulfill all mentioned cuts are accepted and being used in the clusterization process. If a cell has been found to be bad in one of the observables, it has been added to the bad cell lists in OADB/OCDB so that it will be automatically removed from all events before entering any analysis. The following figures contain the corresponding plots for LHC16i, which serves as an example. All data and Monte Carlo periods from this analysis have been checked.

% \begin{figure}[h]
%  	\centering
%  	\includegraphics[width=0.6\textwidth]{figures/ClusterQA/Cells/CellHotCells\_LHC16i\_pass1\_block1.pdf}
%  	\vspace{-0.3cm}
%  	\caption{This plot contains the number of occurrences for each cell summed over all events before applying any cuts in the analysis. If a cell is only hot but does not fire with higher energies, so that Figure \hyperref[fig:cellEnergy]{Figure~\ref*{fig:cellEnergy}} cannot detect the bad cell, this plot has the capability to cut away those noisy cells. The applied cuts are represented by the dotted lines, only cells within are kept.}
%  	\label{fig:cellHotCells}
%  \end{figure}
%
%  \begin{figure}[h]
%  	\vspace{-0.1cm}
% 	\centering
% 	\includegraphics[width=0.6\textwidth]{figures/ClusterQA/Cells/CellEnergyVsSigma\_LHC16i\_pass1\_block1.pdf}
% 	\vspace{-0.3cm}
% 	\caption{For each cell, the mean cell energy is plotted vs sigma energy in this histogram. If cell energy distribution has multiple peaks or if a cell is hot and firing very often with high energy, it will be clearly separated from the good cells which accumulate very well. The cuts are chosen so that no Monte Carlo cells would be cut. The applied cuts are represented by the dotted lines, only cells within are kept.}
% 	\label{fig:cellEnergy}
% \end{figure}
%
% \begin{figure}
% 	\centering
% 	\includegraphics[width=0.6\textwidth]{figures/ClusterQA/Cells/CellTimeVsSigma\_LHC16i\_pass1\_block1.pdf}
% 	\vspace{-0.3cm}
% 	\caption{For each cell, the mean cell time is plotted vs sigma time in this histogram. If cell time distribution has multiple peaks or if a cell is firing very often at random times, it will be clearly separated from the good cells which accumulate very well. The applied cuts are represented by the dotted lines, only cells within are kept.}
% 	\label{fig:cellTime}
% \end{figure}
%
% \begin{figure}[h]
% 	\centering
% 	\includegraphics[width=0.6\textwidth]{figures/ClusterQA/Cells/CellHotCellsTime1D\_LHC16i\_pass1\_block1.pdf}
% 	\vspace{-0.3cm}
% 	\caption{The cell time distribution is taken as basis and is being integrated from $-20~ns$ to $20~ns$. Then the cell time distribution is integrated over the whole time axis. The ratio of both quantities is plotted in this figure. Ratios over 1.5 indicate that the cell primarily fires at basically random times which is the reason to exclude those cells. The applied cuts are represented by the dotted lines, only cells within are kept.}
% 	\label{fig:cellHotCellsTime1D}
% \end{figure}
%
% \begin{figure}[h]
% 	\centering
% 	\includegraphics[width=0.6\textwidth]{figures/ClusterQA/Cells/CellHotCells2D\_LHC16i\_pass1\_block1.pdf}
% 	\vspace{-0.3cm}
% 	\caption{The cell energy distribution is integrated from $0$ to $L$ and then from $L$ to $30~GeV$. The ratios are plotted for different values of $L$ in this figure. If a cell energy distribution has multiple peaks or a very different shape than other cells, it will be clearly visible in this plot at higher ratios. The dotted lines represent the cuts that have been used to exclude bad cells while only cells with lower ratios as the dotted lines are kept.}
% 	\label{fig:cellHotCells2D}
% \end{figure}
%
% \clearpage

\subsection{Runwise Cell QA}
%
% A detailed quality assurance has also been performed on runwise level. In this case, dead/cold and hot cells are identified on a run by run basis and then being marked as bad cells in OADB/OCDB, accordingly. In \hyperref[fig:cellHotCells]{Figure~\ref*{fig:cellHotCells}}, there is a typical energy fraction plot shown which is used to identify dead/cold and hot cells. The cell ID is plotted on x-axis while the y-axis contains the summed cell energy fraction of the cluster energy, summed over all cluster in a event and over all used events. As it turned out, the summed cell energy fraction is a much better discriminator than just looking at the integrated number of occurrences of respective cell IDs in clusters, which, in principle, could be used as well for the bad cell identification.
%
% %  		\begin{figure}[h]
% % 			\centering
% % 			\includegraphics[width=0.75\textwidth]{figures/ClusterQA/Cells/LHC12a\_EFrac\_Cells\_before.eps}
% % 			\caption{The cell ID is plotted versus the cell energy fraction of clusters, summed over all clusters in a event and over all events used in the analysis. One can clearly see some cold and hot cells by eye. The plot has been produced before the removal of bad cell candidates that have been identified.}
% % 			\label{fig:}
% %  	\end{figure}
%
% To be able to identify dead/cold and hot cells, some reference value, $\overline{E}\_{\mbox{\tiny frac, ref}}^{\mbox{\tiny X}}$, has to be computed in order to make the decision if cell at ID 'X' may be bad or not. This value is obtained by averaging the summed energy fraction over the previous 100 good cells.
% \begin{equation}
% \overline{E}\_{\mbox{\tiny frac, ref}}^{\mbox{\tiny X}} = \frac{1}{100}\cdot\sum\_{j=A-100}^{A-1} E\_{\mbox{\tiny frac}}^{\mbox{\tiny j-th cell ID}}
% \end{equation}
% This average is then compared with cell 'X' in order to decide if it is bad or not. For the first 100 cells, it is not possible, obviously, to compute the average using the given formula. For those cells at the beginning, a special procedure is being performed. The average of the first 100 cells is computed and used for the determination. To remove the possible occurrence bad cells in this average, the list is reiterated and the following cells are discarded: if cells are included with an energy fraction of zero or if a cell has double of the mean energy fraction and at least an energy fraction of '10'. Then the next cell IDs, beginning from 101, are included in the reference value for the first 100 cell IDs.
%
% The following criteria apply when identifying bad cells on runwise level:
% \begin{itemize}
% \item []\textbf{Dead cells}:\\
% A cell with ID 'X' is declared as dead in a run, if
% \begin{equation}
% E\_{\mbox{\tiny frac}}^{\mbox{\tiny X}} = 0
% \end{equation}
% and if it is not dead in anchored Monte Carlo productions, according to the given definition, in order to identify only new dead cells which are not yet included in OADB/OCDB.
% \item []\textbf{Cold cells}:\\
% A cell with ID 'X' is declared as cold in a run, if
% \begin{eqnarray}
% \mbox{if}(E\_{\mbox{\tiny frac}}^{\mbox{\tiny X}} \geq 80)& E\_{\mbox{\tiny frac}}^{\mbox{\tiny X}} < \overline{E}\_{\mbox{\tiny frac, ref}}^{\mbox{\tiny X}}/3 \\
% \mbox{else if}(E\_{\mbox{\tiny frac}}^{\mbox{\tiny X}} \geq 40)& E\_{\mbox{\tiny frac}}^{\mbox{\tiny X}} < \overline{E}\_{\mbox{\tiny frac, ref}}^{\mbox{\tiny X}}/5 \\
% \mbox{else if}(E\_{\mbox{\tiny frac}}^{\mbox{\tiny X}} \geq 10)& E\_{\mbox{\tiny frac}}^{\mbox{\tiny X}} < \overline{E}\_{\mbox{\tiny frac, ref}}^{\mbox{\tiny X}}/8 \\
% \mbox{else if}(E\_{\mbox{\tiny frac}}^{\mbox{\tiny X}} < 10)& E\_{\mbox{\tiny frac}}^{\mbox{\tiny X}} < \overline{E}\_{\mbox{\tiny frac, ref}}^{\mbox{\tiny X}}/10
% \end{eqnarray}
% and if it is not cold in anchored Monte Carlo productions, according to the given definition, in order to identify only true cold cells.
%
% \item []\textbf{Hot cells}:\\
% A cell with ID 'X' is declared as hot in a run, if
% \begin{eqnarray}
% \mbox{if}(E\_{\mbox{\tiny frac}}^{\mbox{\tiny X}} > 80)& E\_{\mbox{\tiny frac}}^{\mbox{\tiny X}} < 2 \cdot \overline{E}\_{\mbox{\tiny frac, ref}}^{\mbox{\tiny X}} \\
% \mbox{else if}(E\_{\mbox{\tiny frac}}^{\mbox{\tiny X}} > 20)& E\_{\mbox{\tiny frac}}^{\mbox{\tiny X}} < 3 \cdot \overline{E}\_{\mbox{\tiny frac, ref}}^{\mbox{\tiny X}} \\
% \mbox{else if}(E\_{\mbox{\tiny frac}}^{\mbox{\tiny X}} > 8)& E\_{\mbox{\tiny frac}}^{\mbox{\tiny X}} < 4 \cdot \overline{E}\_{\mbox{\tiny frac, ref}}^{\mbox{\tiny X}} \\
% \mbox{else if}(E\_{\mbox{\tiny frac}}^{\mbox{\tiny X}} > 5)& E\_{\mbox{\tiny frac}}^{\mbox{\tiny X}} < 5 \cdot \overline{E}\_{\mbox{\tiny frac, ref}}^{\mbox{\tiny X}}
% \end{eqnarray}
% and if it is not hot in anchored Monte Carlo productions, according to the given definition, in order to identify only true hot cells.
% \end{itemize}
%
% While looping through all cell IDs, the cells, that have been identified to belong to one of the three just described categories, will not enter the ongoing average determination for the upcoming cell IDs.
%
% For the dead and cold cell determination, only runs are inspected that have at least $1\cdot10^5$ fired cells integrated over the whole run. For the hot cell determination, the threshold to inspect a run is $1\cdot10^4$ to ensure enough statistics to perform a meaningful selection. If there are clear sub-runranges where some cells are bad, they will be marked as bad in those runs including the skipped runs due to statistics. If a cell is suspicious in the majority of runs, it will be directly excluded from the whole period.
% \clearpage

\section{Additional Cluster/Photon QA Plots}
\label{chap:QAClusterShape}
\label{chap:extQALHC16}

\hyperref[fig:clusterM02]{Figure~\ref*{fig:clusterM02}} shows the cluster shape parameter $\lambda\_{0}^{2}$ including a comparison of data and Monte Carlo.
Moreover, runwise distributions are also shown for all productions.
The corresponding cluster shape parameter $\lambda\_{1}^{2}$ can be found in \hyperref[fig:clusterM20]{Figure~\ref*{fig:clusterM20}}, including a comparison of data and Monte Carlo.
Additionally, runwise distributions are shown for data and both minimum bias simulations for the example period LHC12b.

% \begin{figure}[h]
% 		\centering
% 		\includegraphics[width=0.44\textwidth]{figures/TriggerQA/M02\_afterQA.eps}
% 		\hspace{0.6cm}
% 		\includegraphics[width=0.48\textwidth]{figures/ClusterQA/Runwise/LHC12b/LHC12b/ClusterM02\_Runwise\_LHC12b.eps}\\
% 		\includegraphics[width=0.48\textwidth]{figures/ClusterQA/Runwise/LHC12b/LHC15h1b/ClusterM02\_Runwise\_Pythia8.eps}
% 		\includegraphics[width=0.48\textwidth]{figures/ClusterQA/Runwise/LHC12b/LHC15h2b/ClusterM02\_Runwise\_Phojet.eps}
% 		\caption{(top left) Cluster shape parameter $\lambda\_{0}^{2}$ for \acs{EMCal} clusters for data, Pythia8, Phojet and JetJet Monte Carlo. All points are scaled by the number of available events in each production. Only parameters between 0.1 and 0.5 are accepted as it can be taken from \hyperref[tab:clusterSelection]{Table~\ref*{tab:clusterSelection}}. The other three histograms show runwise superpositions of the $\lambda\_{0}^{2}$ distributions for data, Pythia8 and Phojet that are also scaled by the number of available events in each run. }
% 		\label{fig:clusterM02}
%  \end{figure}
%  \begin{figure}[h]
% 		\centering
% 		\includegraphics[width=0.44\textwidth]{figures/TriggerQA/M20\_afterQA.eps}
% 		\hspace{0.6cm}
% 		\includegraphics[width=0.48\textwidth]{figures/ClusterQA/Runwise/LHC12b/LHC12b/ClusterM20\_Runwise\_LHC12b.eps}\\
% 		\includegraphics[width=0.48\textwidth]{figures/ClusterQA/Runwise/LHC12b/LHC15h1b/ClusterM20\_Runwise\_Pythia8.eps}
% 		\includegraphics[width=0.48\textwidth]{figures/ClusterQA/Runwise/LHC12b//LHC15h2b/ClusterM20\_Runwise\_Phojet.eps}
% 		\caption{(top left) Cluster shape parameter $\lambda\_{1}^{2}$ for \acs{EMCal} clusters in data, Pythia8, Phojet and JetJet Monte Carlo. All points are scaled by the number of available events in each production. As it can be taken from \hyperref[tab:clusterSelection]{Table~\ref*{tab:clusterSelection}}, there is only a cut on $\lambda\_{0}^{2}$. Since $\lambda\_{1}^{2}$ and $\lambda\_{0}^{2}$ are correlated the possible values of $\lambda\_{1}^{2}$ are being limited to 0.5. The other three histograms show runwise superpositions of the $\lambda\_{1}^{2}$ distributions for data, Pythia 8 and Phojet that are also scaled by the number of available events in each run. }
% 		\label{fig:clusterM20}
% \end{figure}

The data points follow Monte Carlo equivalents reasonably well, which nicely agree to each other and all runs show the same distributions.
Both the average and the standard deviation of the shower shape parameters $\lambda\_{0}^{2}$ and $\lambda\_{1}^{2}$ stay constant over all available runs as it can be seen in \hyperref[fig:QA-M02]{Figure~\ref*{fig:QA-M02}} and \hyperref[fig:QA-M20]{Figure~\ref*{fig:QA-M20}}.
Both Monte Carlo productions agree while again data is at slightly higher values.
This behavior makes us confident that the cluster shape is reproduced reasonably well in Monte Carlo.

Finally, in \hyperref[fig:clusterM02M20]{Figure~\ref*{fig:clusterM02M20}} a 2-dimensional plot of both shape parameters $\lambda\_{1}^{2}$ and $\lambda\_{0}^{2}$ can be found.

% 	\begin{figure}[h]
%			\centering
%			\includegraphics[width=0.48\textwidth]{figures/ClusterQA/Runwise/hClusterM02-Mean.eps}
%			\includegraphics[width=0.48\textwidth]{figures/ClusterQA/Runwise/hClusterM02-RMS.eps}
%			\caption{(left) Average value of shower shape parameter $\lambda\_{0}^{2}$. (right) Standard deviation of $\lambda\_{0}^{2}$.}
%			\label{fig:QA-M02}
%  	\end{figure}
%
% 	 \begin{figure}[h]
%			\centering
%			\includegraphics[width=0.48\textwidth]{figures/ClusterQA/Runwise/hClusterM20-Mean.eps}
%			\includegraphics[width=0.48\textwidth]{figures/ClusterQA/Runwise/hClusterM20-RMS.eps}
%			\caption{(left) Average value of shower shape parameter $\lambda\_{1}^{2}$. (right) Standard deviation of $\lambda\_{1}^{2}$.}
%			\label{fig:QA-M20}
%  	\end{figure}
%
% \begin{figure}[h]
%		\centering
%		\includegraphics[width=0.5\textwidth]{figures/ClusterQA/M20VsM02\_all\_LHC12.eps}
%		\caption{A two dimensional plot of $\lambda\_{1}^{2}$ vs. $\lambda\_{0}^{2}$ is shown here. The most common cluster shapes can be well distinguished as they form different peaks with very different shapes in the histogram. Since we only use very loose cuts and do not explicitly use cluster shape analysis to enhance the photon sample, no further discussions on that topic are made in this note.}
%		\label{fig:clusterM02M20}
%\end{figure}
%
%\begin{figure}[h]
%		\centering
%		\includegraphics[width=0.48\textwidth]{figures/TriggerQA/ratios/ratio\_Energy\_Cluster\_afterQA.eps}
%		\includegraphics[width=0.48\textwidth]{figures/TriggerQA/ratios/ratio\_NCells\_afterQA.eps}
%		\caption{Concerning \hyperref[fig:clusterEnergy]{Figures~\ref*{fig:clusterEnergy}} and \hyperref[fig:clusterNCells]{\ref*{fig:clusterNCells}}, the corresponding ratios can be found in the left and right plot accordingly.}
%		\label{fig:clusterQAratios}
%\end{figure}
%
%\begin{figure}[h]
%		\centering
%		\includegraphics[width=0.48\textwidth]{figures/TriggerQA/ratios/ratio\_NGoodTracks.eps}
%		\includegraphics[width=0.48\textwidth]{figures/TriggerQA/ratios/ratio\_GammaCandidates\_EMCAL.eps}
%		\caption{Ratios of charged track multiplicities (left) and \acs{EMCal} ratios of cluster multiplicities (right) for LHC12a-i with the corresponding minimum bias productions and JetJet MC productions as well.}
%		\label{fig:clusterQAratios2}
%\end{figure}
%
%\begin{figure}[h]
%		\centering
%		\includegraphics[width=0.48\textwidth]{figures/TriggerQA/ratios/ratio\_NGoodTracks\_EMC7.eps}
%		\includegraphics[width=0.48\textwidth]{figures/TriggerQA/ratios/ratio\_NGoodTracks\_EGA.eps}
%		\caption{Ratios of charged track multiplicities for EMC L0, INT7 (left) and for EMC L1-GA, INT7 (right).}
%		\label{fig:clusterQAratios3}
%\end{figure}
%
%\begin{figure}[h]
%		\centering
%		\includegraphics[width=0.48\textwidth]{figures/PhotonQA/Comparison/Ratios/ratio\_Electron\_FindTPCClusters.eps}
%		\includegraphics[width=0.48\textwidth]{figures/PhotonQA/Comparison/Ratios/ratio\_Positron\_FindTPCClusters.eps}\\							\includegraphics[width=0.48\textwidth]{figures/PhotonQA/Comparison/Ratios/ratio\_Electron\_TPCClusters.eps}
%		\includegraphics[width=0.48\textwidth]{figures/PhotonQA/Comparison/Ratios/ratio\_Positron\_TPCClusters.eps}
%		\caption{Ratios of \hyperref[fig:TPCfindableClusters]{Figure~\ref*{fig:TPCfindableClusters}} and \hyperref[fig:TPCClusters]{Figure~\ref*{fig:TPCClusters}}.}
%		\label{fig:photonQAratios}
%\end{figure}
%
%\begin{figure}[h]
%		\centering
%		\includegraphics[width=0.48\textwidth]{figures/TriggerQA/EMC7/ratio\_GammaCandidates\_EMCAL.eps}
%		\includegraphics[width=0.48\textwidth]{figures/TriggerQA/EGA/ratio\_GammaCandidates\_EMCAL.eps}
%		\caption{Ratios of \acs{EMCal} cluster multiplicities for EMC L0, INT7 (left) and for EMC L1-GA, INT7 (right).}
%		\label{fig:clusterQAratios5}
%\end{figure}
%
%\begin{figure}[h]
%		\centering
%			\includegraphics[width=0.49\textwidth]{figures/ClusterQA/EtaVsPhi\_RatioLHC15h1\_LHC12\_afterClusterQA.eps}
%			\includegraphics[width=0.49\textwidth]{figures/ClusterQA/EtaVsPhi\_RatioLHC15h2\_LHC12\_afterClusterQA.eps}\\
%			\includegraphics[width=0.49\textwidth]{figures/TriggerQA/EtaVsPhi\_RatioLHC16c2\_LHC12ci\_afterClusterQA.eps}
%		\caption{Ratios of \hyperref[fig:clusterEtaPhi]{Figure~\ref*{fig:clusterEtaPhi}}.}
%		\label{fig:clusterQAratios4}
%\end{figure}
%
%\begin{figure}[h]
%		\centering
%			\includegraphics[width=0.49\textwidth]{figures/EventQA/Comparison/Vertex_Z.eps}
%			\includegraphics[width=0.49\textwidth]{figures/EventQA/Comparison/Ratios/ratio_Vertex_Z.eps}
%		\caption{Z-vertex positions for data and minimum bias Monte Carlo compared for all accepted events. Furthermore, ratios to data are shown.}
%		\label{fig:eventQAzvertex}
%\end{figure}
%
%\clearpage
%
%\begin{figure}[h]
%	\centering
%	\includegraphics[width=0.94\textwidth]{figures/ClusterQA/Runwise/LHC12b/LHC12b/hClusterEtaPhi\_scaledNEventsAndMean\_p0.eps}\\
%	\includegraphics[width=0.94\textwidth]{figures/ClusterQA/Runwise/LHC12b/LHC12b/hClusterEtaPhi\_scaledNEventsAndMean\_p1.eps}
%	\caption{Runwise \acs{EMCal} cluster $\eta$-$\phi$ distributions for data, given the example of LHC12b. The number of available events for each single run can be found in left part of figure \ref{fig:QA-NClusters}.}
%	\label{runwiseEtaPhi:data}
%\end{figure}
%\begin{figure}[h]
%	\centering
%	\includegraphics[width=0.94\textwidth]{figures/ClusterQA/Runwise/LHC12b/LHC15h1b/hClusterEtaPhi\_scaledNEventsAndMean\_p0.eps}\\
%	\includegraphics[width=0.94\textwidth]{figures/ClusterQA/Runwise/LHC12b/LHC15h1b/hClusterEtaPhi\_scaledNEventsAndMean\_p1.eps}
%	\caption{Runwise \acs{EMCal} cluster $\eta$-$\phi$ distributions for Pythia8, according to example run list of LHC12b. The number of available events for each single run can be found in left part of figure \ref{fig:QA-NClusters}.}
%	\label{runwiseEtaPhi:pythia}
%\end{figure}
%
%\newpage
%
%\begin{figure}[h]
%	\centering
%	\includegraphics[width=0.94\textwidth]{figures/ClusterQA/Runwise/LHC12b/LHC15h2b/hClusterEtaPhi\_scaledNEventsAndMean\_p0.eps}\\
%	\includegraphics[width=0.94\textwidth]{figures/ClusterQA/Runwise/LHC12b/LHC15h2b/hClusterEtaPhi\_scaledNEventsAndMean\_p1.eps}
%	\caption{Runwise \acs{EMCal} cluster $\eta$-$\phi$ distributions for Phojet, according to example run list of LHC12b. The number of available events for each single run can be found in left part of figure \ref{fig:QA-NClusters}.}
%	\label{runwiseEtaPhi:phojet}
%\end{figure}
%\clearpage
%
% 	 \begin{figure}[h]
%			\centering
%			\includegraphics[width=0.48\textwidth]{figures/ClusterQA/ClusterEnergyVsModule\_LHC12.eps}
%			\includegraphics[width=0.48\textwidth]{figures/ClusterQA/ClusterEnergyVsModule\_Projected\_LHC12.eps}
%			\caption{On the left side, cluster energies are plotted versus supermodule number for INT7. The z-axis displays the number of occurrences of the respective cluster energies for each supermodule summed over all events. On the right side, the projection to the x-axis is shown in order to being able to directly compare the supermodules.}
% 	\end{figure}
%
% 	 \begin{figure}[h]
%			\centering
%			\includegraphics[width=0.48\textwidth]{figures/TriggerQA/EMC7/ClusterEnergyVsModule\_LHC12.eps}
%			\includegraphics[width=0.48\textwidth]{figures/TriggerQA/EMC7/ClusterEnergyVsModule\_Projected\_LHC12.eps}
%			\caption{On the left side, cluster energies are plotted versus supermodule number for EMC L0, INT7. The z-axis displays the number of occurrences of the respective cluster energies for each supermodule summed over all events. On the right side, the projection to the x-axis is shown in order to being able to directly compare the supermodules.}
% 	\end{figure}
%
% 	 	 \begin{figure}[h]
%			\centering
%			\includegraphics[width=0.48\textwidth]{figures/TriggerQA/EGA/ClusterEnergyVsModule\_LHC12.eps}
%			\includegraphics[width=0.48\textwidth]{figures/TriggerQA/EGA/ClusterEnergyVsModule\_Projected\_LHC12.eps}
%			\caption{On the left side, cluster energies are plotted versus supermodule number for EMC L1-GA, INT7. The z-axis displays the number of occurrences of the respective cluster energies for each supermodule summed over all events. On the right side, the projection to the x-axis is shown in order to being able to directly compare the supermodules.}
% 	\end{figure}
%
%     \begin{figure}[h]
%			\centering
%			\includegraphics[width=0.48\textwidth]{figures/ClusterQA/ModuleEnergyVsModule\_LHC12.eps}
%			\includegraphics[width=0.48\textwidth]{figures/ClusterQA/ModuleEnergyVsModule\_Projected\_LHC12\_LOG.eps}
%			\caption{On the left side, the total energy, which is being registered within a supermodule by summing over all energies of all available cells in an event, are plotted versus supermodule number. The z-axis displays the number of occurrences of the respective total energies for each supermodule summed over all events. On the right side, the projection to the x-axis is shown in order to being able to directly compare the supermodules.}
% 	\end{figure}
%
%    \begin{figure}[h]
%			\centering
%			\includegraphics[width=0.48\textwidth]{figures/ClusterQA/NCellsAbove100VsModule\_LHC12.eps}
%			\includegraphics[width=0.48\textwidth]{figures/ClusterQA/NCellsAbove100VsModule\_Projected\_LHC12.eps}
%			\caption{On the left side, the number of cells above 100 MeV, which are being registered within a supermodule in an event, are plotted versus supermodule number. The z-axis displays the number of occurrences of the respective total number of cells for each supermodule summed over all events. On the right side, the projection to the x-axis is shown in order to being able to directly compare the supermodules.}
% 	\end{figure}
%
%   \begin{figure}[h]
%			\centering
%			\includegraphics[width=0.48\textwidth]{figures/ClusterQA/NCellsAbove1500VsModule\_LHC12.eps}
%			\includegraphics[width=0.48\textwidth]{figures/ClusterQA/NCellsAbove1500VsModule\_Projected\_LHC12.eps}
%			\caption{On the left side, the number of cells above 1500 MeV, which are being registered within a supermodule in an event, are plotted versus supermodule number. The z-axis displays the number of occurrences of the respective total number of cells for each supermodule summed over all events. On the right side, the projection to the x-axis is shown in order to being able to directly compare the supermodules.}
% 	\end{figure}
%
% 	 	 \begin{figure}[h]
%			\centering
%			\includegraphics[width=0.47\textwidth]{figures/TriggerQA/Compare/ClusterEnergy.eps}
%			\includegraphics[width=0.51\textwidth]{figures/TriggerQA/Compare/hClusterEnergy-Mean.eps}
%			\caption{(left) Energy distributions of \acs{EMCal} clusters for minimum bias and triggered data are shown here. (right) This plot displays runwise the mean energy of \acs{EMCal} clusters for all available runs.}
%			\label{fig:clusterEnergyTrigger}
% 	\end{figure}
%
% 	\begin{figure}[h]
%			\centering
%			\includegraphics[width=0.47\textwidth]{figures/TriggerQA/Compare/NCells.eps}
%			\includegraphics[width=0.51\textwidth]{figures/TriggerQA/Compare/hClusterNCells-Mean.eps}
%			\caption{The left histogram shows the number of cells of \acs{EMCal} clusters for minimum bias and triggered data. Following on the right side, the respective runwise mean number of cells per cluster is displayed for the different triggers.}
%			\label{fig:clusterNCellsTrigger}
% 	\end{figure}
%
% 	\begin{figure}[h]
%			\centering
%			\includegraphics[width=0.47\textwidth]{figures/TriggerQA/Compare/M02.eps}
%			\includegraphics[width=0.51\textwidth]{figures/TriggerQA/Compare/hClusterM02-Mean.eps}
%			\caption{(left) Cluster shape parameter $\lambda\_{0}^{2}$ for \acs{EMCal} clusters for minimum bias and triggered data. All points are scaled by the number of available events in each production. Only parameters between 0.1 and 0.5 are accepted as it can be taken from \hyperref[tab:clusterSelection]{Table~\ref*{tab:clusterSelection}}. The other histogram shows the runwise superposition of the $\lambda\_{0}^{2}$ distributions for triggered data.}
%			\label{fig:clusterM02Trigger}
% 	\end{figure}
%
% 	 \begin{figure}[h]
%			\centering
%			\includegraphics[width=0.47\textwidth]{figures/TriggerQA/Compare/M20.eps}
%			\includegraphics[width=0.51\textwidth]{figures/TriggerQA/Compare/hClusterM20-Mean.eps}
%			\caption{(left) Cluster shape parameter $\lambda\_{1}^{2}$ for \acs{EMCal} clusters for minimum bias and triggered data. All points are scaled by the number of available events in each production. As it can be taken from \hyperref[tab:clusterSelection]{Table~\ref*{tab:clusterSelection}}, there is only a cut on $\lambda\_{0}^{2}$. Since $\lambda\_{1}^{2}$ and $\lambda\_{0}^{2}$ are correlated the possible values of $\lambda\_{1}^{2}$ are being limited to 0.5. The other histogram shows the runwise superposition of the $\lambda\_{1}^{2}$ distributions for triggered data.}
%			\label{fig:clusterM20Trigger}
% 	\end{figure}
%
% 	\begin{figure}[h]
%			\centering
%			\includegraphics[width=0.49\textwidth]{figures/TriggerQA/Compare/hCluster-FractionMatches.eps}
%			\includegraphics[width=0.49\textwidth]{figures/TriggerQA/Compare/hCluster-FractionMatchesH.eps}
%			\caption{(left) The fraction of matches is shown for each single run used in this analysis for minimum bias and triggered data respectively. The fraction is calculated by dividing the number of successful cluster -- V$^{0}$-matches by the total number of available mother candidates.(right) The fraction of matches is shown for $\it{p}\_{\text{T, Pair}}>2.5$~GeV/$\it{c}$ for each single run used in this analysis for minimum bias and triggered data respectively. The fraction is calculated by dividing the number of successful cluster -- V$^{0}$-matches above the $\it{p}\_{\text{T}}$-threshold by the total number of available mother candidates in that interval.}
%			\label{fig:QA-FractionMatchesTrigger}
% 	\end{figure}
%
% 	\begin{figure}[h]
%			\centering
%			\includegraphics[width=0.45\textwidth]{figures/TriggerQA/Compare/dEta.eps}
%			\includegraphics[width=0.45\textwidth]{figures/TriggerQA/Compare/dPhi.eps}\\
%			\includegraphics[width=0.45\textwidth]{figures/TriggerQA/Compare/dEtaMatched.eps}
%			\includegraphics[width=0.45\textwidth]{figures/TriggerQA/Compare/dPhiMatched.eps}
%			\caption{(top) $\Delta\eta$- and $\Delta\phi$-projections of \hyperref[fig:matchingAllAccepted]{Figure~\ref*{fig:matchingAllAccepted}} for minimum bias and triggered data. (bottom) For minimum bias and triggered data, only the projections of successful matches are shown. All distributions in this figure are respectively scaled to an integral of 1.}
%			\label{fig:matchingDeltaTrigger}
% 	\end{figure}
%
% 	\begin{figure}[h]
%			\centering
%			\includegraphics[width=0.48\textwidth]{figures/TriggerQA/Compare/hClusterDeltaEta-Mean.eps}
%			\includegraphics[width=0.48\textwidth]{figures/TriggerQA/Compare/hClusterDeltaEta-RMS.eps}\\
%			\includegraphics[width=0.48\textwidth]{figures/TriggerQA/Compare/hClusterDeltaPhi-Mean.eps}
%			\includegraphics[width=0.48\textwidth]{figures/TriggerQA/Compare/hClusterDeltaPhi-RMS.eps}\\
%			\caption{On the left side, the average values of $\Delta\eta$ and $\Delta\phi$ are plotted for every run for triggered data while the two histograms on the right side contain the mean standard deviations of the respective observables. All points are obtained by looking at the $\Delta\eta$- and $\Delta\phi$-projections of successful matches.}
%			\label{fig:QA-DeltaEtaPhiTrigger}
% 	\end{figure}
%
% 	\begin{figure}[h]
%			\centering
%			\includegraphics[width=0.48\textwidth]{figures/TriggerQA/Compare/MotherPi0Conv\_Eta.eps}
%			\includegraphics[width=0.48\textwidth]{figures/TriggerQA/Compare/MotherPi0Conv\_Phi.eps}
%			\caption{(left) Comparison of minimum bias and triggered data for $\eta$ distribution of selected conversion photon candidates under the $\pi^{0}$-peak. (right) Comparison of minimum bias and triggered data for $\phi$ distribution of selected conversion photon candidates under the $\pi^{0}$-peak. All histograms in both plots are scaled by the respective number of accepted events.}
%			\label{fig:etaphiPi0UnderpeakTrigger}
% 	\end{figure}
%
% 	\begin{figure}[h]
%			\centering
%			\includegraphics[width=0.48\textwidth]{figures/TriggerQA/Compare/MotherEtaConv\_Eta.eps}
%			\includegraphics[width=0.48\textwidth]{figures/TriggerQA/Compare/MotherEtaConv\_Phi.eps}
%			\caption{(left) Comparison of minimum bias and triggered data for $\eta$ distribution of selected conversion photon candidates under the $\eta$-peak. (right) Comparison of minimum bias and triggered data for $\phi$ distribution of selected conversion photon candidates under the $\eta$-peak. All histograms in both plots are scaled by the respective number of accepted events.}
%			\label{fig:etaphiEtaUnderpeakTrigger}
% 	\end{figure}

\clearpage


\section{Additional Neutral Meson Plots}
\label{chap:AppNeutralMeson}

% 	\begin{figure}[h]
% 		\centering
% 			\includegraphics[width=0.5\textwidth]{figures/Meson/MesonTrigger/Pi0_data_CorrectedYieldFinal_WithMC.eps}\\
% 			\includegraphics[width=0.6\textwidth]{figures/Meson/MesonTrigger/Pi0_data_RatioSpectraToFitFinal_withMC.eps}
% 			\caption{Final measured neutral pion transverse momentum spectra with \ac{EMCal} for \pp collisions at \sth compared to the Monte Carlo input spectra (blue). Indeed, the input spectrum is a realistic one with a reasonable power law behavior, so that it is suitable for this analysis.}
% 			\label{fig:EMCalWithMC}
% 		\end{figure}

% \clearpage

\subsection{\texorpdfstring{Neutral Mesons in \pp Collisions at \s~=~13~TeV, MB INT7}{Neutral Mesons in pp Collisions at 13 TeV, MB INT7}}

% 		\begin{figure}[h]
% 		 	\includegraphics[width=0.96\textwidth]	{figures/Meson/INT7new/Pi0_data_MesonWithBck_00010113_00200009327000008250400000_1111111067032230000_0163103100000010.eps}\\
% 			\includegraphics[width=0.96\textwidth]{figures/Meson/INT7new/Pi0_data_MesonSubtracted_00010113_00200009327000008250400000_1111111067032230000_0163103100000010.eps}
% 			\caption{Invariant-mass distribution of reconstructed photon pairs $M_{\gamma\gamma}$ around the neutral pion mass in \pT slices in pp collisions at \sth  before (upper plot) and after subtraction (lower plot) for minimum bias INT7 triggered data. The black histograms in the upper plot show the combined signal and background distribution and the blue histograms show the calculated and normalized mixed event background. After the subtraction, the invariant mass distributions are fitted with \hyperref[eq:gausexp]{Equation~\ref*{eq:gausexp}}, which is shown in green in the lower plot.}
% 			\label{fig:InvMassPi0PtSliceINT7}
% 		\end{figure}
%
% 		\begin{figure}
% 		 	\includegraphics[width=0.96\textwidth]	{figures/Meson/INT7new/Eta_data_MesonWithBck_00010113_00200009327000008250400000_1111111067032230000_0163103100000010.eps}\\
% 			\includegraphics[width=0.96\textwidth]{figures/Meson/INT7new/Eta_data_MesonSubtracted_00010113_00200009327000008250400000_1111111067032230000_0163103100000010.eps}
% 			\caption{Invariant-mass distribution of reconstructed photon pairs $M_{\gamma\gamma}$ around the eta meson mass in \pT slices in pp collisions at \sth  before (upper plot) and after subtraction (lower plot) for minimum bias INT7 triggered data. The black histograms in the upper plot show the combined signal and background distribution and the blue histograms show the calculated and normalized mixed event background. After the subtraction, the invariant mass distributions are fitted with \hyperref[eq:gausexp]{Equation~\ref*{eq:gausexp}}, which is shown in green in the lower plot.}
% 			\label{fig:InvMassEtaPtSliceINT7}
% 		\end{figure}
%
% 		\begin{figure}
% 		 	\includegraphics[width=0.96\textwidth]	{figures/Meson/INT7new/Pi0_MesonLineShapeCompared_00010113_00200009327000008250400000_1111111067032230000_0163103100000010.eps}\\
% 			\includegraphics[width=0.96\textwidth]{figures/Meson/INT7new/Eta_MesonLineShapeCompared_00010113_00200009327000008250400000_1111111067032230000_0163103100000010.eps}
% 			\caption{Invariant-mass distribution of reconstructed photon pairs $M_{\gamma\gamma}$ around the neutral pion (upper) and eta meson (lower) mass in \pT slices in pp collisions at \sth after subtraction. The black dots represent the distribution obtained after the background subtraction in the data, while the blue dots show the same distribution for the reconstructed Monte Carlo sample. On top of these, the real neutral pion and eta distributions, according to the Monte Carlo, are plotted as red histograms, those were obtained by validating that both reconstructed photons are actually photons (or decay products of one) and originate in the same meson, which has to be a neutral pion or eta meson. All distributions are normalized the integral in the given invariant mass range.}
% 			\label{fig:InvMassShapeComp2760}
% 		\end{figure}
%
% 		\begin{figure}
% 		 	\includegraphics[width=0.96\textwidth]{figures/Meson/INT7new/Pi0_MC_TrueMesonDecomposedPhotonsAndElectron_00010113_00200009327000008250400000_1111111067032230000_0163103100000010.eps}
% 		 	\includegraphics[width=0.96\textwidth]	{figures/Meson/INT7new/Pi0_MC_TrueMesonDecomposedMerged_00010113_00200009327000008250400000_1111111067032230000_0163103100000010.eps}
% 			\caption{Invariant-mass distribution of reconstructed validated photon pairs $M_{\gamma\gamma}$ around the mass of the neutral pion mass (0.135~GeV/$c^2$) in \pT slices in pp collisions at \sth. The upper distributions have been obtained by using the Monte Carlo simulation to validate that the two photons come from the same neutral meson (black bullets). Afterwards it has been checked whether one (purple) or both (cyan) of the reconstructed photons in the calorimeter converted or if both clusters were from a real photon (red).
% 			For the lower plots, the same convention applies for the black points. However, this time is was checked whether the cluster reconstructed in the \ac{EMCal} contains more than one particle from the same neutral pion. In this case, the red bullets show the distribution where there are at least two particles from the same neutral pion in the same cluster, while the blue distribution in addition requires that at least one of them was originating in a conversion.}
% 			\label{fig:InvMassPi0PtSliceMCDec2760}
% 		\end{figure}
% 		\begin{figure}
% 			\includegraphics[width=0.96\textwidth]{figures/Meson/INT7new/Eta_MC_TrueMesonDecomposedPhotonsAndElectron_00010113_00200009327000008250400000_1111111067032230000_0163103100000010.eps}
% 			\includegraphics[width=0.96\textwidth]	{figures/Meson/INT7new/Eta_MC_TrueMesonDecomposedMerged_00010113_00200009327000008250400000_1111111067032230000_0163103100000010.eps}
% 			\caption{Invariant-mass distribution of reconstructed validated photon pairs $M_{\gamma\gamma}$ around the mass of the $\eta$ meson mass (0.548~GeV/$c^2$) in \pT slices in pp collisions at \sth. The upper distributions have been obtained by using the Monte Carlo simulation to validate that the two photons come from the same neutral meson (black bullets). Afterwards it has been checked whether one (purple) or both (cyan) of the reconstructed photons in the calorimeter converted or if both clusters were from a real photon (red). For the lower plots, the same convention applies for the black points. However, this time is was checked whether the cluster reconstructed in the \ac{EMCal} contains more than one particle from the same eta meson. In this case, the red bullets show the distribution where there are at least two particles from the same eta meson in the same cluster, while the blue distribution in addition requires that at least one of them was originating in a conversion.}
% 			\label{fig:InvMassEtaPtSliceMCDec2760}
% 		\end{figure}
%
% 			\begin{figure}
% 			\includegraphics[width=0.49\textwidth]{figures/Meson/INT7new/Weighting/Pi0_data_CorrectedYield.eps}
% 			\includegraphics[width=0.49\textwidth]{figures/Meson/INT7new/Weighting/Eta_data_CorrectedYield.eps}
% 				\caption{Corrected neutral meson spectra for the minimum bias INT7 trigger in \sth. For both mesons (pion on left side, eta on right side), the fully corrected spectra are compared using the multiplicity weighted correction factors and the ones without any modification.}
% 			\label{fig:CompMultWeight}
% 			\end{figure}
%
% 			\begin{figure}
% 			\centering
% 			\includegraphics[width=1\textwidth]{figures/Meson/INT7new/Timing/Pi0_data_MesonSubtracted_00000113_1111111003032220000_0163103100000050.eps}
% 			\caption{This figure shows the background subtracted invariant mass distributions for an \acs{EMCal}-\acs{EMCal} analysis where one cluster was taken from within the cluster time cut window and paired to another cluster that had to be outside the timing cut window. One can clearly see that there is no yield at all visible and therefore, all fits are failing. Thus, the cluster time cut is fully efficient and there is no yield found due to contributions from outside the timing cut.}
% 			\label{fig:EMCALInOutTiming}
%  			\end{figure}
% 		\clearpage


		\subsection{Detailed Systematics}
		\label{sec:SysDetailed}

			\begin{figure}[t]
				\centering
				\includegraphics[width=0.49\textwidth]{figures/Meson/Systematics/EMC/SysMeanNewWithMean_Pi0_13TeVINT7.png}
				\includegraphics[width=0.49\textwidth]{figures/Meson/Systematics/EMC/SysMeanNewWithMean_Eta_13TeVINT7.png}
				\includegraphics[width=0.49\textwidth]{figures/Meson/Systematics/EMC/SysMeanNewWithMean_EtaToPi0_13TeVINT7.png}
				\caption{from the \ac{EMCal} method for \pp collisions at \s~=~13~TeV for the INT7 trigger. The colored points represent the individual error sources, while the black points represent the final systematic error for the $\pi^0$ meson transverse momentum spectrum.}
				\label{fig:SysErrEMCDetailed_INT7}
			\end{figure}
			\begin{figure}[t]
				\centering
			  	\includegraphics[width=0.49\textwidth]{figures/Meson/Systematics/EMC/SysMeanNewWithMean_Pi0_13TeVEG2.png}
			  	\includegraphics[width=0.49\textwidth]{figures/Meson/Systematics/EMC/SysMeanNewWithMean_Eta_13TeVEG2.png}
			  	\includegraphics[width=0.49\textwidth]{figures/Meson/Systematics/EMC/SysMeanNewWithMean_EtaToPi0_13TeVEG2.png}
			  	\caption{from the \ac{EMCal} method for \pp collisions at \s~=~13~TeV for the EG2 trigger. The colored points represent the individual error sources, while the black points represent the final systematic error for the $\pi^0$ meson transverse momentum spectrum.}
			  	\label{fig:SysErrEMCDetailed_EG2}
			 \end{figure}

			 \begin{figure}[t]
			  	\centering
			  	\includegraphics[width=0.49\textwidth]{figures/Meson/Systematics/EMC/SysMeanNewWithMean_Pi0_13TeVEG1.png}
			  	\includegraphics[width=0.49\textwidth]{figures/Meson/Systematics/EMC/SysMeanNewWithMean_Eta_13TeVEG1.png}
			  	\includegraphics[width=0.49\textwidth]{figures/Meson/Systematics/EMC/SysMeanNewWithMean_EtaToPi0_13TeVEG1.png}
			  	\caption{from the \ac{EMCal} method for \pp collisions at \s~=~13~TeV for the EG1 trigger. The colored points represent the individual error sources, while the black points represent the final systematic error for the $\pi^0$ meson transverse momentum spectrum.}
			  	\label{fig:SysErrEMCDetailed_EG1}
			  \end{figure}


% 		\begin{figure}[h]
% 			\centering
% 			\includegraphics[width=0.66\textwidth]{figures/Meson/Systematics/SysMeanNewWithMean_Pi0_8TeV_2017_11_02.eps}\\
% 			\includegraphics[width=0.66\textwidth]{figures/Meson/SystematicsTrigger/SysMeanNewWithMean_Pi0_8TeVEMC7_2017_11_02.eps}\\
% 			\includegraphics[width=0.66\textwidth]{figures/Meson/SystematicsTrigger/SysMeanNewWithMean_Pi0_8TeVEGA_2017_11_02.eps}
% 			\caption{Detailed systematic errors for \pp collisions at \sth. The colored points represent the individual error sources, while the black points represent the final systematic error for the $\pi^0$ meson transverse momentum spectrum.}
% 			\label{fig:SysErrPP2760Pi0Det}
% 		\end{figure}
%
% 		\begin{figure}
% 			\centering
% 			\includegraphics[width=0.7\textwidth]{figures/Meson/Systematics/SysMeanNewWithMean_Eta_8TeV_2017_11_02.eps}\\
% 			\includegraphics[width=0.7\textwidth]{figures/Meson/SystematicsTrigger/SysMeanNewWithMean_Eta_8TeVEMC7_2017_11_02.eps}\\
% 			\includegraphics[width=0.7\textwidth]{figures/Meson/SystematicsTrigger/SysMeanNewWithMean_Eta_8TeVEGA_2017_11_02.eps}
% 			\caption{Detailed systematic errors for \pp collisions at \sth. The colored points represent the individual error sources, while the black points represent the final systematic error for the $\eta$ meson transverse momentum spectrum.}
% 			\label{fig:SysErrPP2760EtaDet}
% 		\end{figure}
%
% 		\begin{figure}
% 			\centering
% 			\includegraphics[width=0.7\textwidth]{figures/Meson/Systematics/SysMeanNewWithMean_Pi0EtaBinning_8TeV_2017_11_02.eps}\\
% 			\includegraphics[width=0.7\textwidth]{figures/Meson/SystematicsTrigger/SysMeanNewWithMean_Pi0EtaBinning_8TeVEMC7_2017_11_02.eps}\\
% 			\includegraphics[width=0.7\textwidth]{figures/Meson/SystematicsTrigger/SysMeanNewWithMean_Pi0EtaBinning_8TeVEGA_2017_11_02.eps}
% 			\caption{Detailed systematic errors for \pp collisions at \sth. The colored points represent the individual error sources, while the black points represent the final systematic error for the \EtaToPi ratio versus transverse momentum.}
% 			\label{fig:SysErrPP2760EtaToPi0Det}
% 		\end{figure}
%
% 		\begin{figure}
% 			\centering
% 			\includegraphics[width=0.49\textwidth]{figures/Meson/Systematics/Pi0_data_Efficiencies_diffMCs.eps}
% 			\includegraphics[width=0.49\textwidth]{figures/Meson/Systematics/Pi0_data_CorrectedYield_MaterialError.eps}
% 			\caption{(left) Reconstruction efficiencies of the neutral pion for the two different generators used as well as the combination of both. (right) Comparison of the invariant yields when considering only \acs{EMCal} clusters with/without \acs{TRD} in front as well as considering the whole available \acs{EMCal} which is just the sum of the both specific cases.}
% 			\label{fig:sysEffMaterial}
% 		\end{figure}
%
% 		\begin{figure}
% 			\centering
% 			\includegraphics[width=0.49\textwidth]{figures/Meson/Systematics/Pi0_data_RAWYield_DistBC.eps}
% 			\includegraphics[width=0.49\textwidth]{figures/Meson/Systematics/Pi0_data_CorrectedYield_DistBC.eps}
% 			\caption{(left) Raw yields for neutral pion reconstruction for different distance to bad channel cuts. The different cuts applied are listed in the legend. (right) Comparison of the invariant yields for the different cuts. As it can be seen, they all agree well within uncertainties, although raw yields are lowered by 30~\% for certain cuts. Thus, a good description in Monte Carlo is proven so that no systematics may be assigned in this case.}
% 			\label{fig:sysDistBC}
% 		\end{figure}
%
% 		\begin{figure}
% 			\centering
% 		    \includegraphics[width=0.49\textwidth]{figures/Meson/Systematics/data_SPDPileUp.eps}
% 			\caption{The histogram is only filled if at least two primary vertices have been reconstructed; it shows the distance of the two primary vertices that have the highest distance in the event (unit of 'cm'). In blue, the rejection by SPD pileup cut is plotted whereas in green the rejection of SPD tracklet-cluster cut is shown. Assuming a Gaussian shape of the underlying distribution, the blue curve is fitted from -15 to 15 with a Gaussian (dotted black), excluding the center. Below 0.9~cm, the SPD pileup cut is not effective any more so that only a fraction of 'in-bunch' pileup can be rejected using the SPD tracklet-cluster cut, which is also fitted with a Gaussian (dotted red). In order to calculate the SPD pileup cut efficiency, the black dotted Gaussian is integrated from -0.9 to 0.9 and summed with the integral of red dotted Gaussian, integrated in the same region. This value is then divided by the black dotted Gaussian which is integrated from -15 to 15, arriving at the quoted efficiency of 92~\%.}
% 			\label{fig:SPDPileUp}
% 		\end{figure}

%		\begin{figure}
%			\centering
%			\includegraphics[width=0.49\textwidth]{figures/Meson/SystematicsTrigger/Pi0EtaBinning_data_CorrectedYield.eps}
%			\includegraphics[width=0.49\textwidth]{figures/Meson/SystematicsTrigger/Pi0_data_Mass.eps}
%			\caption{\newline Left plot: To properly consider the systematic error source due to the number of seven data taking periods, the neutral pion yield has been compared between the periods with enough statistics to do conclusive statements. The bins have been chosen to be wider than in the analysis in order to enhance the statistics per bin as we are dealing at single period level with approximately 10 times less statistics compared to the combined LHC12a-i analysis. Given the fluctuations around 1 and considering the statistical fluctuations, an error of 3\% has been concluded.\newline Right plot: The mass position for the same choice of data taking periods are shown here. Within statistical errors, all mass positions nicely agree to each other. It has to be noted that the statistics per single period is rather limited compared to the complete set of periods, so that especially at low \pT and higher \pT not all fits are converging due to limited number of neutral pion candidates per bin which explains the few given outliers in this plot.}
%			\label{fig:periodError}
%		\end{figure}
% \clearpage

\clearpage

\section{Acronyms and Technical Terms}
\begin{acronym}[ACORDEBB]
	\acro{ACORDE}{ALICE cosmic ray detector}
	\acro{AGS}{Alternating Gradient Synchrotron}
 	\acro{ALICE}{A Large Ion Collider Experiment}
	\acro{AOD}{Analysis Object Data}
	\acro{APD}{Avalanche Photodiode}
	\acro{ATLAS}{A Toroidal LHC Apparatus}
	\acro{BLUE}{Best Linear Unbiased Estimate}
	\acro{BNL}{Brockhaven National Laboratory}
	\acro{CERN}{European Organization for Nuclear Research}
	\acro{CDF}{Collider Detector at Fermilab}
	\acro{CMS}{Compact Muon Solenoid experiment}
	\acro{CCMF}{Conv-Calo mass fit}
	\acro{CMF}{Calo mass fit}
	\acro{CCRF}{Conv-Calo ratio fit}
	\acro{CRF}{Calo ratio fit}
	\acro{CPV}{Charged-Particle-Veto}
	\acro{CTP}{Central Trigger Processor}
	\acro{DAQ}{Data Acquisition}
	\acro{DCA}{distance of closest approach}
 	\acro{DCal}{Di-jet Calorimeter}
	\acro{DPM}{Dual Parton Model}
 	\acro{EMCal}{Electromagnetic Calorimeter}
	\acro{ESD}{Event Summary Data}
 	\acro{FEE}{Front End Electronics}
 	\acro{FMD}{Forward Multiplicity Detector}
 	\acro{FPGA}{Field Programmable Gate Array}
	\acro{FWHM}{full width at half maximum}
 	\acro{GEANT}{Geometry and Tracking Software}
	\acro{GTU}{Global Tracking Unit}
	\acro{HFE}{heavy flavor electron}
 	\acro{HG}{Hadron Gas}
 	\acro{HIJING}{Heavy Ion Jet Interaction Generator}
 	\acro{HLT}{High Level Trigger}
 	\acro{HMPID}{High Momentum Particle Identification Detector}
 	\acro{HYDJET}{Monte Carlo event generator for heavy ion collisions}
	\acro{INEL}{Inelastic}
 	\acro{ITS}{Inner Tracking System}
	\acro{L0}{level-0}
 	\acro{L1}{level-1}
 	\acro{L3}{High Energy Physics Experiment at LEP}
	\acro{LEP}{Large Electron Positron Collider}
	\acro{LEGO}{Lightweight Environment for Grid Operations}
	\acro{LEIR}{Low Energy Ion Ring}
	\acro{LHA}{Les Houches Accord}
	\acro{LHC}{Large Hadron Collider}
	\acro{LHCb}{LHC beauty experiment}
	\acro{LHEF}{Les Houches Event Files}
	\acro{LINAC2}{linear accelerator}
	\acro{LINAC3}{linear accelerator}
	\acro{LO}{Leading Order}
 	\acro{LTU}{Local Trigger Unit}
	\acro{LQCD}{Lattice QCD}
 	\acro{MC}{Monte Carlo simulation}
 	\acro{MRPC}{Multigap Resistive Plate Chamber}
	\acro{MWPC}{multi-wire-proportional chamber}
	\acro{NSD}{Non Single Diffractive Events}
	\acro{NLO}{Next-to-Leading Order}
	\acro{NNLO}{Next-to-Next-to Leading Order}
 	\acro{PDF}{Parton distribution function}
 	\acro{PHOS}{Photon Spectrometer}
 	\acro{RICH}{Ring Imaging Cherenkov detector}
	\acro{PCM}{Photon Conversion Method}
 	\acro{PDG}{Particle Data Group}
	\acro{PDF}{parton density functions}
	\acro{PHENIX}{Pioneering High Energy Nuclear Interactions eXperiment}
 	\acro{PID}{particle identification}
	\acro{PMD}{Photon Multiplicity Detector}
	\acro{pQCD}{perturbative QCD}
 	\acro{PS}{Proton Synchrotron}
	\acro{PSB}{Proton Synchrotron Booster}
	\acro{PYQUEN}{Monte Carlo event generator for pp-collisions}
	\acro{PYTHIA}{Monte Carlo event generator for pp-collisions}
	\acro{QCD}{Quantum Chromodynamics}
	\acro{QED}{Quantum Electrodynamics}
	\acro{QFT}{Quantum Field Theory}
	\acro{QGP}{Quark-Gluon Plasma}
	\acro{RCT}{run condition table}
	\acro{RHIC}{Relativistic Heavy Ion Collider}
	\acro{SDD}{Silicon Drift Detector}
	\acro{SDM}{Symmetric Decay Method}
	\acro{SHIM}{Secondary Hadronic Interaction Method}
	\acro{SM}{Standard Model}
	\acro{SPD}{Silicon Pixel Detector}
	\acro{SPS}{Super Proton Synchrotron}
	\acro{SSD}{Silicon Strip Detector}
	\acro{T0}[TZERO]{Timing and Trigger detector at ALICE}
	\acro{TBC}{Test beam corrected}
	\acro{TOF}{Time-Of-Flight detector}
	\acro{TPC}{Time Projection Chamber}
	\acro{TR}{transition radiation}
	\acro{TRD}{Transition Radiation Detector}
	\acro{TRG}{Trigger System}
	\acro{TRU}{Trigger Region Unit}
	\acro{V0}[V$^0$]{Unknown Particle}
	\acro{VZERO}{VZERO}
	\acro{ZDC}{Zero Degree Calorimeter}
	\acro{ZEM}{Zero Degree Electromagnetic Calorimeter}
	\acro{ZN}{Zero Degree Neutron Calorimeter}
	\acro{ZP}{Zero Degree Proton Calorimeter}
\end{acronym}

\clearpage

 \bibliography{../Bibliography.bib}
 \bibliographystyle{../kp.bst}


\end{appendix}


\clearpage
\showTODOs{EMCal}
% \showTODOs{PCMEMCal}
% \showTODOs{PCMPHOS}
\showTODOs{PCM}
\showTODOs{PHOS}
\showTODOs{All}

\end{document}
